\input{preamble}

% OK, start here.
%
\begin{document}

\title{Diviseurs sur les espaces algébriques}


\maketitle

\phantomsection
\label{section-phantom}

\tableofcontents

\section{Introduction}
\label{section-introduction}

\noindent
Dans ce chapitre, nous étudions les diviseurs sur les espaces algébriques et des sujets connexes.
Une référence de base sur les espaces algébriques est \cite{Kn}.









\section{Points associés et points faiblement associés}
\label{section-associated}

\noindent
Dans le cas des schémas, nous avons introduit deux notions concurrentes
de points associés : les points associés au sens usuel
(Diviseurs, section \ref{divisors-section-associated})
et les points faiblement associés
(Diviseurs, section \ref{divisors-section-weakly-associated}).
Pour un espace algébrique général, la notion de point associé
est pratiquement inutilisable, et nous ne chercherons même pas à la définir.
Si l'espace algébrique est localement noethérien, nous nous autorisons
à employer l'expression « point associé » à la place de
« point faiblement associé », puisque ces notions coïncident pour les
schémas noethériens (Diviseurs, lemme \ref{divisors-lemma-ass-weakly-ass}).
Avant de donner notre définition, nous avons besoin d'un lemme.

\begin{lemma}
\label{lemma-associated}
Soit $S$ un schéma. Soit $X$ un espace algébrique sur $S$.
Soit $\mathcal{F}$ un $\mathcal{O}_X$-module quasi-cohérent.
Soit $x \in |X|$. Les conditions suivantes sont équivalentes :
\begin{enumerate}
\item il existe un morphisme \'etale $f : U \to X$, où $U$ est un schéma,
et un point $u \in U$ d'image $x$ tels que $u$ soit faiblement associé
à $f^*\mathcal{F}$,
\item pour tout morphisme \'etale $f : U \to X$, où $U$ est un schéma,
et tout point $u \in U$ d'image $x$, le point $u$ est faiblement associé
à $f^*\mathcal{F}$,
\item l'idéal maximal de $\mathcal{O}_{X, \overline{x}}$
est un idéal premier faiblement associé au module $\mathcal{F}_{\overline{x}}$.
\end{enumerate}
Si $X$ est localement noethérien, elles sont également équivalentes à :
\begin{enumerate}
\item[(4)] il existe un morphisme \'etale $f : U \to X$, où $U$ est un schéma,
et un point $u \in U$ d'image $x$ tels que $u$ soit associé
à $f^*\mathcal{F}$,
\item[(5)] pour tout morphisme \'etale $f : U \to X$, où $U$ est un schéma,
et tout point $u \in U$ d'image $x$, le point $u$ est associé
à $f^*\mathcal{F}$,
\item[(6)] l'idéal maximal de $\mathcal{O}_{X, \overline{x}}$
est un idéal premier associé au module $\mathcal{F}_{\overline{x}}$.
\end{enumerate}
\end{lemma}

\begin{proof}
Choisissons un schéma $U$, un point $u$ et un morphisme \'etale
$f : U \to X$ qui envoie $u$ sur $x$. Relevons $\overline{x}$ en un point
géométrique de $U$ au-dessus de $u$. Rappelons que
$\mathcal{O}_{X, \overline{x}} = \mathcal{O}_{U, u}^{sh}$
où le hensélisé strict est pris relativement au
relèvement choisi de $\overline{x}$ ; voir
Propriétés des espaces, lemme
\ref{spaces-properties-lemma-describe-etale-local-ring}.
Enfin, nous avons
$$
\mathcal{F}_{\overline{x}} =
(f^*\mathcal{F})_u \otimes_{\mathcal{O}_{U, u}}
\mathcal{O}_{X, \overline{x}} =
(f^*\mathcal{F})_u \otimes_{\mathcal{O}_{U, u}}
\mathcal{O}_{U, u}^{sh}
$$
d'après
Propriétés des espaces, lemme \ref{spaces-properties-lemma-stalk-quasi-coherent}.
L'équivalence de (1), (2) et (3) résulte donc de
Compléments sur la platitude, lemme \ref{flat-lemma-weakly-associated-henselization}.
Si $X$ est localement noethérien,
tout $U$ comme ci-dessus est localement noethérien ;
les conditions (1), resp.\ (2), sont donc équivalentes aux conditions (4), resp.\ (5), d'après
Diviseurs, lemme \ref{divisors-lemma-ass-weakly-ass}.
D'autre part, dans le cas localement noethérien,
l'anneau local $\mathcal{O}_{X, \overline{x}}$ est lui aussi noethérien
(Propriétés des espaces, lemme
\ref{spaces-properties-lemma-Noetherian-local-ring-Noetherian}).
L'équivalence de (3) et (6) résulte donc du même lemme
(ou par Algèbre, lemme \ref{algebra-lemma-ass-weakly-ass}).
\end{proof}

\begin{definition}
\label{definition-weakly-associated}
Soit $S$ un schéma. Soit $X$ un espace algébrique sur $S$.
Soit $\mathcal{F}$ un faisceau quasi-cohérent sur $X$.
Soit $x \in |X|$.
\begin{enumerate}
\item On dit que $x$ est {\it faiblement associé} à $\mathcal{F}$
si les conditions équivalentes (1), (2) et (3) du
lemme \ref{lemma-associated} sont satisfaites.
\item On note $\text{WeakAss}(\mathcal{F})$ l'ensemble des points
faiblement associés à $\mathcal{F}$.
\item Les {\it points faiblement associés à $X$} sont les points
faiblement associés à $\mathcal{O}_X$.
\end{enumerate}
Si $X$ est localement noethérien, on dira
{\it $x$ est associé à $\mathcal{F}$}
si et seulement si $x$ est faiblement associé à $\mathcal{F}$, et l'on pose
$\text{Ass}(\mathcal{F}) = \text{WeakAss}(\mathcal{F})$.
Enfin (toujours sous l'hypothèse que $X$ est localement noethérien),
on dira que {\it $x$ est un point associé de $X$} si et seulement si
$x$ est un point faiblement associé à $X$.
\end{definition}

\noindent
À ce stade, nous pouvons démontrer les lemmes de rigueur.

\begin{lemma}
\label{lemma-weakly-ass-support}
Soit $S$ un schéma. Soit $X$ un espace algébrique sur $S$.
Soit $\mathcal{F}$ un $\mathcal{O}_X$-module quasi-cohérent.
Alors $\text{WeakAss}(\mathcal{F}) \subset \text{Supp}(\mathcal{F})$.
\end{lemma}

\begin{proof}
Cela résulte immédiatement des définitions. Le support d'un faisceau abélien
sur $X$ est défini dans Propriétés des espaces, définition
\ref{spaces-properties-definition-support}.
\end{proof}

\begin{lemma}
\label{lemma-ses-weakly-ass}
Soit $S$ un schéma. Soit $X$ un espace algébrique sur $S$.
Soit $0 \to \mathcal{F}_1 \to \mathcal{F}_2 \to \mathcal{F}_3 \to 0$
une suite exacte de faisceaux quasi-cohérents sur $X$.
Alors
$\text{WeakAss}(\mathcal{F}_2) \subset
\text{WeakAss}(\mathcal{F}_1) \cup \text{WeakAss}(\mathcal{F}_3)$
et
$\text{WeakAss}(\mathcal{F}_1) \subset \text{WeakAss}(\mathcal{F}_2)$.
\end{lemma}

\begin{proof}
Pour tout point géométrique $\overline{x} \in X$,
la suite des fibres
$0 \to \mathcal{F}_{1, \overline{x}} \to
\mathcal{F}_{2, \overline{x}} \to
\mathcal{F}_{3, \overline{x}} \to 0$
est une suite exacte de $\mathcal{O}_{X, \overline{x}}$-modules.
Le lemme résulte donc de
Algèbre, lemme \ref{algebra-lemma-weakly-ass}.
\end{proof}

\begin{lemma}
\label{lemma-weakly-ass-zero}
Soit $S$ un schéma. Soit $X$ un espace algébrique sur $S$.
Soit $\mathcal{F}$ un $\mathcal{O}_X$-module quasi-cohérent.
Alors
$$
\mathcal{F} = (0) \Leftrightarrow \text{WeakAss}(\mathcal{F}) = \emptyset
$$
\end{lemma}

\begin{proof}
Choisissons un schéma $U$ et un morphisme \'etale surjectif $f : U \to X$.
Alors $\mathcal{F}$ est nul si et seulement si $f^*\mathcal{F}$ est nul.
Le lemme résulte donc de la définition et du lemme dans le
cas des schémas ; voir
Diviseurs, lemme \ref{divisors-lemma-weakly-ass-zero}.
\end{proof}

\begin{lemma}
\label{lemma-minimal-support-in-weakly-ass}
Soit $S$ un schéma. Soit $X$ un espace algébrique sur $S$.
Soit $\mathcal{F}$ un $\mathcal{O}_X$-module quasi-cohérent.
Soit $x \in |X|$. Si
\begin{enumerate}
\item $x \in \text{Supp}(\mathcal{F})$
\item $x$ est un point de codimension $0$ de $X$
(Propriétés des espaces, définition
\ref{spaces-properties-definition-dimension-local-ring}).
\end{enumerate}
Alors $x \in \text{WeakAss}(\mathcal{F})$. Si $\mathcal{F}$
est un $\mathcal{O}_X$-module de type fini de support schématique $Z$
(Morphismes d'espaces, définition
\ref{spaces-morphisms-definition-scheme-theoretic-support})
et si $x$ est un point de codimension $0$ de $Z$, alors
$x \in \text{WeakAss}(\mathcal{F})$.
\end{lemma}

\begin{proof}
Puisque $x \in \text{Supp}(\mathcal{F})$, la fibre
$\mathcal{F}_{\overline{x}}$ n'est pas nulle. Ainsi,
$\text{WeakAss}(\mathcal{F}_{\overline{x}})$
est non vide d'après
Algèbre, lemme \ref{algebra-lemma-weakly-ass-zero}.
D'autre part, le spectre de $\mathcal{O}_{X, \overline{x}}$
est réduit à un point. Par définition, $x$ est donc faiblement associé à
$\mathcal{F}$. Pour la dernière assertion, l'homomorphisme
$\mathcal{O}_{X, \overline{x}} \to \mathcal{O}_{Z, \overline{z}}$
est surjectif, le spectre de $\mathcal{O}_{Z, \overline{z}}$
est réduit à un point, et $\mathcal{F}_{\overline{x}}$ est un
module non nul sur $\mathcal{O}_{Z, \overline{z}}$.
\end{proof}

\begin{lemma}
\label{lemma-minimal-support-in-weakly-ass-decent}
Soit $S$ un schéma. Soit $X$ un espace algébrique sur $S$.
Soit $\mathcal{F}$ un $\mathcal{O}_X$-module quasi-cohérent.
Soit $x \in |X|$. Si
\begin{enumerate}
\item $X$ est décent (par exemple quasi-séparé ou localement séparé),
\item $x \in \text{Supp}(\mathcal{F})$
\item $x$ n'est pas une spécialisation d'un autre point dans
$\text{Supp}(\mathcal{F})$.
\end{enumerate}
Alors $x \in \text{WeakAss}(\mathcal{F})$.
\end{lemma}

\begin{proof}
(Un espace algébrique quasi-séparé est décent ; voir
Espaces décents, section \ref{decent-spaces-section-reasonable-decent}.
Un espace algébrique localement séparé est décent ; voir
Espaces décents, lemme \ref{decent-spaces-lemma-locally-separated-decent}.)
Choisissons un schéma $U$, un point $u \in U$ et un morphisme \'etale
$f : U \to X$ envoyant $u$ sur $x$. D'après
Espaces décents, lemme
\ref{decent-spaces-lemma-decent-no-specializations-map-to-same-point}
si $u' \leadsto u$ est une spécialisation non triviale, alors
$f(u') \not = x$. Ainsi, $u \in \text{Supp}(f^*\mathcal{F})$
n'est pas une spécialisation d'un autre point de
$\text{Supp}(f^*\mathcal{F})$.
Par conséquent, $u \in \text{WeakAss}(f^*\mathcal{F})$ d'après
Diviseurs, lemme \ref{lemma-minimal-support-in-weakly-ass}.
\end{proof}

\begin{lemma}
\label{lemma-finite-ass}
Soit $S$ un schéma. Soit $X$ un espace algébrique localement noethérien sur $S$.
Soit $\mathcal{F}$ un $\mathcal{O}_X$-module cohérent.
Alors $\text{Ass}(\mathcal{F}) \cap W$ est fini pour
tout ouvert quasi-compact $W \subset |X|$.
\end{lemma}

\begin{proof}
Choisissons un schéma quasi-compact $U$ et un morphisme \'etale $U \to X$
tel que $W$ soit l'image de $|U| \to |X|$. Alors $U$ est un
schéma noethérien, et l'on conclut en appliquant
Diviseurs, lemme \ref{divisors-lemma-finite-ass}.
\end{proof}

\begin{lemma}
\label{lemma-restriction-injective-open-contains-weakly-ass}
Soit $S$ un schéma. Soit $X$ un espace algébrique sur $S$.
Soit $\mathcal{F}$ un $\mathcal{O}_X$-module quasi-cohérent.
Si $U \to X$ est un morphisme \'etale tel que
$\text{WeakAss}(\mathcal{F}) \subset \Im(|U| \to |X|)$, alors
$\Gamma(X, \mathcal{F}) \to \Gamma(U, \mathcal{F})$ est injectif.
\end{lemma}

\begin{proof}
Soit $s \in \Gamma(X, \mathcal{F})$ une section dont la restriction à $U$ est nulle.
Soit $\mathcal{F}' \subset \mathcal{F}$ l'image de l'homomorphisme
$\mathcal{O}_X \to \mathcal{F}$ défini par $s$. Alors $\mathcal{F}'|_U = 0$.
Cela entraîne
$\text{WeakAss}(\mathcal{F}') \cap \Im(|U| \to |X|) = \emptyset$
(par définition des points faiblement associés).
D'autre part,
$\text{WeakAss}(\mathcal{F}') \subset \text{WeakAss}(\mathcal{F})$
par le lemme \ref{lemma-ses-weakly-ass}. Nous en déduisons que
$\text{WeakAss}(\mathcal{F}') = \emptyset$.
Ainsi, $\mathcal{F}' = 0$ par le lemme \ref{lemma-weakly-ass-zero}.
\end{proof}

\begin{lemma}
\label{lemma-weakass-pushforward}
Soit $S$ un schéma. Soit $f : X \to Y$ un morphisme quasi-compact et quasi-séparé
d'espaces algébriques sur $S$. Soit $\mathcal{F}$ un
$\mathcal{O}_X$-module quasi-cohérent. Soit $y \in |Y|$ un point qui n'appartient pas à
l'image de $|f|$. Alors $y$ n'est pas faiblement associé à $f_*\mathcal{F}$.
\end{lemma}

\begin{proof}
D'après Morphismes d'espaces, lemme \ref{spaces-morphisms-lemma-pushforward},
le $\mathcal{O}_Y$-module $f_*\mathcal{F}$ est quasi-cohérent, de sorte que
l'énoncé a un sens.
Choisissons un schéma affine $V$, un point $v \in V$ et un morphisme \'etale
$V \to Y$ envoyant $v$ sur $y$. Nous pouvons remplacer
$f : X \to Y$, $\mathcal{F}$, $y$ par
$X \times_Y V \to V$, $\mathcal{F}|_{X \times_Y V}$, $v$.
Nous pouvons donc supposer que $Y$ est un schéma affine.
Dans ce cas, $X$ est quasi-compact ; nous pouvons donc choisir
un schéma affine $U$ et un morphisme \'etale surjectif $U \to X$.
Notons $g : U \to Y$ le composé.
Alors $f_*\mathcal{F} \subset g_*(\mathcal{F}|_U)$.
Par le lemme \ref{lemma-ses-weakly-ass},
on se ramène au cas des schémas, traité dans
Diviseurs, lemme \ref{divisors-lemma-weakass-pushforward}.
\end{proof}

\begin{lemma}
\label{lemma-check-injective-on-weakass}
Soit $S$ un schéma. Soit $X$ un espace algébrique sur $S$.
Soit $\varphi : \mathcal{F} \to \mathcal{G}$ un homomorphisme de
$\mathcal{O}_X$-modules quasi-cohérents. Supposons que, pour tout
$x \in |X|$, l'une au moins des conditions suivantes soit satisfaite :
\begin{enumerate}
\item $\mathcal{F}_{\overline{x}} \to \mathcal{G}_{\overline{x}}$
est injectif, ou
\item $x \not \in \text{WeakAss}(\mathcal{F})$.
\end{enumerate}
Alors $\varphi$ est injectif.
\end{lemma}

\begin{proof}
Les hypothèses impliquent que $\text{WeakAss}(\Ker(\varphi)) = \emptyset$
et donc que $\Ker(\varphi) = 0$ par le lemme \ref{lemma-weakly-ass-zero}.
\end{proof}

\begin{lemma}
\label{lemma-weakass-reduced}
Soit $S$ un schéma. Soit $X$ un espace algébrique réduit sur $S$.
Alors les points faiblement associés à $X$ sont exactement les points de codimension $0$
de $X$.
\end{lemma}

\begin{proof}
En travaillant localement pour la topologie \'etale, cela résulte de
Diviseurs, lemme \ref{divisors-lemma-weakass-reduced}
et
Propriétés des espaces, lemme \ref{spaces-properties-lemma-codimension-0-points}.
\end{proof}










\section{Morphismes et points faiblement associés}
\label{section-morphisms-weakly-associated}

\begin{lemma}
\label{lemma-weakly-ass-reverse-functorial}
Soit $S$ un schéma.
Soit $f : X \to Y$ un morphisme affine d'espaces algébriques sur $S$.
Soit $\mathcal{F}$ un $\mathcal{O}_X$-module quasi-cohérent.
On a alors
$$
\text{WeakAss}_S(f_*\mathcal{F}) \subset f(\text{WeakAss}_X(\mathcal{F}))
$$
\end{lemma}

\begin{proof}
Choisissons un schéma $V$ et un morphisme \'etale surjectif $V \to Y$.
Posons $U = X \times_Y V$. Alors $U \to V$ est un morphisme affine
de schémas. Par définition des points faiblement associés,
le problème est ramené au morphisme de schémas $U \to V$. Ce cas est
traité dans Diviseurs, lemme \ref{divisors-lemma-weakly-ass-reverse-functorial}.
\end{proof}

\begin{lemma}
\label{lemma-ass-functorial-equal}
Soit $S$ un schéma.
Soit $f : X \to Y$ un morphisme affine d'espaces algébriques sur $S$.
Soit $\mathcal{F}$ un $\mathcal{O}_X$-module quasi-cohérent.
Si $X$ est localement noethérien, alors on a
$$
\text{WeakAss}_Y(f_*\mathcal{F}) =
f(\text{WeakAss}_X(\mathcal{F}))
$$
\end{lemma}

\begin{proof}
Choisissons un schéma $V$ et un morphisme \'etale surjectif $V \to Y$.
Posons $U = X \times_Y V$. Alors $U \to V$ est un morphisme affine
et $U$ est localement noethérien.
Par définition des points faiblement associés,
le problème est ramené au morphisme de schémas $U \to V$. Ce cas est
traité dans Diviseurs, lemme \ref{divisors-lemma-ass-functorial-equal}.
\end{proof}

\begin{lemma}
\label{lemma-weakly-associated-finite}
Soit $S$ un schéma.
Soit $f : X \to Y$ un morphisme fini d'espaces algébriques sur $S$.
Soit $\mathcal{F}$ un $\mathcal{O}_X$-module quasi-cohérent.
Alors $\text{WeakAss}(f_*\mathcal{F}) = f(\text{WeakAss}(\mathcal{F}))$.
\end{lemma}

\begin{proof}
Choisissons un schéma $V$ et un morphisme \'etale surjectif $V \to Y$.
Posons $U = X \times_Y V$. Alors $U \to V$ est un morphisme fini
de schémas. Par définition des points faiblement associés,
le problème est ramené au morphisme de schémas $U \to V$. Ce cas est
traité dans Diviseurs, lemme \ref{divisors-lemma-weakly-associated-finite}.
\end{proof}

\begin{lemma}
\label{lemma-weakly-ass-pullback}
Soit $S$ un schéma. Soit $f : X \to Y$ un morphisme d'espaces algébriques
sur $S$. Soit $\mathcal{G}$ un $\mathcal{O}_Y$-module quasi-cohérent.
Soit $x \in |X|$ et $y = f(x) \in |Y|$. Si
\begin{enumerate}
\item $y \in \text{WeakAss}_S(\mathcal{G})$,
\item $f$ est plat en $x$, et
\item la dimension de l'anneau local de la fibre de $f$ au point $x$
est nulle (Morphismes d'espaces, définition
\ref{spaces-morphisms-definition-dimension-fibre}),
\end{enumerate}
alors $x \in \text{WeakAss}(f^*\mathcal{G})$.
\end{lemma}

\begin{proof}
Choisissons un schéma $V$, un point $v \in V$ et un morphisme \'etale $V \to Y$
envoyant $v$ sur $y$. Choisissons un schéma $U$, un point $u \in U$ et un
morphisme \'etale $U \to V \times_Y X$ envoyant $v$ sur un point situé au-dessus de
$v$ et de $x$. Cela est possible car il existe $t \in |V \times_Y X|$
s'envoyant sur $(v, y)$ d'après Propriétés des espaces, lemme
\ref{spaces-properties-lemma-points-cartesian}.
Par définition, la dimension de $\mathcal{O}_{U_v, u}$ est nulle.
Ainsi, $u$ est un point générique de la fibre $U_v$.
Par définition des points faiblement associés,
le problème est ramené au morphisme de schémas $U \to V$.
Ce cas est traité dans
Diviseurs, lemme \ref{divisors-lemma-weakly-ass-pullback}.
\end{proof}

\begin{lemma}
\label{lemma-weakly-ass-change-fields}
Soit $K/k$ une extension de corps. Soit $X$ un espace algébrique sur $k$.
Soit $\mathcal{F}$ un $\mathcal{O}_X$-module quasi-cohérent.
Soit $y \in X_K$ d'image $x \in X$. Si $y$ est un point
faiblement associé à l'image inverse $\mathcal{F}_K$, alors $x$
est un point faiblement associé à $\mathcal{F}$.
\end{lemma}

\begin{proof}
C'est la traduction de
Diviseurs, lemme \ref{divisors-lemma-weakly-ass-change-fields}
dans le langage des espaces algébriques. Nous omettons les détails de cette
traduction.
\end{proof}

\begin{lemma}
\label{lemma-finite-flat-weak-assassin-up-down}
Soit $S$ un schéma.
Soit $f : X \to Y$ un morphisme fini et plat d'espaces algébriques.
Soit $\mathcal{G}$ un $\mathcal{O}_Y$-module quasi-cohérent.
Soit $x \in |X|$ un point d'image $y \in |Y|$. Alors
$$
x \in \text{WeakAss}(g^*\mathcal{G})
\Leftrightarrow
y \in \text{WeakAss}(\mathcal{G})
$$
\end{lemma}

\begin{proof}
Cela résulte immédiatement du cas des schémas
(Compléments sur la platitude, lemme \ref{flat-lemma-finite-flat-weak-assassin-up-down})
par localisation \'etale.
\end{proof}

\begin{lemma}
\label{lemma-etale-weak-assassin-up-down}
Soit $S$ un schéma.
Soit $f : X \to Y$ un morphisme \'etale d'espaces algébriques.
Soit $\mathcal{G}$ un $\mathcal{O}_Y$-module quasi-cohérent.
Soit $x \in |X|$ un point d'image $y \in |Y|$. Alors
$$
x \in \text{WeakAss}(f^*\mathcal{G})
\Leftrightarrow
y \in \text{WeakAss}(\mathcal{G})
$$
\end{lemma}

\begin{proof}
Cela résulte immédiatement de la définition des points faiblement associés
et, en fait, le lemme correspondant pour les schémas
(Compléments sur la platitude, lemme \ref{flat-lemma-etale-weak-assassin-up-down})
est à la base de notre définition.
\end{proof}







\section{Assassin faible relatif}
\label{section-relative-weak-assassin}

\noindent
Nous avons besoin de quelques lemmes pour définir cet objet.

\begin{lemma}
\label{lemma-locally-noetherian-fibre}
Soit $S$ un schéma. Soit $f : X \to Y$ un morphisme d'espaces algébriques
sur $S$. Soit $y \in |Y|$. Les conditions suivantes sont équivalentes :
\begin{enumerate}
\item il existe un schéma $V$, un point $v \in V$ et un morphisme \'etale $V \to Y$
envoyant $v$ sur $y$ tels que l'espace algébrique $X_v$ soit localement noethérien,
\item pour tout schéma $V$, tout point $v \in V$ et tout morphisme \'etale $V \to Y$
envoyant $v$ sur $y$, l'espace algébrique $X_v$ est localement noethérien, et
\item il existe un corps $k$ et un morphisme $\Spec(k) \to Y$ représentant
$y$ tels que $X_k$ soit localement noethérien.
\end{enumerate}
S'il existe un corps $k_0$ et un monomorphisme $\Spec(k_0) \to Y$
représentant $y$, ces conditions sont également équivalentes à
\begin{enumerate}
\item[(4)] l'espace algébrique $X_{k_0}$ est localement noethérien.
\end{enumerate}
\end{lemma}

\begin{proof}
Remarquons que $X_v = v \times_Y X = \Spec(\kappa(v)) \times_Y X$.
Les implications (2) $\Rightarrow$ (1) $\Rightarrow$ (3) sont donc claires.
Supposons que $\Spec(k) \to Y$ soit un morphisme depuis le spectre d'un corps
tel que $X_k$ soit localement noethérien. Soit $V \to Y$ un morphisme \'etale
depuis un schéma $V$, et soit $v \in V$ un point envoyé sur $y$.
Alors le schéma $v \times_Y \Spec(k)$ est non vide. Choisissons un
point $w \in v \times_Y \Spec(k)$. Considérons les morphismes
$$
X_v \longleftarrow X_w \longrightarrow X_k
$$
Comme $V \to Y$ est \'etale et que $w$ peut être considéré comme un point de
$V \times_Y \Spec(k)$, l'extension $\kappa(w)/k$
est une extension finie séparable de corps
(Morphismes, lemme \ref{morphisms-lemma-etale-over-field}).
Ainsi, $X_w \to X_k$ est un morphisme fini \'etale obtenu par changement de base de
$w \to \Spec(k)$. Par conséquent, $X_w$ est localement noethérien
(Morphismes des espaces, lemme
\ref{spaces-morphisms-lemma-locally-finite-type-locally-noetherian}).
Le morphisme $X_w \to X_v$ est surjectif, affine et plat,
car il est obtenu par changement de base du morphisme surjectif, affine et plat $w \to v$.
Le fait que $X_w$ soit localement noethérien implique alors que
$X_v$ est localement noethérien. Pour le voir, choisissons un morphisme
\'etale surjectif $U \to X$, puis remarquons que
$U_w \to U_v$ est surjectif, affine et plat. En nous plaçant
localement dans le cas affine sur le schéma $U_v$, nous concluons
que $U_w$ est localement noethérien par
Algèbre, lemme \ref{algebra-lemma-descent-Noetherian}.

\medskip\noindent
Enfin, il suffit de prouver que (3) implique (4)
lorsqu'il existe un monomorphisme $\Spec(k_0) \to Y$ dans la classe de $y$.
Alors $\Spec(k) \to Y$ se factorise en $\Spec(k) \to \Spec(k_0) \to Y$.
L'argument donné ci-dessus montre alors que le fait que $X_k$ soit
localement noethérien implique que $X_{k_0}$ est localement noethérien.
\end{proof}

\begin{definition}
\label{definition-locally-Noetherian-fibre}
Soit $S$ un schéma. Soit $f : X \to Y$ un morphisme d'espaces algébriques
sur $S$. Soit $y \in |Y|$. Nous disons que {\it la fibre de $f$ au-dessus de $y$ est
localement noethérienne} si les conditions équivalentes (1), (2) et (3)
du lemme \ref{lemma-locally-noetherian-fibre} sont satisfaites.
Nous disons que {\it les fibres de $f$ sont localement noethériennes} si cette propriété
est satisfaite pour tout $y \in |Y|$.
\end{definition}

\noindent
Bien entendu, la manière usuelle de garantir que les fibres sont localement noethériennes consiste
à supposer que le morphisme est localement de type fini.

\begin{lemma}
\label{lemma-locally-finite-type-locally-Noetherian-fibres}
Soit $S$ un schéma. Soit $f : X \to Y$ un morphisme d'espaces algébriques
sur $S$. Si $f$ est localement de type fini, alors
les fibres de $f$ sont localement noethériennes.
\end{lemma}

\begin{proof}
Cela résulte de Morphismes d'espaces, lemme
\ref{spaces-morphisms-lemma-locally-finite-type-locally-noetherian}
et du fait que le spectre d'un corps est noethérien.
\end{proof}

\begin{lemma}
\label{lemma-relative-assassin}
Soit $S$ un schéma. Soit $f : X \to Y$ un morphisme d'espaces algébriques
sur $S$. Soit $x \in |X|$ et $y = f(x) \in |Y|$.
Soit $\mathcal{F}$ un
$\mathcal{O}_X$-module quasi-cohérent. Considérons les diagrammes commutatifs
$$
\xymatrix{
X \ar[d] & X \times_Y V \ar[d] \ar[l] & X_v \ar[d] \ar[l] \\
Y & V \ar[l] & v \ar[l]
}
\quad
\xymatrix{
X \ar[d] & U \ar[d] \ar[l] & U_v \ar[d] \ar[l] \\
Y & V \ar[l] & v \ar[l]
}
\quad
\xymatrix{
x \ar@{|->}[d] &
x' \ar@{|->}[d] \ar@{|->}[l] &
u \ar@{|->}[ld] \ar@{|->}[l] \\
y &
v \ar@{|->}[l]
}
$$
où $V$ et $U$ sont des schémas, $V \to Y$ et $U \to X \times_Y V$
sont \'etales, et $v \in V$, $x' \in |X_v|$, $u \in U$ sont des points
reliés comme dans le dernier diagramme.
Notons $\mathcal{F}|_{X_v}$ et $\mathcal{F}|_{U_v}$
les images inverses de $\mathcal{F}$.
Les conditions suivantes sont équivalentes :
\begin{enumerate}
\item il existe $V, v, x'$ comme ci-dessus tels que $x'$ soit un point
faiblement associé à $\mathcal{F}|_{X_v}$,
\item pour tous $V \to Y, v, x'$ comme ci-dessus, $x'$ est un point
faiblement associé à $\mathcal{F}|_{X_v}$,
\item il existe $U, V, u, v$ comme ci-dessus tels que $u$ soit un point
faiblement associé à $\mathcal{F}|_{U_v}$,
\item pour tous $U, V, u, v$ comme ci-dessus, $u$ est un point
faiblement associé à $\mathcal{F}|_{U_v}$,
\item il existe un corps $k$, un morphisme $\Spec(k) \to Y$ représentant $y$
et un point $t \in |X_k|$ s'envoyant sur $x$, tels que $t$ soit un point
faiblement associé à $\mathcal{F}|_{X_k}$.
\end{enumerate}
S'il existe un corps $k_0$ et un monomorphisme $\Spec(k_0) \to Y$
représentant $y$, ces conditions sont également équivalentes à
\begin{enumerate}
\item[(6)] $x_0$ est un point faiblement associé à $\mathcal{F}|_{X_{k_0}}$
où $x_0 \in |X_{k_0}|$ est l'unique point s'envoyant sur $x$.
\end{enumerate}
Si la fibre de $f$ au-dessus de $y$ est localement noethérienne, alors, dans les
conditions (1), (2), (3), (4) et (6), nous pouvons remplacer
``faiblement associé'' par ``associé''.
\end{lemma}

\begin{proof}
Remarquons qu'étant donnés $V, v, x'$ comme dans le lemme, nous pouvons trouver
$U \to X \times_Y V$ et $u \in U$ s'envoyant sur $x'$
et qu'alors le morphisme $U_v \to X_v$ est \'etale.
Il est donc clair que (1) et (3) sont équivalentes,
de même que (2) et (4). Chacune de ces conditions implique (5).
Nous allons montrer que (5) implique (2).
Supposons donnés $V, v, x'$, ainsi que $\Spec(k) \to X$ et $t \in |X_k|$,
de sorte que $t$ soit un point faiblement
associé à $\mathcal{F}|_{X_k}$.
Nous pouvons choisir un point $w \in v \times_Y \Spec(k)$.
Nous obtenons alors les morphismes
$$
X_v \longleftarrow X_w \longrightarrow X_k
$$
Comme $V \to Y$ est \'etale et que $w$ peut être considéré comme un point de
$V \times_Y \Spec(k)$, l'extension $\kappa(w)/k$
est une extension finie séparable de corps
(Morphismes, lemme \ref{morphisms-lemma-etale-over-field}).
Ainsi, $X_w \to X_k$ est un morphisme fini \'etale obtenu par changement de base de
$w \to \Spec(k)$. Tout point $x''$ de $X_w$ situé au-dessus de $t$
est donc un point faiblement associé à $\mathcal{F}|_{X_w}$ par le
lemme \ref{lemma-etale-weak-assassin-up-down}.
Nous pouvons choisir $x''$ s'envoyant sur $x'$
(Propriétés des espaces, lemme \ref{spaces-properties-lemma-points-cartesian}).
Le lemme \ref{lemma-weakly-ass-change-fields}
implique alors que $x'$ est un point faiblement associé
à $\mathcal{F}|_{X_v}$.

\medskip\noindent
Pour achever la preuve, il suffit de montrer que les
conditions équivalentes (1) -- (5) impliquent (6) lorsque l'on se donne
$\Spec(k_0) \to Y$ comme dans (6). Dans ce cas, le morphisme
$\Spec(k) \to Y$ de (5) se factorise de manière unique en $\Spec(k) \to \Spec(k_0) \to Y$.
Alors $x_0$ est l'image de $t$ par le morphisme $X_k \to X_{k_0}$.
Le même lemme que ci-dessus montre donc que (6) est vraie.
\end{proof}

\begin{definition}
\label{definition-relative-weak-assassin}
Soit $S$ un schéma. Soit $f : X \to Y$ un morphisme d'espaces algébriques
sur $S$. Soit $\mathcal{F}$ un $\mathcal{O}_X$-module quasi-cohérent.
L'{\it assassin faible relatif de $\mathcal{F}$ dans $X$ sur $Y$}
est l'ensemble $\text{WeakAss}_{X/Y}(\mathcal{F}) \subset |X|$
des $x \in |X|$ pour lesquels les conditions équivalentes du
lemme \ref{lemma-relative-assassin} sont satisfaites.
Si les fibres de $f$ sont localement noethériennes
(définition \ref{definition-locally-Noetherian-fibre}),
on utilise la notation $\text{Ass}_{X/Y}(\mathcal{F})$.
\end{definition}

\noindent
Avec ces notations, nous pouvons énoncer certains des résultats
déjà démontrés pour les schémas.

\begin{lemma}
\label{lemma-bourbaki}
Soit $S$ un schéma.
Soit $f : X \to Y$ un morphisme d'espaces algébriques sur $S$.
Soit $\mathcal{F}$ un $\mathcal{O}_X$-module quasi-cohérent.
Soit $\mathcal{G}$ un $\mathcal{O}_Y$-module quasi-cohérent.
Supposons que
\begin{enumerate}
\item $\mathcal{F}$ est plat sur $Y$,
\item $X$ et $Y$ sont localement noethériens, et
\item les fibres de $f$ sont localement noethériennes.
\end{enumerate}
Alors
$$
\text{Ass}_X(\mathcal{F} \otimes_{\mathcal{O}_X} f^*\mathcal{G}) =
\{x \in \text{Ass}_{X/Y}(\mathcal{F})\text{ tel que }
f(x) \in \text{Ass}_Y(\mathcal{G}) \}
$$
\end{lemma}

\begin{proof}
Par localisation \'etale, l'énoncé résulte immédiatement du résultat
pour les schémas ; voir
Diviseurs, lemme \ref{divisors-lemma-bourbaki}.
Le résultat pour les schémas n'est plus général que parce que
nous n'avons pas défini les points associés pour les
espaces algébriques non noethériens (nous devons donc supposer que $X$
et les fibres de $X \to Y$ sont localement noethériens pour pouvoir
même formuler cet énoncé).
\end{proof}

\begin{lemma}
\label{lemma-base-change-relative-assassin}
Soit $S$ un schéma. Soit
$$
\xymatrix{
X' \ar[d]_{f'} \ar[r]_{g'} & X \ar[d]^f \\
Y' \ar[r]^g & Y
}
$$
un diagramme cartésien d'espaces algébriques sur $S$.
Soit $\mathcal{F}$ un $\mathcal{O}_X$-module quasi-cohérent
et posons $\mathcal{F}' = (g')^*\mathcal{F}$.
Si $f$ est localement de type fini, alors
\begin{enumerate}
\item $x' \in \text{Ass}_{X'/Y'}(\mathcal{F}')
\Rightarrow g'(x') \in \text{Ass}_{X/Y}(\mathcal{F})$
\item si $x \in \text{Ass}_{X/Y}(\mathcal{F})$, alors, étant donné
$y' \in |Y'|$ tel que $f(x) = g(y')$, il existe
$x' \in \text{Ass}_{X'/Y'}(\mathcal{F}')$
avec $g'(x') = x$ et $f'(x') = y'$.
\end{enumerate}
\end{lemma}

\begin{proof}
Cela résulte du cas des schémas par localisation \'etale.
Donnons tous les détails. Choisissons un schéma
$V$ et un morphisme \'etale surjectif $V \to Y$.
Choisissons un schéma $U$ et un morphisme
\'etale surjectif $U \to V \times_Y X$. Choisissons un schéma $V'$
et un morphisme \'etale surjectif $V' \to V \times_Y Y'$.
Alors $U' = V' \times_V U$ est un schéma et le morphisme
$U' \to X'$ est surjectif et \'etale.

\medskip\noindent
Preuve de (1). Choisissons $u' \in U'$ s'envoyant sur $x'$.
Notons $v' \in V'$ l'image de $u'$.
La condition $x' \in \text{Ass}_{X'/Y'}(\mathcal{F}')$ équivaut
à $u' \in \text{Ass}(\mathcal{F}|_{U'_{v'}})$
par définition (il est légitime d'écrire $\text{Ass}$ au lieu de $\text{WeakAss}$
car $U'_{v'}$ est localement noethérien).
En appliquant Diviseurs, lemme \ref{divisors-lemma-base-change-relative-assassin},
nous voyons que l'image $u \in U$ de $u'$ appartient à
$\text{Ass}(\mathcal{F}|_{U_v})$ où $v \in V$ est l'image de $u$.
Cela signifie à son tour que $g'(x') \in \text{Ass}_{X/Y}(\mathcal{F})$.

\medskip\noindent
Preuve de (2). Choisissons $u \in U$ s'envoyant sur $x$.
Notons $v \in V$ l'image de $u$.
La condition $x \in \text{Ass}_{X/Y}(\mathcal{F})$ équivaut
à $u \in \text{Ass}(\mathcal{F}|_{U_v})$
par définition. Choisissons un point $v' \in V'$ s'envoyant
sur $y' \in |Y'|$ et sur $v \in V$ (ce qui est possible d'après
Propriétés des espaces, lemme \ref{spaces-properties-lemma-points-cartesian}).
Soit $t \in \Spec(\kappa(v') \otimes_{\kappa(v)} \kappa(u))$
un point générique d'une composante irréductible.
Soit $u' \in U'$ l'image de $t$.
En appliquant Diviseurs, lemme \ref{divisors-lemma-base-change-relative-assassin},
nous voyons que $u' \in \text{Ass}(\mathcal{F}'|_{U'_{v'}})$.
Cela signifie à son tour que $x' \in \text{Ass}_{X'/Y'}(\mathcal{F}')$
où $x' \in |X'|$ est l'image de $u'$.
\end{proof}

\begin{lemma}
\label{lemma-base-change-relative-assassin-quasi-finite}
Conservons les notations et les hypothèses du
lemme \ref{lemma-base-change-relative-assassin}.
Supposons que $g$ soit localement quasi-fini, ou plus généralement que
pour tout $y' \in |Y'|$ le degré de transcendance de $y'/g(y')$ soit $0$.
Alors $\text{Ass}_{X'/Y'}(\mathcal{F}')$ est l'image réciproque de
$\text{Ass}_{X/Y}(\mathcal{F})$.
\end{lemma}

\begin{proof}
Le degré de transcendance d'un point sur son image est défini dans
Morphismes d'espaces, définition
\ref{spaces-morphisms-definition-dimension-fibre}.
Soit $x' \in |X'|$ d'image $x \in |X|$.
Choisissons un schéma $V$ et un morphisme \'etale surjectif $V \to Y$.
Choisissons un schéma $U$ et un morphisme
\'etale surjectif $U \to V \times_Y X$. Choisissons un schéma $V'$
et un morphisme \'etale surjectif $V' \to V \times_Y Y'$.
Alors $U' = V' \times_V U$ est un schéma et le morphisme
$U' \to X'$ est surjectif et \'etale.
Choisissons $u \in U$ s'envoyant sur $x$.
Notons $v \in V$ l'image de $u$.
La condition $x \in \text{Ass}_{X/Y}(\mathcal{F})$ équivaut
à $u \in \text{Ass}(\mathcal{F}|_{U_v})$
par définition. Choisissons un point $u' \in U'$ s'envoyant
sur $x' \in |X'|$ et sur $u \in U$ (ce qui est possible d'après
Propriétés des espaces, lemme \ref{spaces-properties-lemma-points-cartesian}).
Soit $v' \in V'$ l'image de $u'$.
La condition $x' \in \text{Ass}_{X'/Y'}(\mathcal{F}')$ équivaut
à $u' \in \text{Ass}(\mathcal{F}'|_{U'_{v'}})$
par définition.
Le lemme résulte alors de la discussion de
Diviseurs, remarque \ref{divisors-remark-base-change-relative-assassin},
appliquée à $u' \in \Spec(\kappa(v') \otimes_{\kappa(v)} \kappa(u))$.
\end{proof}

\begin{lemma}
\label{lemma-relative-weak-assassin-finite}
Soit $S$ un schéma.
Soit $f : X \to Y$ un morphisme d'espaces algébriques sur $S$.
Soit $i : Z \to X$ un morphisme fini.
Soit $\mathcal{G}$ un $\mathcal{O}_Z$-module quasi-cohérent.
Alors $\text{WeakAss}_{X/Y}(i_*\mathcal{G}) =
i(\text{WeakAss}_{Z/Y}(\mathcal{G}))$.
\end{lemma}

\begin{proof}
Cela résulte du cas des schémas
(Diviseurs, lemme \ref{divisors-lemma-relative-weak-assassin-finite})
par localisation \'etale. Nous omettons les détails.
\end{proof}

\begin{lemma}
\label{lemma-relative-assassin-constructible}
Soit $Y$ un schéma. Soit $X$ un espace algébrique de présentation finie
sur $Y$. Soit $\mathcal{F}$ un $\mathcal{O}_X$-module quasi-cohérent
de présentation finie. Soit $U \subset X$ un sous-espace ouvert
tel que $U \to Y$ soit quasi-compact. Alors l'ensemble
$$
E = \{y \in Y \mid \text{Ass}_{X_y}(\mathcal{F}_y) \subset |U_y|\}
$$
est une partie localement constructible de $Y$.
\end{lemma}

\begin{proof}
Remarquons que, puisque $Y$ est un schéma, on peut considérer les fibres
$X_y = \Spec(\kappa(y)) \times_Y X$. (De plus, d'après nos définitions,
l'ensemble $\text{Ass}_{X_y}(\mathcal{F}_y)$ est exactement la fibre de
$\text{Ass}_{X/Y}(\mathcal{F}) \to Y$ au-dessus de $y$, mais nous ne l'utiliserons pas.)
La question est locale sur $Y$ : en effet, il suffit de montrer que
$E$ est constructible lorsque $Y$ est affine.
Dans ce cas, $X$ est quasi-compact. Choisissons un schéma affine $W$
et un morphisme \'etale surjectif $\varphi : W \to X$.
Alors $\text{Ass}_{X_y}(\mathcal{F}_y)$ est l'image de
$\text{Ass}_{W_y}(\varphi^*\mathcal{F}_y)$ pour tous les $y \in Y$.
Par conséquent, le lemme résulte du cas des schémas, appliqué à
l'ouvert $\varphi^{-1}(U) \subset W$ et au morphisme $W \to Y$.
Le cas des schémas est traité dans
Compléments sur les morphismes, lemme
\ref{more-morphisms-lemma-relative-assassin-constructible}.
\end{proof}









\section{Idéaux de Fitting}
\label{section-fitting-ideals}

\noindent
Cette section prolonge la discussion de
Diviseurs, section \ref{divisors-section-fitting-ideals}.
Soit $S$ un schéma. Soit $X$ un espace algébrique sur $S$.
Soit $\mathcal{F}$ un $\mathcal{O}_X$-module quasi-cohérent de type fini.
Nous pouvons alors construire les idéaux de Fitting
$$
0 = \text{Fit}_{-1}(\mathcal{F}) \subset \text{Fit}_0(\mathcal{F}) \subset
\text{Fit}_1(\mathcal{F}) \subset \ldots \subset \mathcal{O}_X
$$
comme la suite d'idéaux quasi-cohérents caractérisée par la
propriété suivante : pour tout affine $U = \Spec(A)$ \'etale sur $X$,
si $\mathcal{F}|_U$ correspond au $A$-module $M$, alors
$\text{Fit}_i(\mathcal{F})|_U$
correspond à l'idéal $\text{Fit}_i(M) \subset A$.
Cette définition est indépendante des choix et donne un idéal quasi-cohérent, car
si $A \to B$ est un homomorphisme d'anneaux \'etale, alors le $i$-ième idéal de Fitting
de $M \otimes_A B$ sur $B$ est égal à $\text{Fit}_i(M) B$ d'après
Compléments sur l'algèbre, lemme \ref{more-algebra-lemma-fitting-ideal-basics}, partie (3).
Plus précisément (peut-être), l'existence du faisceau quasi-cohérent d'idéaux
$\text{Fit}_0(\mathcal{O}_X)$ résulte (par exemple) de
la description des faisceaux quasi-cohérents dans
Propriétés des espaces, lemme
\ref{spaces-properties-lemma-characterize-quasi-coherent-small-etale}
et de la propriété d'image inverse donnée dans
Diviseurs, lemme \ref{divisors-lemma-base-change-fitting-ideal}.

\medskip\noindent
L'avantage de cette construction des idéaux de Fitting
est que l'on voit immédiatement que leur formation
commute à la localisation \'etale ; par conséquent, de nombreuses propriétés des
idéaux de Fitting se ramènent immédiatement aux propriétés correspondantes
dans le cas des schémas. Nous utiliserons souvent la
discussion de Propriétés des espaces, section
\ref{spaces-properties-section-properties-modules}
pour passer des propriétés des faisceaux quasi-cohérents
sur les schémas à celles sur les espaces algébriques.

\begin{lemma}
\label{lemma-base-change-fitting-ideal}
Soit $S$ un schéma.
Soit $f : X \to Y$ un morphisme d'espaces algébriques sur $S$.
Soit $\mathcal{F}$ un $\mathcal{O}_Y$-module quasi-cohérent de type fini.
Alors
$f^{-1}\text{Fit}_i(\mathcal{F}) \cdot \mathcal{O}_X =
\text{Fit}_i(f^*\mathcal{F})$.
\end{lemma}

\begin{proof}
L'énoncé se ramène à
Diviseurs, lemme \ref{divisors-lemma-base-change-fitting-ideal}
par localisation \'etale.
\end{proof}

\begin{lemma}
\label{lemma-fitting-ideal-of-finitely-presented}
Soit $S$ un schéma. Soit $X$ un espace algébrique sur $S$.
Soit $\mathcal{F}$ un $\mathcal{O}_X$-module de présentation finie.
Alors $\text{Fit}_r(\mathcal{F})$ est un idéal quasi-cohérent de type fini.
\end{lemma}

\begin{proof}
L'énoncé se ramène à
Diviseurs, lemme \ref{divisors-lemma-fitting-ideal-of-finitely-presented}
par localisation \'etale.
\end{proof}

\begin{lemma}
\label{lemma-on-subscheme-cut-out-by-Fit-0}
Soit $S$ un schéma. Soit $X$ un espace algébrique sur $S$.
Soit $\mathcal{F}$ un $\mathcal{O}_X$-module quasi-cohérent de type fini.
Soit $Z_0 \subset X$ le sous-espace fermé défini par
$\text{Fit}_0(\mathcal{F})$.
Soit $Z \subset X$ le support schématique de $\mathcal{F}$.
Alors
\begin{enumerate}
\item $Z \subset Z_0 \subset X$ comme sous-espaces fermés,
\item $|Z| = |Z_0| = \text{Supp}(\mathcal{F})$ comme parties fermées de $|X|$,
\item il existe un $\mathcal{O}_{Z_0}$-module quasi-cohérent de type fini
$\mathcal{G}_0$ avec
$$
(Z_0 \to X)_*\mathcal{G}_0 = \mathcal{F}.
$$
\end{enumerate}
\end{lemma}

\begin{proof}
Rappelons que la formation de $Z$ commute à la localisation \'etale ; voir
Morphismes d'espaces, définition
\ref{spaces-morphisms-definition-scheme-theoretic-support}
(qui utilise Morphismes d'espaces, lemme
\ref{spaces-morphisms-lemma-scheme-theoretic-support}
pour définir $Z$). Ainsi (1) et (2) découlent du cas des schémas ; voir
Diviseurs, lemme \ref{divisors-lemma-on-subscheme-cut-out-by-Fit-0}.
Pour obtenir $\mathcal{G}_0$ comme dans la partie (3), nous pouvons utiliser
le module $\mathcal{G}$ sur $Z$ fourni par Morphismes d'espaces, lemme
\ref{spaces-morphisms-lemma-scheme-theoretic-support}
et poser $\mathcal{G}_0 = (Z \to Z_0)_*\mathcal{G}$.
\end{proof}

\begin{lemma}
\label{lemma-fitting-ideal-generate-locally}
Soit $S$ un schéma. Soit $X$ un espace algébrique sur $S$.
Soit $\mathcal{F}$ un module quasi-cohérent de type fini
sur $\mathcal{O}_X$. Soit $x \in |X|$. Alors $\mathcal{F}$ peut être
engendré par $r$ éléments dans un voisinage \'etale de $x$ si et seulement
si $\text{Fit}_r(\mathcal{F})_{\overline{x}} = \mathcal{O}_{X, \overline{x}}$.
\end{lemma}

\begin{proof}
L'énoncé se ramène à
Diviseurs, lemme \ref{divisors-lemma-fitting-ideal-generate-locally}
par localisation \'etale (ainsi qu'à la description de l'anneau
local dans Propriétés des espaces, section
\ref{spaces-properties-section-stalks-structure-sheaf}
et au fait que le hensélisé strict d'un anneau local
est fidèlement plat sur celui-ci, ce qui montre que l'égalité sur le hensélisé
strict équivaut à l'égalité sur l'anneau local).
\end{proof}

\begin{lemma}
\label{lemma-fitting-ideal-finite-locally-free}
Soit $S$ un schéma. Soit $X$ un espace algébrique sur $S$.
Soit $\mathcal{F}$ un module quasi-cohérent de type fini
sur $\mathcal{O}_X$. Soit $r \geq 0$. Les conditions suivantes sont équivalentes :
\begin{enumerate}
\item $\mathcal{F}$ est fini localement libre de rang $r$,
\item $\text{Fit}_{r - 1}(\mathcal{F}) = 0$ et
$\text{Fit}_r(\mathcal{F}) = \mathcal{O}_X$, et
\item $\text{Fit}_k(\mathcal{F}) = 0$ pour $k < r$ et
$\text{Fit}_k(\mathcal{F}) = \mathcal{O}_X$ pour $k \geq r$.
\end{enumerate}
\end{lemma}

\begin{proof}
L'énoncé se ramène à
Diviseurs, lemme \ref{divisors-lemma-fitting-ideal-finite-locally-free}
par localisation \'etale.
\end{proof}

\begin{lemma}
\label{lemma-locally-free-rank-r-pullback}
Soit $S$ un schéma. Soit $X$ un espace algébrique sur $S$.
Soit $\mathcal{F}$ un module quasi-cohérent de type fini
sur $\mathcal{O}_X$. Les sous-espaces fermés
$$
X = Z_{-1} \supset Z_0 \supset Z_1 \supset Z_2 \ldots
$$
définis par les idéaux de Fitting de $\mathcal{F}$ ont les propriétés
suivantes :
\begin{enumerate}
\item L'intersection $\bigcap Z_r$ est vide.
\item Le foncteur $(\Sch/X)^{opp} \to \textit{Ens}$ défini par la règle
$$
T \longmapsto
\left\{
\begin{matrix}
\{*\} & \text{si }\mathcal{F}_T\text{ est engendré localement par }
\leq r\text{ sections} \\
\emptyset & \text{sinon}
\end{matrix}
\right.
$$
est représentable par le sous-espace ouvert $X \setminus Z_r$.
\item Le foncteur $F_r : (\Sch/X)^{opp} \to \textit{Ens}$ défini par la règle
$$
T \longmapsto
\left\{
\begin{matrix}
\{*\} & \text{si }\mathcal{F}_T\text{ est localement libre de rang }r\\
\emptyset & \text{sinon}
\end{matrix}
\right.
$$
est représentable par le sous-espace localement fermé $Z_{r - 1} \setminus Z_r$
de $X$.
\end{enumerate}
Si $\mathcal{F}$ est de présentation finie, alors
$Z_r \to X$, $X \setminus Z_r \to X$ et $Z_{r - 1} \setminus Z_r \to X$
sont de présentation finie.
\end{lemma}

\begin{proof}
L'énoncé se ramène à
Diviseurs, lemme \ref{divisors-lemma-locally-free-rank-r-pullback}
par localisation \'etale.
\end{proof}

\begin{lemma}
\label{lemma-finite-presentation-module}
Soit $S$ un schéma. Soit $X$ un espace algébrique sur $S$.
Soit $\mathcal{F}$ un $\mathcal{O}_X$-module
de présentation finie. Soit $X = Z_{-1} \subset Z_0 \subset Z_1 \subset \ldots$
comme dans le lemme \ref{lemma-locally-free-rank-r-pullback}.
Posons $X_r = Z_{r - 1} \setminus Z_r$.
Alors $X' = \coprod_{r \geq 0} X_r$ représente le foncteur
$$
F_{flat} : \Sch/X \longrightarrow \textit{Ens},\quad\quad
T \longmapsto
\left\{
\begin{matrix}
\{*\} & \text{si }\mathcal{F}_T\text{ est plat sur }T\\
\emptyset & \text{sinon}
\end{matrix}
\right.
$$
De plus, $\mathcal{F}|_{X_r}$ est localement libre de rang $r$ et les
morphismes $X_r \to X$ et $X' \to X$ sont de présentation finie.
\end{lemma}

\begin{proof}
L'énoncé se ramène à
Diviseurs, lemme \ref{divisors-lemma-finite-presentation-module}
par localisation \'etale.
\end{proof}














\section{Diviseurs de Cartier effectifs}
\label{section-effective-Cartier-divisors}

\noindent
Pour des raisons de commodité, il semble préférable de définir d'abord la notion de
diviseur de Cartier effectif. Rappelons que dans
Morphismes d'espaces, section \ref{spaces-morphisms-section-closed-immersions},
nous avons étudié la correspondance entre les sous-espaces fermés et les faisceaux
quasi-cohérents d'idéaux. De plus, dans
Propriétés des espaces, section
\ref{spaces-properties-section-properties-modules}, nous avons étudié les propriétés
des modules quasi-cohérents, notamment celle d'être « engendré localement par $1$ élément ».
Ces références montrent que la définition suivante est
compatible avec celle donnée pour les schémas.

\begin{definition}
\label{definition-effective-Cartier-divisor}
Soit $S$ un schéma. Soit $X$ un espace algébrique sur $S$.
\begin{enumerate}
\item Un {\it sous-espace fermé localement principal} de $X$ est un sous-espace fermé
dont le faisceau d'idéaux est engendré localement par $1$ élément.
\item Un {\it diviseur de Cartier effectif} sur $X$ est un sous-espace fermé
$D \subset X$ tel que le faisceau d'idéaux $\mathcal{I}_D \subset \mathcal{O}_X$
soit un $\mathcal{O}_X$-module inversible.
\end{enumerate}
\end{definition}

\noindent
Ainsi, un diviseur de Cartier effectif est un sous-espace fermé localement principal,
mais la réciproque n'est pas toujours vraie. Les diviseurs de Cartier effectifs sont des
sous-espaces fermés de codimension pure $1$ au sens le plus fort possible :
ils sont en effet définis localement par un seul élément non diviseur de zéro.
Ils sont en particulier nulle part denses.

\begin{lemma}
\label{lemma-characterize-effective-Cartier-divisor}
Soit $S$ un schéma. Soit $X$ un espace algébrique sur $S$.
Soit $D \subset X$ un sous-espace fermé.
Les conditions suivantes sont équivalentes :
\begin{enumerate}
\item Le sous-espace $D$ est un diviseur de Cartier effectif sur $X$.
\item Il existe un schéma $U$ et un morphisme \'etale surjectif $U \to X$ tels que
l'image inverse $D \times_X U$ soit un diviseur de Cartier effectif sur $U$.
\item Pour tout schéma $U$ et tout morphisme \'etale $U \to X$,
l'image inverse $D \times_X U$ est un diviseur de Cartier effectif sur $U$.
\item Pour tout $x \in |D|$, il existe un morphisme \'etale
$(U, u) \to (X, x)$ d'espaces algébriques pointés tel que $U = \Spec(A)$
et $D \times_X U = \Spec(A/(f))$, où $f \in A$ est non diviseur de zéro.
\end{enumerate}
\end{lemma}

\begin{proof}
L'équivalence de (1) -- (3) résulte de la
définition \ref{definition-effective-Cartier-divisor}
et des références qui la précèdent.
Supposons (1) et soit $x \in |D|$. Choisissons un schéma $W$ et un
morphisme \'etale surjectif
$W \to X$. Choisissons $w \in D \times_X W$ d'image $x$.
D'après (3), $D \times_X W$ est un diviseur de Cartier
effectif sur $W$. On obtient donc un voisinage \'etale affine $U$
en choisissant un voisinage ouvert affine de $w$ dans $W$, comme dans
Diviseurs, lemme \ref{divisors-lemma-characterize-effective-Cartier-divisor}.

\medskip\noindent
Supposons (4). Alors $\mathcal{I}_D|_U$ est inversible d'après
Diviseurs, lemme \ref{divisors-lemma-characterize-effective-Cartier-divisor}.
Puisque l'on peut trouver un recouvrement \'etale de $X$ par la famille
de tous ces $U$ et de $X \setminus D$, on en déduit que
$\mathcal{I}_D$ est un $\mathcal{O}_X$-module inversible.
\end{proof}

\begin{lemma}
\label{lemma-complement-locally-principal-closed-subscheme}
Soit $S$ un schéma. Soit $X$ un espace algébrique sur $S$.
Soit $Z \subset X$ un sous-espace fermé localement
principal. Soit $U = X \setminus Z$. Alors $U \to X$ est un morphisme affine.
\end{lemma}

\begin{proof}
La question est locale sur $X$ pour la topologie \'etale ; voir
Morphismes d'espaces, lemmes \ref{spaces-morphisms-lemma-affine-local}
et le
lemme \ref{lemma-characterize-effective-Cartier-divisor}.
L'énoncé découle donc du cas des schémas, traité dans
Diviseurs, lemme
\ref{divisors-lemma-complement-locally-principal-closed-subscheme}.
\end{proof}

\begin{lemma}
\label{lemma-complement-effective-Cartier-divisor}
Soit $S$ un schéma. Soit $X$ un espace algébrique sur $S$.
Soit $D \subset X$ un diviseur de Cartier effectif.
Soit $U = X \setminus D$. Alors $U \to X$ est un morphisme affine et $U$
est schématiquement dense dans $X$.
\end{lemma}

\begin{proof}
L'affinité résulte du lemme \ref{lemma-complement-locally-principal-closed-subscheme}.
La question de la densité est locale sur $X$ pour la topologie \'etale d'après
Morphismes d'espaces, définition
\ref{spaces-morphisms-definition-scheme-theoretically-dense}.
L'énoncé découle donc du cas des schémas, traité dans
Diviseurs, lemme
\ref{divisors-lemma-complement-effective-Cartier-divisor}.
\end{proof}

\begin{lemma}
\label{lemma-effective-Cartier-makes-dimension-drop}
Soit $S$ un schéma. Soit $X$ un espace algébrique sur $S$.
Soit $D \subset X$ un diviseur de Cartier effectif.
Soit $x \in |D|$.
Si $\dim_x(X) < \infty$, alors $\dim_x(D) < \dim_x(X)$.
\end{lemma}

\begin{proof}
La définition d'un diviseur de Cartier effectif et celle de la
dimension d'un espace algébrique en un point
(Propriétés des espaces, définition
\ref{spaces-properties-definition-dimension-at-point})
sont toutes deux locales pour la topologie \'etale. Le lemme découle donc du cas des schémas
traité dans
Diviseurs, lemme \ref{divisors-lemma-effective-Cartier-makes-dimension-drop}.
\end{proof}

\begin{definition}
\label{definition-sum-effective-Cartier-divisors}
Soit $S$ un schéma. Soit $X$ un espace algébrique sur $S$.
Étant donnés deux diviseurs de Cartier effectifs
$D_1$, $D_2$ sur $X$, on définit $D = D_1 + D_2$ comme le
sous-espace fermé de $X$ correspondant au faisceau quasi-cohérent
d'idéaux
$\mathcal{I}_{D_1}\mathcal{I}_{D_2} \subset \mathcal{O}_S$.
On appelle ce sous-espace la {\it somme des diviseurs de Cartier effectifs
$D_1$ et $D_2$}.
\end{definition}

\noindent
Il est clair que l'on peut de même définir la somme $\sum n_iD_i$ à partir
d'un nombre fini de diviseurs de Cartier effectifs $D_i$ sur $X$
et d'entiers positifs ou nuls $n_i$.

\begin{lemma}
\label{lemma-sum-effective-Cartier-divisors}
La somme de deux diviseurs de Cartier effectifs est un diviseur de
Cartier effectif.
\end{lemma}

\begin{proof}
Démonstration omise. Localement pour la topologie \'etale, l'énoncé se ramène au
fait d'algèbre élémentaire suivant : si $f_1, f_2 \in A$ sont non diviseurs de zéro dans un anneau $A$, alors
$f_1f_2 \in A$ est non diviseur de zéro.
\end{proof}

\begin{lemma}
\label{lemma-sum-closed-subschemes-effective-Cartier}
Soit $S$ un schéma. Soit $X$ un espace algébrique sur $S$.
Soient $Z, Y$ deux sous-espaces fermés de $X$
de faisceaux d'idéaux $\mathcal{I}$ et $\mathcal{J}$. Si $\mathcal{I}\mathcal{J}$
définit un diviseur de Cartier effectif $D \subset X$, alors $Z$ et $Y$
sont des diviseurs de Cartier effectifs et $D = Z + Y$.
\end{lemma}

\begin{proof}
D'après le lemme \ref{lemma-characterize-effective-Cartier-divisor},
l'énoncé se ramène au cas des schémas, traité dans
Diviseurs, lemme \ref{divisors-lemma-sum-closed-subschemes-effective-Cartier}.
\end{proof}

\noindent
Rappelons que nous avons défini l'image inverse d'un sous-espace fermé
par un morphisme quelconque d'espaces algébriques dans
Morphismes d'espaces, définition
\ref{spaces-morphisms-definition-inverse-image-closed-subspace}.

\begin{lemma}
\label{lemma-pullback-locally-principal}
Soit $S$ un schéma.
Soit $f : X' \to X$ un morphisme d'espaces algébriques sur $S$.
Soit $Z \subset X$ un sous-espace fermé localement principal.
Alors l'image inverse $f^{-1}(Z)$ est un sous-espace fermé localement
principal de $X'$.
\end{lemma}

\begin{proof}
Démonstration omise.
\end{proof}

\begin{definition}
\label{definition-pullback-effective-Cartier-divisor}
Soit $S$ un schéma.
Soit $f : X' \to X$ un morphisme d'espaces algébriques sur $S$.
Soit $D \subset X$
un diviseur de Cartier effectif. On dit que {\it l'image inverse de
$D$ par $f$ est définie} si le sous-espace fermé $f^{-1}(D) \subset X'$
est un diviseur de Cartier effectif. Dans ce cas, on la note indifféremment
$f^*D$ ou $f^{-1}(D)$ et on l'appelle
{\it l'image inverse du diviseur de Cartier effectif}.
\end{definition}

\noindent
La condition selon laquelle $f^{-1}(D)$ est un diviseur de Cartier effectif
est souvent satisfaite en pratique.

\begin{lemma}
\label{lemma-pullback-effective-Cartier-defined}
Soit $S$ un schéma.
Soit $f : X \to Y$ un morphisme d'espaces algébriques sur $S$.
Soit $D \subset Y$ un diviseur de Cartier effectif.
L'image inverse de $D$ par $f$ est définie dans chacun des cas suivants :
\begin{enumerate}
\item $f(x) \not \in |D|$ pour tout point faiblement associé $x$ de $X$,
\item $f$ est plat, et
\item ajouter ici d'autres cas au besoin.
\end{enumerate}
\end{lemma}

\begin{proof}
En travaillant localement pour la topologie \'etale, ce lemme se ramène au cas des schémas ; voir
Diviseurs, lemme \ref{divisors-lemma-pullback-effective-Cartier-defined}.
\end{proof}

\begin{lemma}
\label{lemma-pullback-effective-Cartier-divisors-additive}
Soit $S$ un schéma.
Soit $f : X' \to X$ un morphisme d'espaces algébriques sur $S$.
Soient $D_1$, $D_2$ des diviseurs de Cartier effectifs sur $X$.
Si les images inverses de $D_1$ et $D_2$ sont définies, alors celle
de $D = D_1 + D_2$ est définie et
$f^*D = f^*D_1 + f^*D_2$.
\end{lemma}

\begin{proof}
Démonstration omise.
\end{proof}




\section{Diviseurs de Cartier effectifs et faisceaux inversibles}
\label{section-effective-Cartier-invertible}

\noindent
Puisqu'un diviseur de Cartier effectif possède un faisceau d'idéaux inversible
(définition \ref{definition-effective-Cartier-divisor}), la
définition suivante a un sens.

\begin{definition}
\label{definition-invertible-sheaf-effective-Cartier-divisor}
Soit $S$ un schéma. Soit $X$ un espace algébrique sur $S$
et soit $D \subset X$ un diviseur de Cartier effectif dont le faisceau
d'idéaux est $\mathcal{I}_D$.
\begin{enumerate}
\item Le {\it faisceau inversible $\mathcal{O}_X(D)$ associé à $D$}
est défini par
$$
\mathcal{O}_X(D) =
\SheafHom_{\mathcal{O}_X}(\mathcal{I}_D, \mathcal{O}_X) =
\mathcal{I}_D^{\otimes -1}.
$$
\item La section canonique, habituellement notée $1$ ou $1_D$, est la
section globale de $\mathcal{O}_X(D)$ correspondant au
morphisme d'inclusion $\mathcal{I}_D \to \mathcal{O}_X$.
\item On écrit
$\mathcal{O}_X(-D) = \mathcal{O}_X(D)^{\otimes -1} = \mathcal{I}_D$.
\item Étant donné un second diviseur de Cartier effectif $D' \subset X$, on définit
$\mathcal{O}_X(D - D') =
\mathcal{O}_X(D) \otimes_{\mathcal{O}_X} \mathcal{O}_X(-D')$.
\end{enumerate}
\end{definition}

\noindent
Quelques remarques. Nous verrons plus bas que l'application
$D \mapsto \mathcal{O}_X(D)$ transforme l'addition des diviseurs de Cartier
effectifs (définition \ref{definition-sum-effective-Cartier-divisors})
en l'addition dans le groupe de Picard de $X$
(lemme \ref{lemma-invertible-sheaf-sum-effective-Cartier-divisors}).
Cependant, l'expression $D - D'$ de la définition précédente n'a
aucune signification géométrique. Plus précisément, on peut considérer
l'ensemble des diviseurs de Cartier effectifs sur $X$ comme un monoïde commutatif
$\text{EffCart}(X)$ dont l'élément nul est le diviseur de Cartier effectif vide.
L'application $(D, D') \mapsto \mathcal{O}_X(D - D')$ définit alors
un homomorphisme de groupes
$$
\text{EffCart}(X)^{gp} \longrightarrow \Pic(X)
$$
où le membre de gauche est la complétion en groupe de
$\text{EffCart}(X)$. Autrement dit, lorsqu'on écrit $\mathcal{O}_X(D - D')$,
on peut considérer $D - D'$ comme un élément de $\text{EffCart}(X)^{gp}$.

\begin{lemma}
\label{lemma-conormal-effective-Cartier-divisor}
Soit $S$ un schéma. Soit $X$ un espace algébrique sur $S$.
Soit $D \subset X$ un diviseur de Cartier effectif.
Alors le faisceau conormal vérifie $\mathcal{C}_{D/X} = \mathcal{I}_D|D =
\mathcal{O}_X(D)^{\otimes -1}|_D$.
\end{lemma}

\begin{proof}
Démonstration omise.
\end{proof}

\begin{lemma}
\label{lemma-invertible-sheaf-sum-effective-Cartier-divisors}
Soit $S$ un schéma. Soit $X$ un espace algébrique sur $S$.
Soient $D_1$, $D_2$ des diviseurs de Cartier effectifs sur $X$.
Soit $D = D_1 + D_2$.
Il existe alors un unique isomorphisme
$$
\mathcal{O}_X(D_1) \otimes_{\mathcal{O}_X} \mathcal{O}_X(D_2)
\longrightarrow
\mathcal{O}_X(D)
$$
qui envoie $1_{D_1} \otimes 1_{D_2}$ sur $1_D$.
\end{lemma}

\begin{proof}
Démonstration omise.
\end{proof}

\begin{definition}
\label{definition-regular-section}
Soit $S$ un schéma. Soit $X$ un espace algébrique sur $S$.
Soit $\mathcal{L}$ un faisceau inversible sur $X$.
Une section globale $s \in \Gamma(X, \mathcal{L})$ est appelée
{\it section régulière} si l'application $\mathcal{O}_X \to \mathcal{L}$,
$f \mapsto fs$, est injective.
\end{definition}

\begin{lemma}
\label{lemma-regular-section-structure-sheaf}
Soit $S$ un schéma.
Soit $X$ un espace algébrique sur $S$.
Soit $f \in \Gamma(X, \mathcal{O}_X)$.
Les conditions suivantes sont équivalentes :
\begin{enumerate}
\item $f$ est une section régulière, et
\item pour tout $x \in X$, l'image $f \in \mathcal{O}_{X, \overline{x}}$
est non diviseur de zéro.
\item pour tout $U = \Spec(A)$ affine et \'etale sur $X$,
la restriction $f|_U$ est un élément non diviseur de zéro de $A$, et
\item il existe un schéma $U$ et un morphisme \'etale surjectif
$U \to X$ tels que $f|_U$ soit une section régulière de $\mathcal{O}_U$.
\end{enumerate}
\end{lemma}

\begin{proof}
Démonstration omise.
\end{proof}

\noindent
Notons qu'une section globale $s$ d'un $\mathcal{O}_X$-module inversible
$\mathcal{L}$ peut être considérée comme un homomorphisme de $\mathcal{O}_X$-modules
$s : \mathcal{O}_X \to \mathcal{L}$. Son dual est donc un
homomorphisme $s : \mathcal{L}^{\otimes -1} \to \mathcal{O}_X$.
(Voir Modules sur les sites, lemme \ref{sites-modules-lemma-constructions-invertible},
pour le faisceau inversible dual.)

\begin{definition}
\label{definition-zero-scheme-s}
Soit $S$ un schéma. Soit $X$ un espace algébrique sur $S$.
Soit $\mathcal{L}$ un faisceau inversible.
Soit $s \in \Gamma(X, \mathcal{L})$.
Le {\it schéma des zéros} de $s$ est le sous-espace fermé $Z(s) \subset X$
défini par le faisceau quasi-cohérent d'idéaux
$\mathcal{I} \subset \mathcal{O}_X$ qui est l'image de
l'homomorphisme $s : \mathcal{L}^{\otimes -1} \to \mathcal{O}_X$.
\end{definition}

\begin{lemma}
\label{lemma-zero-scheme}
Soit $S$ un schéma. Soit $X$ un espace algébrique sur $S$.
Soit $\mathcal{L}$ un $\mathcal{O}_X$-module inversible.
Soit $s \in \Gamma(X, \mathcal{L})$.
\begin{enumerate}
\item Considérons les immersions fermées $i : Z \to X$ telles que
$i^*s \in \Gamma(Z, i^*\mathcal{L}))$ soit nulle,
ordonnées par inclusion. Le schéma des zéros $Z(s)$ est l'
élément maximal de cet ensemble ordonné.
\item Pour tout morphisme d'espaces algébriques $f : Y \to X$ sur $S$,
on a $f^*s = 0$ dans $\Gamma(Y, f^*\mathcal{L})$ si et seulement si
$f$ se factorise par $Z(s)$.
\item Le schéma des zéros $Z(s)$ est un sous-espace fermé localement principal de $X$.
\item Le schéma des zéros $Z(s)$ est un diviseur de Cartier effectif sur $X$
si et seulement si $s$ est une section régulière de $\mathcal{L}$.
\end{enumerate}
\end{lemma}

\begin{proof}
Démonstration omise.
\end{proof}

\begin{lemma}
\label{lemma-characterize-OD}
Soit $S$ un schéma. Soit $X$ un espace algébrique sur $S$.
\begin{enumerate}
\item Si $D \subset X$ est un diviseur de Cartier effectif, alors
la section canonique $1_D$ de $\mathcal{O}_X(D)$ est régulière.
\item Réciproquement, si $s$ est une section régulière du faisceau
inversible $\mathcal{L}$, alors il existe un unique diviseur de Cartier
effectif $D = Z(s) \subset X$ et un unique isomorphisme
$\mathcal{O}_X(D) \to \mathcal{L}$ qui envoie $1_D$ sur $s$.
\end{enumerate}
Les constructions
$D \mapsto (\mathcal{O}_X(D), 1_D)$ et $(\mathcal{L}, s) \mapsto Z(s)$
définissent des applications réciproques
$$
\left\{
\begin{matrix}
\text{diviseurs de Cartier effectifs sur }X
\end{matrix}
\right\}
\leftrightarrow
\left\{
\begin{matrix}
\text{couples }(\mathcal{L}, s)\text{ formés d'un}\\
\mathcal{O}_X\text{-module inversible et d'une section globale régulière}
\end{matrix}
\right\}
$$
\end{lemma}

\begin{proof}
Démonstration omise.
\end{proof}





\section{Diviseurs de Cartier effectifs sur les espaces noethériens}
\label{section-Noetherian-effective-Cartier}

\noindent
Dans le cadre localement noethérien, l'essentiel de l'étude des
diviseurs de Cartier effectifs et des sections régulières se simplifie quelque peu.

\begin{lemma}
\label{lemma-effective-Cartier-divisor-Sk}
Soient $S$ un schéma et $X$ un espace algébrique localement noethérien
sur $S$. Soit $D \subset X$ un diviseur de Cartier effectif. Si $X$ vérifie
$(S_k)$, alors $D$ vérifie $(S_{k - 1})$.
\end{lemma}

\begin{proof}
D'après notre définition de la propriété $(S_k)$ pour les espaces algébriques
(Propriétés des espaces, section
\ref{spaces-properties-section-types-properties})
et le
lemme \ref{lemma-characterize-effective-Cartier-divisor},
l'énoncé découle du cas des schémas
(Diviseurs, lemme \ref{divisors-lemma-effective-Cartier-divisor-Sk}).
\end{proof}

\begin{lemma}
\label{lemma-normal-effective-Cartier-divisor-S1}
Soient $S$ un schéma et $X$ un espace algébrique normal localement noethérien
sur $S$. Soit $D \subset X$ un
diviseur de Cartier effectif. Alors $D$ vérifie $(S_1)$.
\end{lemma}

\begin{proof}
D'après notre définition de la normalité pour les espaces algébriques
(Propriétés des espaces, section
\ref{spaces-properties-section-types-properties})
et le
lemme \ref{lemma-characterize-effective-Cartier-divisor},
l'énoncé découle du cas des schémas
(Diviseurs, lemme \ref{divisors-lemma-normal-effective-Cartier-divisor-S1}).
\end{proof}

\noindent
Le lemme suivant permet parfois de construire des diviseurs de
Cartier effectifs.

\begin{lemma}
\label{lemma-complement-open-affine-effective-cartier-divisor}
Soit $S$ un schéma. Soit $X$ un espace algébrique régulier, noethérien et séparé
sur $S$. Soit $U \subset X$ un ouvert affine dense. Il existe alors un
diviseur de Cartier effectif $D \subset X$ tel que $U = X \setminus D$.
\end{lemma}

\begin{proof}
Nous affirmons que l'espace algébrique $D$ muni de la structure réduite induite
sur $X \setminus U$ (Propriétés des espaces, définition
\ref{spaces-properties-definition-reduced-induced-space})
est le diviseur de Cartier effectif recherché. La construction
de $D$ commute à la localisation \'etale ; voir la démonstration de
Propriétés des espaces, lemme
\ref{spaces-properties-lemma-reduced-closed-subspace}.
Soit $X' \to X$ un morphisme \'etale surjectif avec $X'$ affine.
Puisque $X$ est séparé, on voit que $U' = X' \times_X U$ est
affine. Puisque l'application $|X'| \to |X|$ est ouverte, on voit que $U'$
est dense dans $X'$. Puisque $D' = X' \times_X D$ est la structure de schéma
réduite induite sur $X' \setminus U'$, on en déduit que
$D'$ est un diviseur de Cartier effectif d'après
Diviseurs, lemme
\ref{divisors-lemma-complement-open-affine-effective-cartier-divisor}
et sa démonstration. C'est ce qu'il fallait démontrer.
\end{proof}

\begin{lemma}
\label{lemma-Noetherian-regular-separated-pic-effective-Cartier}
Soit $S$ un schéma. Soit $X$ un espace algébrique régulier, noethérien et séparé
sur $S$. Alors tout $\mathcal{O}_X$-module inversible est isomorphe à
$$
\mathcal{O}_X(D - D') =
\mathcal{O}_X(D) \otimes_{\mathcal{O}_X} \mathcal{O}_X(D')^{\otimes -1}
$$
pour certains diviseurs de Cartier effectifs $D, D'$ sur $X$.
\end{lemma}

\begin{proof}
Soit $\mathcal{L}$ un $\mathcal{O}_X$-module inversible.
Choisissons un ouvert affine dense $U \subset X$ tel que
$\mathcal{L}|_U$ soit trivial. Cela est possible car
$X$ possède un sous-espace ouvert dense qui est un schéma ; voir
Propriétés des espaces, proposition
\ref{spaces-properties-proposition-locally-quasi-separated-open-dense-scheme}.
Notons $s : \mathcal{O}_U \to \mathcal{L}|_U$ la trivialisation.
Le complémentaire de $U$ est un diviseur de Cartier effectif
$D$. Nous affirmons que, pour un certain $n > 0$, l'homomorphisme $s$ se prolonge de manière unique en un homomorphisme
$$
s : \mathcal{O}_X(-nD) \longrightarrow \mathcal{L}
$$
Cette affirmation implique le lemme car elle montre que
$\mathcal{L} \otimes_{\mathcal{O}_X} \mathcal{O}_X(nD)$
possède une section globale régulière, donc est isomorphe à
$\mathcal{O}_X(D')$ pour un certain diviseur de Cartier effectif $D'$
d'après le lemme \ref{lemma-characterize-OD}.
Pour démontrer l'affirmation, on peut travailler localement pour la topologie \'etale. On peut donc supposer que
$X$ est un schéma affine noethérien. Puisque
$\mathcal{O}_X(-nD) = \mathcal{I}^n$, où $\mathcal{I} = \mathcal{O}_X(-D)$
est le faisceau d'idéaux de $D$ dans $X$, ce cas résulte de
Cohomologie des schémas, lemme \ref{coherent-lemma-homs-over-open}.
\end{proof}

\noindent
Le lemme suivant relève en réalité d'une autre section.

\begin{lemma}
\label{lemma-smooth-over-valuation-ring-effective-Cartier}
Soit $R$ un anneau de valuation de corps des fractions $K$.
Soit $X$ un espace algébrique sur $R$ tel que $X \to \Spec(R)$
soit lisse. Pour tout diviseur de Cartier effectif $D \subset X_K$,
il existe un diviseur de Cartier effectif $D' \subset X$
tel que $D'_K = D$.
\end{lemma}

\begin{proof}
Soit $D' \subset X$ l'image schématique de $D \to X_K \to X$.
Puisque ce morphisme est quasi-compact, la formation de $D'$
commute au changement de base plat ; voir
Morphismes d'espaces, lemme
\ref{spaces-morphisms-lemma-flat-base-change-scheme-theoretic-image}.
On obtient en particulier $D'_K = D$. On peut donc
supposer $X$ affine. Écrivons $X = \Spec(A)$.
Alors $X_K = \Spec(A \otimes_R K)$ et $D$ correspond à
un idéal $I \subset A \otimes_R K$. Il faut montrer que
$J = I \cap A$ définit un diviseur de Cartier effectif dans $X$.
Observons d'abord que $A/J$ est plat sur $R$ (car c'est un
$R$-module sans torsion ; voir Compléments d'algèbre, lemme
\ref{more-algebra-lemma-valuation-ring-torsion-free-flat}),
donc $J$ est de type fini d'après
Compléments d'algèbre, lemme
\ref{more-algebra-lemma-flat-finite-type-valuation-ring-finite-presentation}
et le
lemme d'Algèbre \ref{algebra-lemma-extension}.
Il suffit donc de montrer que $J_\mathfrak q \subset A_\mathfrak q$
est engendré par un seul élément pour tout idéal premier $\mathfrak q \subset A$.
Soit $\mathfrak p = R \cap \mathfrak q$. Alors
$R_\mathfrak p$ est un anneau de valuation
(Algèbre, lemme \ref{algebra-lemma-make-valuation-rings}).
Observons en outre que $A_\mathfrak q/\mathfrak p A_\mathfrak q$
est un anneau régulier d'après Algèbre, lemme
\ref{algebra-lemma-characterize-smooth-over-field}.
On peut donc appliquer Compléments d'algèbre, lemme
\ref{more-algebra-lemma-picard-group-generic-fibre-regular}
pour voir que $I(A_\mathfrak q \otimes_R K)$ est engendré par
un seul élément $f \in A_\mathfrak p \otimes_R K$.
Après avoir chassé les dénominateurs, on peut supposer $f \in A_\mathfrak q$.
Soit $\mathfrak c \subset R_\mathfrak p$ l'idéal de contenu de $f$
(voir Compléments d'algèbre, définition \ref{more-algebra-definition-content-ideal}
et Compléments sur la platitude, lemme
\ref{flat-lemma-flat-finite-type-local-valuation-ring-has-content}).
Puisque $R_\mathfrak p$ est un anneau de valuation et
que $\mathfrak c$ est de type fini
(Compléments d'algèbre, lemme \ref{more-algebra-lemma-content-finitely-generated}),
on a $\mathfrak c = (\pi)$ pour un certain $\pi \in R_\mathfrak p$
(Algèbre, lemme \ref{algebra-lemma-characterize-valuation-ring}).
Après avoir remplacé $f$ par $\pi^{-1}f$, on a $f \in A_\mathfrak q$
et $f \not \in \mathfrak pA_\mathfrak q$.
Affirmation : $I_\mathfrak q = (f)$, ce qui achève la démonstration.
Pour la vérifier, observons que $f \in I_\mathfrak q$.
On a donc une surjection $A_\mathfrak q/(f) \to A_\mathfrak q/I_\mathfrak q$
qui devient un isomorphisme après tensorisation sur $R$ avec $K$.
Il suffit donc de montrer que
$A_\mathfrak q/(f)$ est plat sur $R_\mathfrak p$.
Cela résulte de
Algèbre, lemme \ref{algebra-lemma-grothendieck-general}
et du choix de $f$.
\end{proof}









\section{Diviseurs de Cartier effectifs relatifs}
\label{section-effective-Cartier-morphisms}

\noindent
Le lemme suivant montre qu'un diviseur de Cartier effectif qui est
plat sur la base est véritablement une « famille de diviseurs de Cartier effectifs »
sur la base. Par exemple, sa restriction à toute fibre est un diviseur de
Cartier effectif.

\begin{lemma}
\label{lemma-relative-Cartier}
Soit $S$ un schéma. Soit $f : X \to Y$ un morphisme
d'espaces algébriques sur $S$. Soit $D \subset X$ un sous-espace fermé.
Supposons que
\begin{enumerate}
\item $D$ soit un diviseur de Cartier effectif, et
\item $D \to Y$ soit un morphisme plat.
\end{enumerate}
Alors, pour tout morphisme de schémas $g : Y' \to Y$, l'image inverse
$(g')^{-1}D$ est un diviseur de Cartier effectif sur $X' = Y' \times_Y X$,
où $g' : X' \to X$ est la projection.
\end{lemma}

\begin{proof}
D'après le lemme \ref{lemma-characterize-effective-Cartier-divisor},
la propriété d'être un diviseur de Cartier effectif est locale pour la topologie \'etale.
Le lemme se ramène donc immédiatement au cas des schémas,
qui est Diviseurs, lemme \ref{divisors-lemma-relative-Cartier}.
\end{proof}

\noindent
Ce lemme justifie la définition suivante.

\begin{definition}
\label{definition-relative-effective-Cartier-divisor}
Soit $S$ un schéma.
Soit $f : X \to Y$ un morphisme d'espaces algébriques sur $S$.
Un {\it diviseur de Cartier effectif relatif} sur $X/Y$ est un
diviseur de Cartier effectif $D \subset X$ tel que $D \to Y$
soit un morphisme plat d'espaces algébriques.
\end{definition}








\section{Fonctions et sections méromorphes}
\label{section-meromorphic-functions}

\noindent
Cette section est l'analogue de
Diviseurs, section \ref{divisors-section-meromorphic-functions}.
Attention : il est encore plus facile de commettre des erreurs sur ce sujet dans le
cas des espaces algébriques que dans celui des schémas !

\medskip\noindent
Soit $S$ un schéma. Soit $X$ un espace algébrique sur $S$.
Pour tout schéma $U$ \'etale sur $X$, nous avons défini l'ensemble
$\mathcal{S}(U) \subset \mathcal{O}_X(U)$
des sections régulières de $\mathcal{O}_X$ sur $U$ ; voir la
définition \ref{definition-regular-section}. La restriction
d'une section régulière à $V/U$ \'etale est régulière. Par conséquent,
$\mathcal{S} : U \mapsto \mathcal{S}(U)$ est un sous-faisceau (d'ensembles)
de $\mathcal{O}_X$. On note parfois $\mathcal{S} = \mathcal{S}_X$
pour indiquer la dépendance en $X$.
De plus, $\mathcal{S}(U)$
est une partie multiplicative de l'anneau $\mathcal{O}_X(U)$ pour
tout $U$. On peut donc considérer
le préfaisceau d'anneaux
$$
U \longmapsto \mathcal{S}(U)^{-1} \mathcal{O}_X(U),
$$
sur $X_\etale$ et son faisceau associé ; voir
Modules sur les sites, section \ref{sites-modules-section-localizing-sheaves-rings}.

\begin{definition}
\label{definition-sheaf-meromorphic-functions}
Soit $S$ un schéma. Soit $X$ un espace algébrique sur $S$.
Le {\it faisceau des fonctions méromorphes sur $X$} est
le faisceau {\it $\mathcal{K}_X$} sur $X_\etale$ associé au préfaisceau
affiché ci-dessus. Une {\it fonction méromorphe} sur $X$
est une section globale de $\mathcal{K}_X$.
\end{definition}

\noindent
Puisque tout élément de chaque $\mathcal{S}(U)$ est non diviseur de zéro dans
$\mathcal{O}_X(U)$, on voit que l'homomorphisme naturel de faisceaux
d'anneaux $\mathcal{O}_X \to \mathcal{K}_X$ est injectif.
De plus, par compatibilité du passage au faisceau associé et de la prise des fibres,
on voit que
$$
\mathcal{K}_{X, \overline{x}} =
\mathcal{S}_{\overline{x}}^{-1}\mathcal{O}_{X, \overline{x}}
$$
pour tout point géométrique $\overline{x}$ de $X$. L'ensemble
$\mathcal{S}_{\overline{x}}$ est une partie de
l'ensemble des éléments non diviseurs de zéro de $\mathcal{O}_{X, \overline{x}}$, mais
ne lui est pas égal en général.

\begin{lemma}
\label{lemma-describe-S}
Soit $S$ un schéma. Soit $X$ un espace algébrique sur $S$.
Pour $U$ affine et \'etale sur $X$, l'ensemble
$\mathcal{S}_X(U)$ est l'ensemble des éléments non diviseurs de zéro de
$\mathcal{O}_X(U)$.
\end{lemma}

\begin{proof}
Cela résulte du lemme \ref{lemma-regular-section-structure-sheaf}.
\end{proof}

\noindent
Soit ensuite $\mathcal{F}$ un faisceau de $\mathcal{O}_X$-modules
sur $X_\etale$.
Considérons le préfaisceau $U \mapsto \mathcal{S}(U)^{-1}\mathcal{F}(U)$.
Le faisceau qui lui est associé est
$\mathcal{F} \otimes_{\mathcal{O}_X} \mathcal{K}_X$ ; voir
Modules sur les sites, lemme \ref{sites-modules-lemma-simple-invert-module}.

\begin{definition}
\label{definition-meromorphic-section}
Soit $S$ un schéma. Soit $X$ un espace algébrique sur $S$.
Soit $\mathcal{F}$ un faisceau de $\mathcal{O}_X$-modules
sur $X_\etale$.
\begin{enumerate}
\item On note $\mathcal{K}_X(\mathcal{F})$ le faisceau de
$\mathcal{K}_X$-modules associé au préfaisceau
$U \mapsto \mathcal{S}(U)^{-1}\mathcal{F}(U)$. De manière équivalente,
$\mathcal{K}_X(\mathcal{F}) =
\mathcal{F} \otimes_{\mathcal{O}_X} \mathcal{K}_X$ (voir ci-dessus).
\item Une {\it section méromorphe de $\mathcal{F}$}
est une section globale de $\mathcal{K}_X(\mathcal{F})$.
\end{enumerate}
\end{definition}

\noindent
On a en particulier
$$
\mathcal{K}_X(\mathcal{F})_{\overline{x}}
=
\mathcal{F}_{\overline{x}}
\otimes_{\mathcal{O}_{X, \overline{x}}} \mathcal{K}_{X, \overline{x}}
=
\mathcal{S}_{\overline{x}}^{-1}\mathcal{F}_{\overline{x}}
$$
pour tout point géométrique $\overline{x}$ de $X$. Il faut toutefois prendre garde
au fait que $\mathcal{S}_{\overline{x}}$ peut ne pas être l'ensemble des
éléments non diviseurs de zéro de l'anneau local \'etale $\mathcal{O}_{X, \overline{x}}$,
comme nous l'avons signalé plus haut. Les faisceaux de sections méromorphes ne sont pas
des modules quasi-cohérents en général, mais ils ont certaines propriétés communes
avec les modules quasi-cohérents.

\begin{lemma}
\label{lemma-meromorphic-quasi-coherent}
Soit $S$ un schéma. Soit $X$ un espace algébrique sur $S$. Supposons que
\begin{enumerate}
\item[(a)] tout point faiblement associé de $X$
soit un point de codimension $0$, et
\item[(b)] $X$ vérifie les conditions équivalentes
du lemme de Morphismes d'espaces
\ref{spaces-morphisms-lemma-prepare-normalization}.
\end{enumerate}
Alors
\begin{enumerate}
\item $\mathcal{K}_X$ est un faisceau quasi-cohérent de $\mathcal{O}_X$-algèbres,
\item pour $U \in X_\etale$ affine, $\mathcal{K}_X(U)$
est l'anneau total des fractions de $\mathcal{O}_X(U)$,
\item pour un point géométrique $\overline{x}$, l'ensemble
$\mathcal{S}_{\overline{x}}$
est l'ensemble des éléments non diviseurs de zéro de $\mathcal{O}_{X, \overline{x}}$, et
\item pour un point géométrique $\overline{x}$, l'anneau
$\mathcal{K}_{X, \overline{x}}$ est l'anneau total des fractions de
$\mathcal{O}_{X, \overline{x}}$.
\end{enumerate}
\end{lemma}

\begin{proof}
Le lemme \ref{lemma-regular-section-structure-sheaf}
montre que, pour $U \in X_\etale$ affine,
$\mathcal{S}_X(U) \subset \mathcal{O}_X(U)$
est l'ensemble des éléments non diviseurs de zéro de $\mathcal{O}_X(U)$.
Ainsi, le préfaisceau $\mathcal{S}^{-1}\mathcal{O}_X$ est égal à
$$
U \longmapsto Q(\mathcal{O}_X(U))
$$
sur $X_{affine, \etale}$, avec les notations de
Algèbre, exemple \ref{algebra-example-localize-at-prime}.
Observons que les points de codimension $0$ de $X$ correspondent
aux points génériques de $U$ ; voir
Propriétés des espaces, lemme \ref{spaces-properties-lemma-codimension-0-points}.
Ainsi, si $U = \Spec(A)$, alors $A$ est un anneau ayant un nombre fini d'idéaux premiers minimaux
tel que tout idéal premier faiblement associé de $A$ soit minimal.
Il en est de même de toute extension \'etale de $A$ (car son
spectre est un schéma affine \'etale sur $X$ et peut donc
jouer le rôle de $A$ dans la phrase précédente).
Pour montrer que notre préfaisceau est un faisceau quasi-cohérent,
il suffit de montrer que
$$
Q(A) \otimes_A B \longrightarrow Q(B)
$$
est un isomorphisme lorsque $A \to B$ est un homomorphisme \'etale d'anneaux ; voir
Propriétés des espaces, lemme
\ref{spaces-properties-lemma-characterize-quasi-coherent-small-etale}.
(Pour définir la flèche affichée, observons que, puisque $A \to B$
est plat, il envoie les éléments non diviseurs de zéro sur des éléments non diviseurs de zéro.)
D'après les lemmes d'Algèbre
\ref{algebra-lemma-total-ring-fractions-no-embedded-points} et
\ref{algebra-lemma-weakly-ass-zero-divisors}.
on a
$$
Q(A) = \prod\nolimits_{\mathfrak p \subset A\text{ minimal}} A_\mathfrak p
\quad\text{et}\quad
Q(B) = \prod\nolimits_{\mathfrak q \subset B\text{ minimal}} B_\mathfrak q
$$
Puisque $A \to B$ est \'etale, les idéaux premiers minimaux de $B$ sont exactement les
idéaux premiers de $B$ au-dessus des idéaux premiers minimaux
de $A$ (par exemple d'après Compléments d'algèbre, lemme
\ref{more-algebra-lemma-dimension-etale-extension}).
D'après les lemmes d'Algèbre \ref{algebra-lemma-local-dimension-zero-henselian},
\ref{algebra-lemma-characterize-henselian} (13), et
\ref{algebra-lemma-mop-up}, on voit que $A_\mathfrak p \otimes_A B$
est un produit fini d'anneaux locaux finis \'etales sur $A_\mathfrak p$.
Cela implique clairement que $A_\mathfrak p \otimes_A B =
\prod_{\mathfrak q\text{ au-dessus de }\mathfrak p} B_\mathfrak q$,
comme souhaité.

\medskip\noindent
À ce stade, on sait que (1) et (2) sont satisfaites.
Démonstration de (3). Soit $s \in \mathcal{O}_{X, \overline{x}}$
un élément non diviseur de zéro. On peut trouver un voisinage \'etale
$(U, \overline{u}) \to (X, \overline{x})$
et $f \in \mathcal{O}_X(U)$ d'image $s$.
Soit $u \in U$ le point déterminé par $\overline{u}$.
Puisque $\mathcal{O}_{U, u} \to \mathcal{O}_{X, \overline{x}}$
est fidèlement plat (car c'est un hensélisé strict), on voit que
$f$ s'envoie sur un élément non diviseur de zéro de $\mathcal{O}_{U, u}$.
D'après Diviseurs, lemme \ref{divisors-lemma-meromorphic-weakass-finite},
quitte à rétrécir $U$, on obtient que $f$ est non diviseur de zéro
et donc une section de $\mathcal{S}_X(U)$.
L'assertion (4) résulte de (3) par calcul des fibres.
\end{proof}

\begin{lemma}
\label{lemma-meromorphic-quasi-coherent-agree}
Soit $S$ un schéma. Soit $X$ un espace algébrique sur $S$. Supposons que
\begin{enumerate}
\item[(a)] tout point faiblement associé de $X$
soit un point de codimension $0$, et
\item[(b)] $X$ vérifie les conditions équivalentes
du lemme de Morphismes d'espaces
\ref{spaces-morphisms-lemma-prepare-normalization}.
\item[(c)] $X$ soit représentable par un schéma $X_0$
(notation maladroite mais provisoire).
\end{enumerate}
Alors le faisceau des fonctions méromorphes $\mathcal{K}_X$
est le faisceau quasi-cohérent de $\mathcal{O}_X$-algèbres
associé au faisceau quasi-cohérent des fonctions méromorphes
$\mathcal{K}_{X_0}$.
\end{lemma}

\begin{proof}
Pour l'équivalence entre $\QCoh(\mathcal{O}_X)$ et
$\QCoh(\mathcal{O}_{X_0})$, voir
Propriétés des espaces, section \ref{spaces-properties-section-quasi-coherent}.
Le lemme est vrai parce que $\mathcal{K}_X$ et $\mathcal{K}_{X_0}$
sont quasi-cohérents et prennent la même valeur sur les ouverts affines correspondants
de $X$ et $X_0$, d'après le lemme \ref{lemma-meromorphic-quasi-coherent} et
Diviseurs, lemme \ref{divisors-lemma-meromorphic-weakass-finite}.
\end{proof}

\begin{definition}
\label{definition-pullback-meromorphic-sections}
Soit $S$ un schéma. Soit $f : X \to Y$ un morphisme
d'espaces algébriques sur $S$. On dit que {\it les images inverses des fonctions
méromorphes sont définies pour $f$} si, pour tout diagramme commutatif
$$
\xymatrix{
U \ar[r] \ar[d] & X \ar[d] \\
V \ar[r] & Y
}
$$
où $U \in X_\etale$ et $V \in Y_\etale$, et pour toute
section $s \in \mathcal{S}_Y(V)$, l'image inverse
$f^\sharp(s) \in \mathcal{O}_X(U)$ est un élément
de $\mathcal{S}_X(U)$.
\end{definition}

\noindent
Dans ce cas, il existe un morphisme induit
$f^\sharp : f_{small}^{-1}\mathcal{K}_Y \to \mathcal{K}_X$,
autrement dit, on obtient un diagramme commutatif de morphismes
de topos annelés
$$
\xymatrix{
(\Sh(X_\etale), \mathcal{K}_X) \ar[r] \ar[d]^{f_{small}} &
(\Sh(X_\etale), \mathcal{O}_X) \ar[d]^{f_{small}} \\
(\Sh(Y_\etale), \mathcal{K}_Y) \ar[r] &
(\Sh(Y_\etale), \mathcal{O}_Y)
}
$$
On note parfois $f^*(s) = f^\sharp(s)$ pour une
section $s \in \Gamma(Y, \mathcal{K}_Y)$.

\begin{lemma}
\label{lemma-pullback-meromorphic-sections-defined}
Soit $S$ un schéma. Soit $f : X \to Y$ un morphisme d'espaces algébriques
sur $S$. Les images inverses des sections méromorphes sont définies
dans chacun des cas suivants :
\begin{enumerate}
\item les points faiblement associés de $X$ sont envoyés
sur des points de codimension $0$ de $Y$,
\item $f$ est plat,
\item ajouter ici d'autres cas au besoin.
\end{enumerate}
\end{lemma}

\begin{proof}
En travaillant localement pour la topologie \'etale, cela se ramène au cas des schémas ; voir
Diviseurs, lemme \ref{divisors-lemma-pullback-meromorphic-sections-defined}.
Pour effectuer cette réduction, utiliser le lemme \ref{lemma-regular-section-structure-sheaf}
(description des sections régulières),
la définition \ref{definition-weakly-associated} (définition
des points faiblement associés), et
Propriétés des espaces, lemme \ref{spaces-properties-lemma-codimension-0-points}
(description des points de codimension $0$).
\end{proof}

\begin{lemma}
\label{lemma-compute-meromorphic}
Soit $S$ un schéma. Soit $X$ un espace algébrique sur $S$. Supposons que
\begin{enumerate}
\item[(a)] tout point faiblement associé de $X$
soit un point de codimension $0$, et
\item[(b)] $X$ vérifie les conditions équivalentes
du lemme de Morphismes d'espaces
\ref{spaces-morphisms-lemma-prepare-normalization},
\item[(c)] tout point de codimension $0$ de $X$ puisse être représenté
par un monomorphisme $\Spec(k) \to X$.
\end{enumerate}
Soit $X^0 \subset |X|$ l'ensemble des points de codimension $0$ de $X$.
Alors on a
$$
\mathcal{K}_X =
\bigoplus\nolimits_{\eta \in X^0} j_{\eta, *}\mathcal{O}_{X, \eta} =
\prod\nolimits_{\eta \in X^0} j_{\eta, *}\mathcal{O}_{X, \eta}
$$
où $j_\eta : \Spec(\mathcal{O}_{X, \eta}) \to X$ est le morphisme canonique
de Schémas, section \ref{schemes-section-points} ; cela a un sens parce que
$X^0$ est contenu dans le lieu schématique de $X$. De même,
pour tout $\mathcal{O}_X$-module quasi-cohérent $\mathcal{F}$,
on obtient la formule
$$
\mathcal{K}_X(\mathcal{F}) =
\bigoplus\nolimits_{\eta \in X^0} j_{\eta, *}\mathcal{F}_\eta =
\prod\nolimits_{\eta \in X^0} j_{\eta, *}\mathcal{F}_\eta
$$
pour le faisceau des sections méromorphes de $\mathcal{F}$.
Enfin, l'anneau des fonctions rationnelles de $X$ est l'anneau des fonctions
méromorphes sur $X$, autrement dit : $R(X) = \Gamma(X, \mathcal{K}_X)$.
\end{lemma}

\begin{proof}
D'après Espaces décents, lemme \ref{decent-spaces-lemma-get-reasonable} et
section \ref{decent-spaces-section-reasonable-decent},
on voit que $X$ est décent\footnote{Réciproquement, si $X$ est décent,
alors la condition (c) est automatiquement satisfaite.}.
Ainsi, $X^0 \subset |X|$ est l'ensemble des points génériques des composantes
irréductibles
(Espaces décents, lemme \ref{decent-spaces-lemma-decent-generic-points}),
et $X^0$ est localement fini dans $|X|$ d'après (b).
Il en résulte que $X^0$ est contenu dans tout ouvert dense de $|X|$.
En particulier, $X^0$ est contenu dans le lieu schématique
(Espaces décents, théorème \ref{decent-spaces-theorem-decent-open-dense-scheme}).
Ainsi, les anneaux locaux $\mathcal{O}_{X, \eta}$ et les morphismes $j_\eta$
sont définis.

\medskip\noindent
Observons qu'une somme directe localement finie de faisceaux de modules
est égale au produit. Ce fait, joint au fait que $X^0$ est localement
fini dans $|X|$, explique les égalités entre
les sommes directes et les produits dans l'énoncé.
Ensuite, puisque $\mathcal{K}_X(\mathcal{F}) =
\mathcal{F} \otimes_{\mathcal{O}_X} \mathcal{K}_X$,
on voit que la seconde égalité résulte de la première.

\medskip\noindent
Soit
$j : Y = \coprod\nolimits_{\eta \in X^0} \Spec(\mathcal{O}_{X, \eta}) \to X$
le morphisme induit par les morphismes $j_\eta$.
Il faut montrer que $\mathcal{K}_X = j_*\mathcal{O}_Y$.
Observons que $\mathcal{K}_Y = \mathcal{O}_Y$, car $Y$ est une somme
disjointe de spectres d'anneaux locaux de dimension $0$ : dans un anneau local
de dimension zéro, tout élément non diviseur de zéro est inversible.
Ensuite, remarquons que les images inverses des fonctions méromorphes
sont définies pour $j$ d'après le
lemme \ref{lemma-pullback-meromorphic-sections-defined}.
Cela donne un morphisme
$$
\mathcal{K}_X \longrightarrow j_*\mathcal{O}_Y.
$$
Soit $U \in X_\etale$ affine. D'après le lemme \ref{lemma-meromorphic-quasi-coherent},
le membre de gauche prend pour valeur l'anneau total des fractions de $\mathcal{O}_X(U)$.
D'autre part, le membre de droite est égal au produit des
anneaux locaux de $U$ aux points de codimension $0$, c'est-à-dire aux points génériques
de $U$. Ces deux anneaux sont égaux (comme on l'a déjà vu dans la démonstration
du lemme \ref{lemma-meromorphic-quasi-coherent}), d'après les lemmes
d'Algèbre
\ref{algebra-lemma-total-ring-fractions-no-embedded-points} et
\ref{algebra-lemma-weakly-ass-zero-divisors}.
Ainsi, notre morphisme est un isomorphisme.

\medskip\noindent
Enfin, il reste à montrer que $R(X) = \Gamma(X, \mathcal{K}_X)$.
Cela résulte du cas des schémas
(Diviseurs, lemme \ref{divisors-lemma-meromorphic-weakass-finite})
appliqué au lieu schématique $X' \subset X$.
En effet, l'anneau des fonctions rationnelles de $X$ est, par définition,
le même que l'anneau des fonctions rationnelles sur $X'$
car celui-ci est un sous-espace ouvert dense de $X$ (voir ci-dessus).
Bien entendu, $R(X')$ coïncide avec l'anneau des fonctions rationnelles
lorsque $X'$ est considéré comme un schéma.
D'autre part, d'après notre description de $\mathcal{K}_X$ ci-dessus,
et le fait, établi plus haut, que $X^0 \subset |X'|$
est contenu dans tout ouvert dense, on voit que
$\Gamma(X, \mathcal{K}_X) = \Gamma(X', \mathcal{K}_{X'})$.
Enfin, utilisons la compatibilité consignée dans le
lemme \ref{lemma-meromorphic-quasi-coherent-agree}.
\end{proof}

\begin{definition}
\label{definition-regular-meromorphic-section}
Soit $S$ un schéma.
Soit $X$ un espace algébrique sur $S$.
Soit $\mathcal{L}$ un $\mathcal{O}_X$-module inversible.
Une section méromorphe $s$ de $\mathcal{L}$ est dite {\it régulière}
si l'homomorphisme induit $\mathcal{K}_X \to \mathcal{K}_X(\mathcal{L})$
est injectif.
\end{definition}

\noindent
Précisons quand on peut prendre l'image inverse des sections méromorphes (régulières).

\begin{lemma}
\label{lemma-meromorphic-sections-pullback}
Soit $S$ un schéma.
Soit $f : X \to Y$ un morphisme d'espaces algébriques sur $S$.
Supposons que les images inverses des fonctions méromorphes soient définies
pour $f$ (voir la
définition \ref{definition-pullback-meromorphic-sections}).
\begin{enumerate}
\item Soit $\mathcal{F}$ un faisceau de $\mathcal{O}_Y$-modules.
Il existe une application canonique d'image inverse
$f^* : \Gamma(Y, \mathcal{K}_Y(\mathcal{F})) \to
\Gamma(X, \mathcal{K}_X(f^*\mathcal{F}))$
pour les sections méromorphes de $\mathcal{F}$.
\item Soit $\mathcal{L}$ un $\mathcal{O}_X$-module inversible.
Une section méromorphe régulière $s$ de $\mathcal{L}$ a pour image inverse
la section méromorphe régulière $f^*s$ de $f^*\mathcal{L}$.
\end{enumerate}
\end{lemma}

\begin{proof}
Omise.
\end{proof}

\begin{lemma}
\label{lemma-regular-meromorphic-section-exists}
Soit $S$ un schéma. Soit $X$ un espace algébrique sur $S$
vérifiant les conditions (a), (b) et (c) du
lemme \ref{lemma-compute-meromorphic}.
Alors tout $\mathcal{O}_X$-module inversible $\mathcal{L}$ possède
une section méromorphe régulière.
\end{lemma}

\begin{proof}
Avec les notations du lemme \ref{lemma-compute-meromorphic},
la fibre $\mathcal{L}_\eta$ de $\mathcal{L}$ est définie pour tout
$\eta \in X^0$ et c'est un module libre de rang $1$ sur $\mathcal{O}_{X, \eta}$.
Choisissons un générateur $s_\eta \in \mathcal{L}_\eta$ pour tout $\eta \in X^0$.
Il résulte immédiatement de la description de
$\mathcal{K}_X$ et $\mathcal{K}_X(\mathcal{L})$
dans le lemme \ref{lemma-compute-meromorphic}
que $s = \prod s_\eta$ est une section méromorphe régulière de $\mathcal{L}$.
\end{proof}












\section{Proj relatif}
\label{section-relative-proj}

\noindent
Cette section reprend la construction du Proj relatif
dans le cadre des espaces algébriques. Son contenu
correspond à celui de Constructions, section
\ref{constructions-section-relative-proj}
et de Diviseurs, section \ref{divisors-section-relative-proj},
dans le cas des schémas.

\begin{situation}
\label{situation-relative-proj}
Ici, $S$ est un schéma, $X$ est un espace algébrique sur $S$ et
$\mathcal{A}$ est une $\mathcal{O}_X$-algèbre graduée quasi-cohérente.
\end{situation}

\noindent
Dans la situation \ref{situation-relative-proj}, nous allons définir
un foncteur $F : (\Sch/S)_{fppf}^{opp} \to \textit{Ens}$ qui se révélera
être un espace algébrique. Nous suivrons, mutatis mutandis,
la méthode de
Constructions, section \ref{constructions-section-relative-proj}.
Tout d'abord, pour un schéma $T$ sur $S$, nous appelons
{\it quadruplet au-dessus de $T$} un système
$(d, f : T \to X, \mathcal{L}, \psi)$
\begin{enumerate}
\item $d \geq 1$ est un entier,
\item $f : T \to X$ est un morphisme au-dessus de $S$,
\item $\mathcal{L}$ est un $\mathcal{O}_T$-module inversible, et
\item
$\psi : f^*\mathcal{A}^{(d)} \to \bigoplus_{n \geq 0}\mathcal{L}^{\otimes n}$
est un homomorphisme de $\mathcal{O}_T$-algèbres graduées
tel que $f^*\mathcal{A}_d \to \mathcal{L}$ soit surjectif.
\end{enumerate}
Nous disons que les deux quadruplets $(d, f, \mathcal{L}, \psi)$ et
$(d', f', \mathcal{L}', \psi')$ sont {\it équivalents}\footnote{Cette
définition est motivée par
Constructions, lemme \ref{constructions-lemma-equivalent-relative}.
Son avantage est qu'elle définit manifestement
une relation d'équivalence.}
si et seulement si
$f = f'$ et s'il existe un entier positif $m = ad = a'd'$
ainsi qu'un isomorphisme
$\beta : \mathcal{L}^{\otimes a} \to (\mathcal{L}')^{\otimes a'}$
tel que $\beta \circ \psi|_{f^*\mathcal{A}^{(m)}}$ et
$\psi'|_{f^*\mathcal{A}^{(m)}}$ coïncident
comme homomorphismes d'anneaux gradués
$f^*\mathcal{A}^{(m)} \to \bigoplus_{n \geq 0} (\mathcal{L}')^{\otimes mn}$.
Étant donnés un quadruplet $(d, f, \mathcal{L}, \psi)$
et un morphisme $h : T' \to T$, on dispose du quadruplet image inverse
$(d, f \circ h, h^*\mathcal{L}, h^*\psi)$. L'image inverse respecte
la relation d'équivalence. Enfin, pour un schéma {\it quasi-compact} $T$
sur $S$, on pose
$$
F(T) = \text{l'ensemble des classes d'équivalence de quadruplets au-dessus de }T
$$
et, pour un schéma arbitraire $T$ sur $S$, on pose
$$
F(T)
=
\lim_{V \subset T\text{ ouvert quasi-compact}} F(V).
$$
Autrement dit, un élément $\xi$ de $F(T)$ correspond à un système
compatible de choix d'éléments $\xi_V \in F(V)$, où $V$ parcourt les
ouverts quasi-compacts de $T$. Nous avons ainsi défini notre foncteur
\begin{equation}
\label{equation-proj}
F : \Sch^{opp} \longrightarrow \textit{Ens}
\end{equation}
Il existe un morphisme de foncteurs $F \to X$ qui envoie le quadruplet
$(d, f, \mathcal{L}, \psi)$ sur $f$.

\begin{lemma}
\label{lemma-relative-proj}
Dans la situation \ref{situation-relative-proj}, le foncteur $F$ ci-dessus est un
espace algébrique. Pour tout morphisme $g : Z \to X$, où $Z$ est un schéma,
il existe un isomorphisme canonique
$\underline{\text{Proj}}_Z(g^*\mathcal{A}) = Z \times_X F$
compatible avec tout changement de base ultérieur.
\end{lemma}

\begin{proof}
Il suffit de démontrer la seconde assertion, voir
Espaces, lemme \ref{spaces-lemma-representable-over-space}.
Soit $g : Z \to X$ un morphisme, où $Z$ est un schéma.
Soit $F'$ le foncteur des quadruplets associé
à la $\mathcal{O}_Z$-algèbre graduée quasi-cohérente $g^*\mathcal{A}$.
Il existe alors un isomorphisme canonique $F' = Z \times_X F$ qui envoie
un quadruplet $(d, f : T \to Z, \mathcal{L}, \psi)$ pour $F'$
sur $(d, g \circ f, \mathcal{L}, \psi)$ (détails omis, voir la démonstration de
Constructions, lemme \ref{constructions-lemma-proj-base-change}).
D'après les lemmes
\ref{constructions-lemma-equivalent-relative},
\ref{constructions-lemma-relative-proj} et
\ref{constructions-lemma-glueing-gives-functor-proj} de Constructions, ainsi que la
définition \ref{constructions-definition-relative-proj},
$F'$ est représentable par
$\underline{\text{Proj}}_Z(g^*\mathcal{A})$.
\end{proof}

\noindent
Le lemme ci-dessus montre que la définition suivante a un sens.

\begin{definition}
\label{definition-relative-proj}
Soient $S$ un schéma et $X$ un espace algébrique sur $S$.
Soit $\mathcal{A}$ un faisceau quasi-cohérent
d'algèbres graduées sur $\mathcal{O}_X$. Le
{\it spectre homogène relatif de $\mathcal{A}$ sur $X$},
ou le {\it spectre homogène de $\mathcal{A}$ sur $X$}, ou le
{\it Proj relatif de $\mathcal{A}$ sur $X$} est l'espace algébrique
$F$ sur $X$ du lemme \ref{lemma-relative-proj}.
On le note $\pi : \underline{\text{Proj}}_X(\mathcal{A}) \to X$.
\end{definition}

\noindent
En particulier, le morphisme structural du Proj relatif est représentable
par construction. On peut aussi décrire le Proj relatif par recollement. Soit
$\varphi : U \to X$ un morphisme \'etale surjectif, où $U$ est un schéma.
Posons $R = U \times_X U$, avec pour morphismes de projection $s, t : R  \to U$.
D'après le lemme \ref{lemma-relative-proj}, il existe un isomorphisme canonique
$$
\gamma : 
\underline{\text{Proj}}_U(\varphi^*\mathcal{A})
\longrightarrow
\underline{\text{Proj}}_X(\mathcal{A}) \times_X U
$$
au-dessus de $U$. Soit $\alpha : t^*\varphi^*\mathcal{A} \to s^*\varphi^*\mathcal{A}$
l'isomorphisme canonique fourni par la
proposition de Propriétés des espaces
\ref{spaces-properties-proposition-quasi-coherent}.
Alors le diagramme
$$
\xymatrix{
&
\underline{\text{Proj}}_U(\varphi^*\mathcal{A}) \times_{U, s} R
\ar@{=}[r] &
\underline{\text{Proj}}_R(s^*\varphi^*\mathcal{A})
\ar[dd]_{\text{induit par }\alpha} \\
\underline{\text{Proj}}_X(\mathcal{A}) \times_X R
\ar[ru]_{s^*\gamma} \ar[rd]^{t^*\gamma} \\
&
\underline{\text{Proj}}_U(\varphi^*\mathcal{A}) \times_{U, t} R
\ar@{=}[r] &
\underline{\text{Proj}}_R(t^*\varphi^*\mathcal{A})
}
$$
est commutatif (les signes d'égalité proviennent du
lemme \ref{constructions-lemma-relative-proj-base-change} de Constructions).
Ainsi, si l'on note $\mathcal{A}_U$, $\mathcal{A}_R$
les images inverses de $\mathcal{A}$ sur $U$, $R$, alors
$P = \underline{\text{Proj}}_X(\mathcal{A})$ admet pour revêtement \'etale
le schéma $P_U = \underline{\text{Proj}}_U(\mathcal{A}_U)$, et
$P_U \times_P P_U$ est égal à
$P_R = \underline{\text{Proj}}_R(\mathcal{A}_R)$.
Ces remarques permettent de raisonner de la manière habituelle par
localisation \'etale afin de transporter les résultats sur le Proj relatif du cas
des schémas à celui des espaces algébriques.

\begin{lemma}
\label{lemma-twists-of-structure-sheaf}
Dans la situation \ref{situation-relative-proj}, le Proj relatif est muni
d'un faisceau quasi-cohérent d'algèbres graduées par $\mathbf{Z}$
$\bigoplus_{n \in \mathbf{Z}}
\mathcal{O}_{\underline{\text{Proj}}_X(\mathcal{A})}(n)$
et d'un homomorphisme canonique d'algèbres graduées
$$
\psi :
\pi^*\mathcal{A}
\longrightarrow
\bigoplus\nolimits_{n \geq 0}
\mathcal{O}_{\underline{\text{Proj}}_X(\mathcal{A})}(n)
$$
dont le changement de base à tout schéma sur $X$ coïncide avec
celui du lemme \ref{constructions-lemma-glue-relative-proj-twists} de Constructions.
\end{lemma}

\begin{proof}
Comme dans la discussion qui suit la définition \ref{definition-relative-proj},
choisissons un schéma $U$ et un morphisme \'etale surjectif
$U \to X$, puis posons $R = U \times_X U$, avec projections $s, t : R \to U$,
$\mathcal{A}_U = \mathcal{A}|_U$, $\mathcal{A}_R = \mathcal{A}|_R$,
ainsi que $\pi : P = \underline{\text{Proj}}_X(\mathcal{A}) \to X$,
$\pi_U : P_U = \underline{\text{Proj}}_U(\mathcal{A}_U)$ et
$\pi_R : P_R = \underline{\text{Proj}}_U(\mathcal{A}_R)$.
D'après le
lemme \ref{constructions-lemma-glue-relative-proj-twists} de Constructions,
on dispose d'un faisceau quasi-cohérent d'algèbres graduées par $\mathbf{Z}$
sur $\mathcal{O}_{P_U}$,
$\bigoplus_{n \in \mathbf{Z}} \mathcal{O}_{P_U}(n)$
ainsi que d'un homomorphisme canonique
$\psi_U : \pi_U^*\mathcal{A}_U \to \bigoplus_{n \geq 0} \mathcal{O}_{P_U}(n)$
et de même pour $P_R$. D'après le
lemme \ref{constructions-lemma-relative-proj-base-change} de Constructions,
les images inverses de $\mathcal{O}_{P_U}(n)$ et de $\psi_U$ par l'une ou l'autre projection
$P_R \to P_U$ sont $\mathcal{O}_{P_R}(n)$ et $\psi_R$.
D'après la proposition de Propriétés des espaces
\ref{spaces-properties-proposition-quasi-coherent}
on obtient $\mathcal{O}_{P}(n)$ et $\psi$.
Nous omettons de vérifier la compatibilité avec l'image inverse sur
un schéma arbitraire au-dessus de $X$.
\end{proof}

\noindent
Le Proj relatif étant construit, passons à quelques propriétés
élémentaires.

\begin{lemma}
\label{lemma-relative-proj-base-change}
Soit $S$ un schéma. Soit $g : X' \to X$ un morphisme d'espaces algébriques
sur $S$, et soit $\mathcal{A}$ un faisceau quasi-cohérent
d'algèbres graduées sur $\mathcal{O}_X$. Il existe alors un isomorphisme canonique
$$
r :
\underline{\text{Proj}}_{X'}(g^*\mathcal{A})
\longrightarrow
X' \times_X \underline{\text{Proj}}_X(\mathcal{A})
$$
ainsi qu'un isomorphisme correspondant
$$
\theta :
r^*\text{pr}_2^*\left(\bigoplus\nolimits_{d \in \mathbf{Z}}
\mathcal{O}_{\underline{\text{Proj}}_X(\mathcal{A})}(d)\right)
\longrightarrow
\bigoplus\nolimits_{d \in \mathbf{Z}}
\mathcal{O}_{\underline{\text{Proj}}_{X'}(g^*\mathcal{A})}(d)
$$
d'algèbres graduées par $\mathbf{Z}$
sur $\mathcal{O}_{\underline{\text{Proj}}_{X'}(g^*\mathcal{A})}$.
\end{lemma}

\begin{proof}
Soit $F$ le foncteur (\ref{equation-proj}) et soit $F'$ le
foncteur correspondant défini à l'aide de $g^*\mathcal{A}$ sur $X'$.
Nous affirmons qu'il existe un isomorphisme canonique $r : F' \to X' \times_X F$
de foncteurs (et, bien sûr, $r$ est l'isomorphisme du lemme).
Il suffit de construire la bijection
$r : F'(T) \to X'(T) \times_{X(T)} F(T)$ pour les schémas quasi-compacts $T$
sur $S$. Tout d'abord, si $\xi = (d', f', \mathcal{L}', \psi')$ est un
quadruplet au-dessus de $T$ pour $F'$, on peut poser
$r(\xi) = (f', (d', g \circ f', \mathcal{L}', \psi'))$. Cette expression a un sens
car $(g \circ f')^*\mathcal{A}^{(d)} = (f')^*(g^*\mathcal{A})^{(d)}$.
L'application inverse envoie la paire $(f', (d, f, \mathcal{L}, \psi))$
sur le quadruplet $(d, f', \mathcal{L}, \psi)$. Nous omettons la démonstration
de la dernière assertion (indication : se ramener au cas des schémas par localisation
\'etale et appliquer le lemme de Constructions
\ref{constructions-lemma-relative-proj-base-change}).
\end{proof}

\begin{lemma}
\label{lemma-relative-proj-separated}
Dans la situation \ref{situation-relative-proj}, le morphisme
$\pi : \underline{\text{Proj}}_X(\mathcal{A}) \to X$
est séparé.
\end{lemma}

\begin{proof}
D'après le lemme \ref{spaces-morphisms-lemma-separated-local} de Morphismes d'espaces
et la construction du Proj relatif, cela résulte du
cas des schémas, traité dans
Constructions, lemme \ref{constructions-lemma-relative-proj-separated}.
\end{proof}

\begin{lemma}
\label{lemma-relative-proj-quasi-compact}
Dans la situation \ref{situation-relative-proj}, si l'une des conditions suivantes est satisfaite :
\begin{enumerate}
\item $\mathcal{A}$ est de type fini en tant que faisceau
d'algèbres sur $\mathcal{A}_0$,
\item $\mathcal{A}$ est engendrée par $\mathcal{A}_1$ en tant
qu'algèbre sur $\mathcal{A}_0$, et $\mathcal{A}_1$ est un
$\mathcal{A}_0$-module de type fini,
\item il existe un sous-$\mathcal{A}_0$-module quasi-cohérent de type fini
$\mathcal{F} \subset \mathcal{A}_{+}$ tel que
$\mathcal{A}_{+}/\mathcal{F}\mathcal{A}$ soit un faisceau d'idéaux localement nilpotent
de $\mathcal{A}/\mathcal{F}\mathcal{A}$,
\end{enumerate}
alors $\pi : \underline{\text{Proj}}_X(\mathcal{A}) \to X$ est quasi-compact.
\end{lemma}

\begin{proof}
D'après le lemme \ref{spaces-morphisms-lemma-quasi-compact-local} de Morphismes d'espaces
et la construction du Proj relatif, cela résulte du
cas des schémas, traité dans
Diviseurs, lemme \ref{divisors-lemma-relative-proj-quasi-compact}.
\end{proof}

\begin{lemma}
\label{lemma-relative-proj-finite-type}
Dans la situation \ref{situation-relative-proj},
si $\mathcal{A}$ est de type fini en tant que faisceau
d'algèbres sur $\mathcal{O}_X$, alors
$\pi : \underline{\text{Proj}}_X(\mathcal{A}) \to X$ est de type fini.
\end{lemma}

\begin{proof}
D'après le lemme \ref{spaces-morphisms-lemma-finite-type-local} de Morphismes d'espaces
et la construction du Proj relatif, cela résulte du
cas des schémas, traité dans
Diviseurs, lemme \ref{divisors-lemma-relative-proj-finite-type}.
\end{proof}

\begin{lemma}
\label{lemma-relative-proj-universally-closed}
Dans la situation \ref{situation-relative-proj}, si
$\mathcal{O}_X \to \mathcal{A}_0$
est un homomorphisme d'algèbres entier\footnote{Autrement dit, la clôture intégrale
de $\mathcal{O}_X$ dans $\mathcal{A}_0$, voir
Morphismes d'espaces, définition
\ref{spaces-morphisms-definition-integral-closure}, est égale à
$\mathcal{A}_0$.} et si $\mathcal{A}$ est de type fini en tant
qu'algèbre sur $\mathcal{A}_0$, alors
$\pi : \underline{\text{Proj}}_X(\mathcal{A}) \to X$ est universellement fermé.
\end{lemma}

\begin{proof}
D'après le lemme de Morphismes d'espaces
\ref{spaces-morphisms-lemma-universally-closed-local}
et la construction du Proj relatif, cela résulte du
cas des schémas, traité dans
Diviseurs, lemme \ref{divisors-lemma-relative-proj-universally-closed}.
\end{proof}

\begin{lemma}
\label{lemma-relative-proj-proper}
Dans la situation \ref{situation-relative-proj},
les conditions suivantes sont équivalentes :
\begin{enumerate}
\item $\mathcal{A}_0$ est un $\mathcal{O}_X$-module de type fini
et $\mathcal{A}$ est de type fini en tant qu'algèbre sur $\mathcal{A}_0$,
\item $\mathcal{A}_0$ est un $\mathcal{O}_X$-module de type fini
et $\mathcal{A}$ est de type fini en tant qu'algèbre sur $\mathcal{O}_X$.
\end{enumerate}
Si ces conditions sont satisfaites, alors
$\pi : \underline{\text{Proj}}_X(\mathcal{A}) \to X$
est propre.
\end{lemma}

\begin{proof}
D'après le lemme de Morphismes d'espaces
\ref{spaces-morphisms-lemma-proper-local}
et la construction du Proj relatif, cela résulte du
cas des schémas, traité dans
Diviseurs, lemme \ref{divisors-lemma-relative-proj-universally-closed}.
\end{proof}

\begin{lemma}
\label{lemma-relative-proj-generated-in-degree-1}
Soient $S$ un schéma et $X$ un espace algébrique sur $S$.
Soit $\mathcal{A}$ un faisceau quasi-cohérent d'algèbres graduées sur $\mathcal{O}_X$
engendré en tant qu'algèbre sur $\mathcal{A}_0$ par $\mathcal{A}_1$.
Avec $P = \underline{\text{Proj}}_X(\mathcal{A})$, on a
\begin{enumerate}
\item $P$ représente le foncteur $F_1$ qui associe à
$T$ sur $S$ l'ensemble des classes d'isomorphie de
triplets $(f, \mathcal{L}, \psi)$, où $f : T \to X$ est un morphisme
au-dessus de $S$, $\mathcal{L}$ est un $\mathcal{O}_T$-module inversible, et
$\psi : f^*\mathcal{A} \to \bigoplus_{n \geq 0} \mathcal{L}^{\otimes n}$
est un homomorphisme de $\mathcal{O}_T$-algèbres graduées induisant une surjection
$f^*\mathcal{A}_1 \to \mathcal{L}$,
\item l'homomorphisme canonique $\pi^*\mathcal{A}_1 \to \mathcal{O}_P(1)$ est
surjectif, et
\item chaque $\mathcal{O}_P(n)$ est inversible,
et les homomorphismes de multiplication induisent des isomorphismes
$\mathcal{O}_P(n) \otimes_{\mathcal{O}_P} \mathcal{O}_P(m) =
\mathcal{O}_P(n + m)$.
\end{enumerate}
\end{lemma}

\begin{proof}
Omise.
Voir Constructions, lemme \ref{constructions-lemma-apply-relative},
pour le cas des schémas.
\end{proof}







\section{Fonctorialité du Proj relatif}
\label{section-functoriality-relative-proj}

\noindent
Cette section est l'analogue de
Constructions, section \ref{constructions-section-functoriality-relative-proj}.

\begin{lemma}
\label{lemma-morphism-relative-proj}
Soient $S$ un schéma et $X$ un espace algébrique sur $S$.
Soit $\psi : \mathcal{A} \to \mathcal{B}$ un homomorphisme
d'algèbres graduées quasi-cohérentes sur $\mathcal{O}_X$. Posons
$P = \underline{\text{Proj}}_X(\mathcal{A}) \to X$ et
$Q = \underline{\text{Proj}}_X(\mathcal{B}) \to X$.
Il existe un sous-espace ouvert canonique
$U(\psi) \subset Q$ et un morphisme canonique
d'espaces algébriques
$$
r_\psi :
U(\psi)
\longrightarrow
P
$$
au-dessus de $X$ et un homomorphisme d'algèbres graduées par $\mathbf{Z}$ sur $\mathcal{O}_{U(\psi)}$
$$
\theta = \theta_\psi :
r_\psi^*\left(
\bigoplus\nolimits_{d \in \mathbf{Z}} \mathcal{O}_P(d)
\right)
\longrightarrow
\bigoplus\nolimits_{d \in \mathbf{Z}} \mathcal{O}_{U(\psi)}(d).
$$
Le triplet $(U(\psi), r_\psi, \theta)$ est caractérisé par la propriété
que, pour tout schéma $W$ \'etale sur $X$, le triplet
$$
(U(\psi) \times_X W,\quad
r_\psi|_{U(\psi) \times_X W} : U(\psi) \times_X W \to  P \times_X W,\quad
\theta|_{U(\psi) \times_X W})
$$
est égal au triplet associé à $\psi : \mathcal{A}|_W \to \mathcal{B}|_W$
par Constructions, lemme \ref{constructions-lemma-morphism-relative-proj}.
\end{lemma}

\begin{proof}
Ce lemme résulte de la localisation \'etale et du cas des schémas, voir la
discussion qui suit la
définition \ref{definition-relative-proj}. Détails omis.
\end{proof}

\begin{lemma}
\label{lemma-morphism-relative-proj-transitive}
Soient $S$ un schéma et $X$ un espace algébrique sur $S$.
Soient $\mathcal{A}$, $\mathcal{B}$ et $\mathcal{C}$ des
algèbres graduées quasi-cohérentes sur $\mathcal{O}_X$.
Posons $P = \underline{\text{Proj}}_X(\mathcal{A})$,
$Q = \underline{\text{Proj}}_X(\mathcal{B})$ et
$R = \underline{\text{Proj}}_X(\mathcal{C})$.
Soient $\varphi : \mathcal{A} \to \mathcal{B}$,
$\psi : \mathcal{B} \to \mathcal{C}$ des homomorphismes d'algèbres graduées sur $\mathcal{O}_X$.
Alors
$$
U(\psi \circ \varphi) = r_\varphi^{-1}(U(\psi))
\quad
\text{et}
\quad
r_{\psi \circ \varphi}
=
r_\varphi \circ r_\psi|_{U(\psi \circ \varphi)}.
$$
De plus,
$$
\theta_\psi \circ r_\psi^*\theta_\varphi
=
\theta_{\psi \circ \varphi}
$$
avec les notations évidentes.
\end{lemma}

\begin{proof}
Omise.
\end{proof}

\begin{lemma}
\label{lemma-surjective-graded-rings-map-relative-proj}
Sous les hypothèses et avec les notations du lemme \ref{lemma-morphism-relative-proj}
ci-dessus, supposons que $\mathcal{A}_d \to \mathcal{B}_d$ soit surjectif pour
$d \gg 0$. Alors
\begin{enumerate}
\item $U(\psi) = Q$,
\item $r_\psi : Q \to R$ est une immersion fermée, et
\item les homomorphismes $\theta : r_\psi^*\mathcal{O}_P(n) \to \mathcal{O}_Q(n)$
sont surjectifs mais ne sont pas en général des isomorphismes (même si
$\mathcal{A} \to \mathcal{B}$ est surjectif).
\end{enumerate}
\end{lemma}

\begin{proof}
Cela résulte du cas des schémas
(Constructions, lemme
\ref{constructions-lemma-surjective-graded-rings-map-relative-proj})
par localisation \'etale.
\end{proof}

\begin{lemma}
\label{lemma-eventual-iso-graded-rings-map-relative-proj}
Sous les hypothèses et avec les notations du lemme \ref{lemma-morphism-relative-proj}
ci-dessus, supposons que $\mathcal{A}_d \to \mathcal{B}_d$ soit un isomorphisme pour tout
$d \gg 0$. Alors
\begin{enumerate}
\item $U(\psi) = Q$,
\item $r_\psi : Q \to P$ est un isomorphisme, et
\item les homomorphismes $\theta : r_\psi^*\mathcal{O}_P(n) \to \mathcal{O}_Q(n)$
sont des isomorphismes.
\end{enumerate}
\end{lemma}

\begin{proof}
Cela résulte du cas des schémas
(Constructions, lemme
\ref{constructions-lemma-eventual-iso-graded-rings-map-relative-proj})
par localisation \'etale.
\end{proof}

\begin{lemma}
\label{lemma-surjective-generated-degree-1-map-relative-proj}
Sous les hypothèses et avec les notations du lemme \ref{lemma-morphism-relative-proj}
ci-dessus, supposons que $\mathcal{A}_d \to \mathcal{B}_d$ soit surjectif pour $d \gg 0$
et que $\mathcal{A}$ soit engendrée par $\mathcal{A}_1$ sur $\mathcal{A}_0$.
Alors
\begin{enumerate}
\item $U(\psi) = Q$,
\item $r_\psi : Q \to P$ est une immersion fermée, et
\item les homomorphismes $\theta : r_\psi^*\mathcal{O}_P(n) \to \mathcal{O}_Q(n)$
sont des isomorphismes.
\end{enumerate}
\end{lemma}

\begin{proof}
Cela résulte du cas des schémas
(Constructions, lemme
\ref{constructions-lemma-surjective-generated-degree-1-map-relative-proj})
par localisation \'etale.
\end{proof}










\section{Faisceaux inversibles et morphismes vers le Proj relatif}
\label{section-invertible-relative-proj}

\noindent
Il semble que le lemme suivant puisse être utile quelque part.
La situation est la suivante :
\begin{enumerate}
\item Soient $S$ un schéma et $Y$ un espace algébrique sur $S$.
\item Soit $\mathcal{A}$ une algèbre graduée quasi-cohérente sur $\mathcal{O}_Y$.
\item Notons $\pi : \underline{\text{Proj}}_Y(\mathcal{A}) \to Y$ le Proj
relatif de $\mathcal{A}$ sur $Y$.
\item Soit $f : X \to Y$ un morphisme d'espaces algébriques sur $S$.
\item Soit $\mathcal{L}$ un $\mathcal{O}_X$-module inversible.
\item Soit
$\psi : f^*\mathcal{A} \to \bigoplus_{d \geq 0} \mathcal{L}^{\otimes d}$
un homomorphisme d'algèbres graduées sur $\mathcal{O}_X$.
\end{enumerate}
Ces données étant fixées, soit $U(\psi) \subset X$ le sous-espace ouvert tel que
$$
|U(\psi)| = \bigcup\nolimits_{d \geq 1}
\{\text{lieu où }f^*\mathcal{A}_d \to \mathcal{L}^{\otimes d}
\text{ est surjectif}\}
$$
La formation de $U(\psi) \subset X$ commute au
changement de base par tout morphisme $X' \to X$.

\begin{lemma}
\label{lemma-invertible-map-into-relative-proj}
Sous les hypothèses et avec les notations ci-dessus, le morphisme
$\psi$ induit un morphisme canonique d'espaces algébriques sur $Y$
$$
r_{\mathcal{L}, \psi} :
U(\psi) \longrightarrow \underline{\text{Proj}}_Y(\mathcal{A})
$$
ainsi qu'un homomorphisme d'algèbres graduées sur $\mathcal{O}_{U(\psi)}$
$$
\theta :
r_{\mathcal{L}, \psi}^*\left(
\bigoplus\nolimits_{d \geq 0}
\mathcal{O}_{\underline{\text{Proj}}_Y(\mathcal{A})}(d)
\right)
\longrightarrow
\bigoplus\nolimits_{d \geq 0} \mathcal{L}^{\otimes d}|_{U(\psi)}
$$
caractérisé par les propriétés suivantes :
\begin{enumerate}
\item Pour $V \to Y$ \'etale et $d \geq 0$, le diagramme
$$
\xymatrix{
\mathcal{A}_d(V) \ar[d]_{\psi} \ar[r]_{\psi} &
\Gamma(V \times_Y X, \mathcal{L}^{\otimes d}) \ar[d]^{restriction} \\
\Gamma(V \times_Y \underline{\text{Proj}}_Y(\mathcal{A}),
\mathcal{O}_{\underline{\text{Proj}}_Y(\mathcal{A})}(d)) \ar[r]^-\theta &
\Gamma(V \times_Y U(\psi), \mathcal{L}^{\otimes d})
}
$$
est commutatif.
\item Pour tout $d \geq 1$ et tout morphisme $W \to X$, où $W$ est un schéma,
tel que $\psi|_W : f^*\mathcal{A}_d|_W \to \mathcal{L}^{\otimes d}|_W$
soit surjectif, on a (a) $W \to X$ se factorise par $U(\psi)$ et
(b) le composé de $W \to U(\psi)$ avec $r_{\mathcal{L}, \psi}$
coïncide avec le morphisme $W \to \underline{\text{Proj}}_Y(\mathcal{A})$
dont l'existence résulte de la construction de $\underline{\text{Proj}}_Y(\mathcal{A})$,
voir la définition \ref{definition-relative-proj}.
\item Considérons un diagramme commutatif
$$
\xymatrix{
X' \ar[r]_{g'} \ar[d]_{f'} & X \ar[d]^f \\
Y' \ar[r]^g & Y
}
$$
où $X'$ et $Y'$ sont des schémas, posons $\mathcal{A}' = g^*\mathcal{A}$
et $\mathcal{L}' = (g')^*\mathcal{L}$ et notons
$\psi' : (f')^*\mathcal{A} \to \bigoplus_{d \geq 0} (\mathcal{L}')^{\otimes d}$
l'image inverse de $\psi$. Notons $U(\psi')$, $r_{\psi', \mathcal{L}'}$
et $\theta'$ le sous-espace ouvert, le morphisme et l'homomorphisme construits
dans Constructions, lemme \ref{lemma-invertible-map-into-relative-proj}.
Alors $U(\psi') = (g')^{-1}(U(\psi))$
et $r_{\psi', \mathcal{L}'}$ coïncide avec le changement de base
de $r_{\psi, \mathcal{L}}$ via l'isomorphisme
$\underline{\text{Proj}}_{Y'}(\mathcal{A}') =
Y' \times_Y \underline{\text{Proj}}_Y(\mathcal{A})$
du lemme \ref{lemma-relative-proj-base-change}.
De plus, $\theta'$ est l'image inverse de $\theta$.
\end{enumerate}
\end{lemma}

\begin{proof}
Omise. Indications :
Observons d'abord que, pour un schéma quasi-compact
$W$ sur $X$, les conditions suivantes sont équivalentes :
\begin{enumerate}
\item $W \to X$ se factorise par $U(\psi)$, et
\item il existe un $d$ tel que
$\psi|_W : f^*\mathcal{A}_d|_W \to \mathcal{L}^{\otimes d}|_W$
soit surjectif.
\end{enumerate}
On obtient ainsi une description de $U(\psi)$ comme sous-foncteur de $X$
sur notre catégorie de base $(\Sch/S)_{fppf}$. Pour de tels $W$ et $d$,
considérons le quadruplet
$(d, W \to Y, \mathcal{L}|_W, \psi^{(d)}|_W)$.
D'après la définition de $\underline{\text{Proj}}_Y(\mathcal{A})$,
on obtient un morphisme $W \to \underline{\text{Proj}}_Y(\mathcal{A})$.
Notre notion d'équivalence des quadruplets montre que
ce morphisme est indépendant du choix de $d$.
Cela définit manifestement une transformation de foncteurs
$r_{\psi, \mathcal{L}} : U(\psi) \to \underline{\text{Proj}}_Y(\mathcal{A})$,
c'est-à-dire un morphisme d'espaces algébriques.
Par construction, ce morphisme vérifie (2).
Comme le morphisme construit dans
Constructions, lemme \ref{constructions-lemma-invertible-map-into-relative-proj}
vérifie la même propriété, (3) est vraie.

\medskip\noindent
Pour construire $\theta$ et vérifier la compatibilité (1) du
lemme, travaillons localement pour la topologie \'etale sur $Y$ et $X$, en raisonnant comme
dans la discussion qui suit la
définition \ref{definition-relative-proj}.
\end{proof}








\section{Faisceaux relativement amples}
\label{section-relatively-ample}

\noindent
Cette section est l'analogue de
Morphismes, section \ref{morphisms-section-relatively-ample}
pour les espaces algébriques.
Notre définition d'un faisceau inversible relativement ample est
la suivante.

\begin{definition}
\label{definition-relatively-ample}
Soit $S$ un schéma.
Soit $f : X \to Y$ un morphisme d'espaces algébriques sur $S$.
Soit $\mathcal{L}$ un $\mathcal{O}_X$-module inversible.
On dit que $\mathcal{L}$ est {\it relativement ample}, ou {\it ample relativement à $f$},
ou {\it ample sur $X/Y$}, ou {\it $f$-ample}, si $f : X \to Y$
est représentable et si, pour tout morphisme $Z \to Y$
où $Z$ est un schéma, l'image inverse $\mathcal{L}_Z$ de $\mathcal{L}$
sur $X_Z = Z \times_Y X$ est ample sur $X_Z/Z$ au sens de
Morphismes, définition \ref{morphisms-definition-relatively-ample}.
\end{definition}

\noindent
Nous ramènerons presque toujours les questions sur les faisceaux inversibles
relativement amples au cas des schémas. Ainsi, dans cette section, nous n'avons
essentiellement que des vérifications de cohérence.

\begin{lemma}
\label{lemma-relatively-ample-sanity-check}
Soit $S$ un schéma.
Soit $f : X \to Y$ un morphisme d'espaces algébriques sur $S$.
Soit $\mathcal{L}$ un $\mathcal{O}_X$-module inversible.
Supposons que $Y$ soit un schéma. Les conditions suivantes sont équivalentes :
\begin{enumerate}
\item $\mathcal{L}$ est ample sur $X/Y$ au sens de la
définition \ref{definition-relatively-ample}, et
\item $X$ est un schéma et $\mathcal{L}$ est ample sur $X/Y$
au sens de
Morphismes, définition \ref{morphisms-definition-relatively-ample}.
\end{enumerate}
\end{lemma}

\begin{proof}
Cela résulte des définitions et de
Morphismes, lemme \ref{morphisms-lemma-ample-base-change}
(qui affirme que, pour les schémas, la propriété d'être relativement ample
est préservée par changement de base).
\end{proof}

\begin{lemma}
\label{lemma-ample-base-change}
Soit $S$ un schéma.
Soit $f : X \to Y$ un morphisme d'espaces algébriques sur $S$.
Soit $\mathcal{L}$ un $\mathcal{O}_X$-module inversible.
Soit $Y' \to Y$ un morphisme d'espaces algébriques sur $S$.
Soit $f' : X' \to Y'$ le changement de base de $f$ et notons
$\mathcal{L}'$ l'image inverse de $\mathcal{L}$ sur $X'$.
Si $\mathcal{L}$ est $f$-ample, alors $\mathcal{L}'$ est $f'$-ample.
\end{lemma}

\begin{proof}
Cela résulte immédiatement de la définition !
(Indication : transitivité du changement de base.)
\end{proof}

\begin{lemma}
\label{lemma-relatively-ample-properties}
Soit $S$ un schéma.
Soit $f : X \to Y$ un morphisme d'espaces algébriques sur $S$.
S'il existe un faisceau inversible $f$-ample, alors
$f$ est représentable, quasi-compact et séparé.
\end{lemma}

\begin{proof}
C'est clair d'après les définitions et
Morphismes, lemme \ref{morphisms-lemma-relatively-ample-separated}.
(En cas de doute, voir le principe exposé dans
Espaces algébriques, lemme
\ref{spaces-lemma-representable-transformations-property-implication}.)
\end{proof}

\begin{lemma}
\label{lemma-descend-relatively-ample}
Soit $V \to U$ un morphisme \'etale surjectif de schémas affines.
Soit $X$ un espace algébrique sur $U$.
Soit $\mathcal{L}$ un $\mathcal{O}_X$-module inversible.
Soit $Y = V \times_U X$ et soit $\mathcal{N}$
l'image inverse de $\mathcal{L}$ sur $Y$.
Les conditions suivantes sont équivalentes :
\begin{enumerate}
\item $\mathcal{L}$ est ample sur $X/U$, et
\item $\mathcal{N}$ est ample sur $Y/V$.
\end{enumerate}
\end{lemma}

\begin{proof}
L'implication (1) $\Rightarrow$ (2) résulte du
lemme \ref{lemma-ample-base-change}.
Supposons (2). Cela implique que $Y \to V$ est
quasi-compact et séparé (lemme \ref{lemma-relatively-ample-properties})
et que $Y$ est un schéma. Il s'ensuit que le morphisme $f : X \to U$ est
quasi-compact et séparé
(Morphismes d'espaces, lemmes
\ref{spaces-morphisms-lemma-quasi-compact-local} et
\ref{spaces-morphisms-lemma-separated-local}).
Posons $\mathcal{A} = \bigoplus_{d \geq 0} f_*\mathcal{L}^{\otimes d}$.
C'est un faisceau quasi-cohérent d'algèbres graduées sur $\mathcal{O}_U$
(Morphismes d'espaces, lemme \ref{spaces-morphisms-lemma-pushforward}).
Par adjonction, nous avons un morphisme
$\psi : f^*\mathcal{A} \to \bigoplus_{d \geq 0} \mathcal{L}^{\otimes d}$.
En appliquant le lemme \ref{lemma-invertible-map-into-relative-proj},
nous obtenons un sous-espace ouvert $U(\psi) \subset X$ et un morphisme
$$
r_{\mathcal{L}, \psi} : U(\psi) \to \underline{\text{Proj}}_U(\mathcal{A})
$$
Puisque $h : V \to U$ est \'etale, nous avons
$\mathcal{A}|_V = (Y \to V)_*(\bigoplus_{d \geq 0} \mathcal{N}^{\otimes d})$,
voir Propriétés des espaces, lemme
\ref{spaces-properties-lemma-pushforward-etale-base-change-modules}.
Il s'ensuit que l'image inverse $\psi'$ de $\psi$ sur
$Y$ est le morphisme d'adjonction associé à la situation $(Y \to V, \mathcal{N})$
considérée dans Morphismes, lemme \ref{morphisms-lemma-characterize-relatively-ample}, partie (5).
Comme $\mathcal{N}$ est ample sur $Y/V$, le lemme que nous venons de
citer donne $U(\psi') = Y$ et montre que $r_{\mathcal{N}, \psi'}$
est une immersion ouverte.
Puisque le lemme \ref{lemma-invertible-map-into-relative-proj}
nous dit que la formation de $r_{\mathcal{L}, \psi}$
commute au changement de base, nous en concluons que
$U(\psi) = X$ et que nous avons un diagramme commutatif
$$
\xymatrix{
Y \ar[r]_-{r'} \ar[d] &
\underline{\text{Proj}}_V(\mathcal{A}|_V) \ar[d] \ar[r] &
V \ar[d] \\
X \ar[r]^-r &
\underline{\text{Proj}}_U(\mathcal{A}) \ar[r] &
U
}
$$
dont les carrés sont cartésiens. Nous en concluons que $r$ est une
immersion ouverte par
Morphismes d'espaces, lemme \ref{spaces-morphisms-lemma-closed-immersion-local}.
Ainsi, $X$ est un schéma. Nous pouvons alors appliquer
Morphismes, lemme \ref{morphisms-lemma-characterize-relatively-ample}, partie (5),
pour conclure que $\mathcal{L}$ est ample sur $X/U$.
\end{proof}

\begin{lemma}
\label{lemma-relatively-ample-local}
Soit $S$ un schéma. Soit $f : X \to Y$ un morphisme d'espaces algébriques
sur $S$. Soit $\mathcal{L}$ un $\mathcal{O}_X$-module inversible.
Les conditions suivantes sont équivalentes :
\begin{enumerate}
\item $\mathcal{L}$ est ample sur $X/Y$,
\item pour tout schéma $Z$ et tout morphisme $Z \to Y$,
l'espace algébrique $X_Z = Z \times_Y X$ est un schéma
et l'image inverse $\mathcal{L}_Z$ est ample sur $X_Z/Z$,
\item pour tout schéma affine $Z$ et tout morphisme $Z \to Y$,
l'espace algébrique $X_Z = Z \times_Y X$ est un schéma
et l'image inverse $\mathcal{L}_Z$ est ample sur $X_Z/Z$,
\item il existe un schéma $V$ et un morphisme \'etale surjectif
$V \to Y$ tels que l'espace algébrique $X_V = V \times_Y X$ soit un schéma
et que l'image inverse $\mathcal{L}_V$ soit ample sur $X_V/V$.
\end{enumerate}
\end{lemma}

\begin{proof}
Les assertions (1) et (2) sont équivalentes par définition.
L'implication (2) $\Rightarrow$ (3) est immédiate.
Si (3) est satisfaite et si $Z \to Y$ est comme dans (2), nous voyons
que la représentabilité de $X_Z \to Z$ se vérifie localement sur les ouverts affines de $Z$.
Ainsi, $X_Z$ est un schéma, par exemple par
Propriétés des espaces, lemme \ref{spaces-properties-lemma-subscheme}.
Il en résulte alors que $\mathcal{L}_Z$ est ample sur $X_Z/Z$, car
cela est vrai localement sur $Z$ et nous pouvons utiliser
Morphismes, lemme \ref{morphisms-lemma-characterize-relatively-ample}.
Ainsi, (1), (2) et (3) sont équivalentes. Ces conditions
impliquent clairement (4).

\medskip\noindent
Supposons (4). Soit $Z \to Y$ un morphisme avec $Z$ affine.
Alors $U = V \times_Y Z \to Z$ est un morphisme \'etale surjectif
tel que l'image inverse de $\mathcal{L}_Z$ par $X_U \to X_Z$
soit relativement ample sur $X_U/U$.
Nous pouvons remplacer $U$ par une somme disjointe finie d'ouverts affines dont les images recouvrent encore la base.
Il s'ensuit que $\mathcal{L}_Z$ est ample sur $X_Z/Z$ par le
lemme \ref{lemma-descend-relatively-ample}.
Ainsi (4) $\Rightarrow$ (3), ce qui achève la démonstration.
\end{proof}

\begin{lemma}
\label{lemma-quasi-affine-O-ample}
Soit $S$ un schéma. Soit $f : X \to Y$ un morphisme d'espaces algébriques
sur $S$. Alors $f$ est quasi-affine si et seulement si $\mathcal{O}_X$
est ample relativement à $f$.
\end{lemma}

\begin{proof}
Cela résulte du cas des schémas, voir
Morphismes, lemme \ref{morphisms-lemma-quasi-affine-O-ample}.
\end{proof}



\section{Amplitude relative et cohomologie}
\label{section-ample-and-proper}

\noindent
Cette section contient quelques résultats apparentés à ceux
de Cohomologie des schémas, sections
\ref{coherent-section-applications-formal-functions} et
\ref{coherent-section-ample-cohomology}.

\medskip\noindent
Le lemme suivant n'est qu'un exemple de ce que l'on peut faire.

\begin{lemma}
\label{lemma-vanshing-gives-ample}
Soit $R$ un anneau noethérien. Soit $X$ un espace algébrique sur $R$
tel que le morphisme structural $f : X \to \Spec(R)$ soit propre.
Soit $\mathcal{L}$ un $\mathcal{O}_X$-module inversible.
Les conditions suivantes sont équivalentes :
\begin{enumerate}
\item $\mathcal{L}$ est ample sur $X/R$
(définition \ref{definition-relatively-ample}),
\item pour tout $\mathcal{O}_X$-module cohérent $\mathcal{F}$,
il existe un $n_0 \geq 0$ tel que
$H^p(X, \mathcal{F} \otimes \mathcal{L}^{\otimes n}) = 0$
pour tous $n \geq n_0$ et $p > 0$.
\end{enumerate}
\end{lemma}

\begin{proof}
L'implication (1) $\Rightarrow$ (2) résulte de
Cohomologie des schémas, lemme \ref{coherent-lemma-coherent-proper-ample},
car l'hypothèse (1) implique que $X$ est un schéma.
L'implication (2) $\Rightarrow$ (1) est le
lemme de Cohomologie des espaces
\ref{spaces-cohomology-lemma-Noetherian-h1-zero-invertible}.
\end{proof}

\begin{lemma}
\label{lemma-ample-on-fibre}
Soit $Y$ un schéma noethérien. Soit $X$ un espace algébrique sur $Y$
tel que le morphisme structural $f : X \to Y$ soit propre.
Soit $\mathcal{L}$ un $\mathcal{O}_X$-module inversible.
Soit $\mathcal{F}$ un $\mathcal{O}_X$-module cohérent.
Soit $y \in Y$ un point tel que $X_y$ soit un schéma et que
$\mathcal{L}_y$ soit ample sur $X_y$.
Il existe alors un $d_0$ tel que, pour tout $d \geq d_0$, on ait
$$
R^pf_*(\mathcal{F} \otimes_{\mathcal{O}_X} \mathcal{L}^{\otimes d})_y = 0
\text{ pour }p > 0
$$
et que l'homomorphisme
$$
f_*(\mathcal{F} \otimes_{\mathcal{O}_X} \mathcal{L}^{\otimes d})_y
\longrightarrow
H^0(X_y, \mathcal{F}_y \otimes_{\mathcal{O}_{X_y}} \mathcal{L}_y^{\otimes d})
$$
soit surjectif.
\end{lemma}

\begin{proof}
Observons que $\mathcal{O}_{Y, y}$ est un anneau local noethérien.
Considérons le morphisme canonique
$c : \Spec(\mathcal{O}_{Y, y}) \to Y$, voir
Schémas, équation (\ref{schemes-equation-canonical-morphism}).
Ce morphisme est plat, puisqu'il identifie les anneaux locaux.
Notons provisoirement $f' : X' \to \Spec(\mathcal{O}_{Y, y})$
le morphisme obtenu de $f$ par ce changement de base. Nous avons
$c^*R^pf_*\mathcal{F} = R^pf'_*\mathcal{F}'$ d'après
Cohomologie des espaces, lemme
\ref{spaces-cohomology-lemma-flat-base-change-cohomology}.
De plus, les fibres $X_y$ et $X'_y$ s'identifient.
Nous pouvons donc supposer que $Y = \Spec(A)$ est le spectre d'un
anneau local noethérien $(A, \mathfrak m, \kappa)$ et que $y \in Y$
correspond à $\mathfrak m$. Dans ce cas,
$R^pf_*(\mathcal{F} \otimes_{\mathcal{O}_X} \mathcal{L}^{\otimes d})_y =
H^p(X, \mathcal{F} \otimes_{\mathcal{O}_X} \mathcal{L}^{\otimes d})$
pour tout $p \geq 0$. Notons $f_y : X_y \to \Spec(\kappa)$ la projection.

\medskip\noindent
Posons $B = \text{Gr}_\mathfrak m(A) =
\bigoplus_{n \geq 0} \mathfrak m^n/\mathfrak m^{n + 1}$.
Considérons le faisceau $\mathcal{B} = f_y^*\widetilde{B}$
d'algèbres graduées quasi-cohérentes sur $\mathcal{O}_{X_y}$.
Nous emploierons les notations de Cohomologie des espaces, section
\ref{spaces-cohomology-section-theorem-formal-functions},
en remplaçant $I$ par $\mathfrak m$.
Puisque $X_y$ est le sous-espace fermé de $X$ défini par
$\mathfrak m\mathcal{O}_X$, nous pouvons considérer
$\mathfrak m^n\mathcal{F}/\mathfrak m^{n + 1}\mathcal{F}$
comme un $\mathcal{O}_{X_y}$-module cohérent, voir
Cohomologie des espaces, lemme \ref{spaces-cohomology-lemma-i-star-equivalence}.
Alors
$\bigoplus_{n \geq 0} \mathfrak m^n\mathcal{F}/\mathfrak m^{n + 1}\mathcal{F}$
est un $\mathcal{B}$-module gradué quasi-cohérent de type fini,
car il est engendré en degré zéro sur $\mathcal{B}$
et sa partie de degré zéro est
$\mathcal{F}_y = \mathcal{F}/\mathfrak m \mathcal{F}$,
qui est un $\mathcal{O}_{X_y}$-module cohérent.
Ainsi, d'après Cohomologie des schémas, lemme
\ref{coherent-lemma-graded-finiteness}, partie (2),
il existe un $d_0$ tel que
$$
H^p(X_y, \mathfrak m^n \mathcal{F}/ \mathfrak m^{n + 1}\mathcal{F}
\otimes_{\mathcal{O}_{X_y}} \mathcal{L}_y^{\otimes d}) = 0
$$
pour tous $p > 0$, $d \geq d_0$ et $n \geq 0$. D'après
Cohomologie des espaces, lemme
\ref{spaces-cohomology-lemma-relative-affine-cohomology},
cela revient à dire que
$
H^p(X, \mathfrak m^n \mathcal{F}/ \mathfrak m^{n + 1}\mathcal{F}
\otimes_{\mathcal{O}_X} \mathcal{L}^{\otimes d}) = 0
$
pour tous $p > 0$, $d \geq d_0$ et $n \geq 0$.

\medskip\noindent
Considérons les suites exactes courtes
$$
0 \to \mathfrak m^n\mathcal{F}/\mathfrak m^{n + 1} \mathcal{F}
\to \mathcal{F}/\mathfrak m^{n + 1} \mathcal{F}
\to \mathcal{F}/\mathfrak m^n \mathcal{F} \to 0
$$
de $\mathcal{O}_X$-modules cohérents. La tensorisation par $\mathcal{L}^{\otimes d}$
est un foncteur exact et nous obtenons des suites exactes courtes
$$
0 \to
\mathfrak m^n\mathcal{F}/\mathfrak m^{n + 1} \mathcal{F}
\otimes_{\mathcal{O}_X} \mathcal{L}^{\otimes d}
\to \mathcal{F}/\mathfrak m^{n + 1} \mathcal{F}
\otimes_{\mathcal{O}_X} \mathcal{L}^{\otimes d}
\to \mathcal{F}/\mathfrak m^n \mathcal{F}
\otimes_{\mathcal{O}_X} \mathcal{L}^{\otimes d} \to 0
$$
En utilisant la suite exacte longue de cohomologie et l'annulation ci-dessus,
nous concluons (par récurrence) que
\begin{enumerate}
\item $H^p(X, \mathcal{F}/\mathfrak m^n \mathcal{F}
\otimes_{\mathcal{O}_X} \mathcal{L}^{\otimes d}) = 0$
pour tous $p > 0$, $d \geq d_0$ et $n \geq 0$, et
\item $H^0(X, \mathcal{F}/\mathfrak m^n \mathcal{F}
\otimes_{\mathcal{O}_X} \mathcal{L}^{\otimes d}) \to
H^0(X_y, \mathcal{F}_y \otimes_{\mathcal{O}_{X_y}} \mathcal{L}_y^{\otimes d})$
est surjectif pour tous $d \geq d_0$ et $n \geq 1$.
\end{enumerate}
D'après le théorème des fonctions formelles
(Cohomologie des espaces, théorème
\ref{spaces-cohomology-theorem-formal-functions}),
nous trouvons que le complété $\mathfrak m$-adique de
$H^p(X, \mathcal{F} \otimes_{\mathcal{O}_X} \mathcal{L}^{\otimes d})$
est nul pour tous $d \geq d_0$ et $p > 0$.
Comme $H^p(X, \mathcal{F} \otimes_{\mathcal{O}_X} \mathcal{L}^{\otimes d})$
est un $A$-module de type fini d'après
Cohomologie des espaces, lemme
\ref{spaces-cohomology-lemma-proper-over-affine-cohomology-finite},
il résulte du lemme de Nakayama (Algèbre, lemme \ref{algebra-lemma-NAK})
que $H^p(X, \mathcal{F} \otimes_{\mathcal{O}_X} \mathcal{L}^{\otimes d})$
est nul pour tous $d \geq d_0$ et $p > 0$.
Pour $p = 0$, nous déduisons du
lemme de Cohomologie des espaces
\ref{spaces-cohomology-lemma-ML-cohomology-powers-ideal}, partie (3),
que $H^0(X, \mathcal{F} \otimes_{\mathcal{O}_X} \mathcal{L}^{\otimes d}) \to
H^0(X_y, \mathcal{F}_y \otimes_{\mathcal{O}_{X_y}} \mathcal{L}_y^{\otimes d})$
est surjectif, ce qui donne la dernière assertion du lemme.
\end{proof}

\begin{lemma}
\label{lemma-ample-in-neighbourhood}
(Pour une version plus générale, voir
Descente sur les espaces, lemme \ref{spaces-descent-lemma-ample-in-neighbourhood}).
Soit $Y$ un schéma noethérien. Soit $X$ un espace algébrique sur $Y$
tel que le morphisme structural $f : X \to Y$ soit propre.
Soit $\mathcal{L}$ un $\mathcal{O}_X$-module inversible.
Soit $y \in Y$ un point tel que $X_y$ soit un schéma et que
$\mathcal{L}_y$ soit ample sur $X_y$.
Il existe alors un voisinage ouvert $V \subset Y$
de $y$ tel que $\mathcal{L}|_{f^{-1}(V)}$ soit ample sur $f^{-1}(V)/V$
(au sens de la définition \ref{definition-relatively-ample}).
\end{lemma}

\begin{proof}
Choisissons $d_0$ comme dans le lemme \ref{lemma-ample-on-fibre} pour
$\mathcal{F} = \mathcal{O}_X$. Choisissons $d \geq d_0$
de sorte qu'il existe un $r \geq 0$ et des sections
$s_{y, 0}, \ldots, s_{y, r} \in H^0(X_y, \mathcal{L}_y^{\otimes d})$
qui définissent une immersion fermée
$$
\varphi_y =
\varphi_{\mathcal{L}_y^{\otimes d}, (s_{y, 0}, \ldots, s_{y, r})} :
X_y \to \mathbf{P}^r_{\kappa(y)}.
$$
C'est possible d'après Morphismes, lemme
\ref{morphisms-lemma-finite-type-over-affine-ample-very-ample},
mais nous utilisons aussi
Morphismes, lemme \ref{morphisms-lemma-image-proper-scheme-closed},
pour voir que $\varphi_y$ est une immersion fermée, et
Constructions, section \ref{constructions-section-projective-space},
pour la description des morphismes vers l'espace projectif
au moyen de faisceaux inversibles et de sections.
Par notre choix de $d_0$, après avoir remplacé $Y$ par un voisinage ouvert
de $y$, nous pouvons choisir
$s_0, \ldots, s_r \in H^0(X, \mathcal{L}^{\otimes d})$
qui s'envoient sur $s_{y, 0}, \ldots, s_{y, r}$.
Soit $X_{s_i} \subset X$ le sous-espace ouvert où $s_i$
est un générateur de $\mathcal{L}^{\otimes d}$. Puisque
les $s_{y, i}$ engendrent $\mathcal{L}_y^{\otimes d}$, nous avons
$|X_y| \subset U = \bigcup |X_{s_i}|$. Comme $X \to Y$ est fermé,
il existe un voisinage ouvert $y \in V \subset Y$
tel que $|f|^{-1}(V) \subset U$. Après avoir remplacé $Y$ par $V$, nous pouvons
supposer que les $s_i$ engendrent $\mathcal{L}^{\otimes d}$. Nous
obtenons ainsi un morphisme
$$
\varphi = \varphi_{\mathcal{L}^{\otimes d}, (s_0, \ldots, s_r)} :
X \longrightarrow \mathbf{P}^r_Y
$$
avec $\mathcal{L}^{\otimes d} \cong \varphi^*\mathcal{O}_{\mathbf{P}^r_Y}(1)$
dont le changement de base en $y$ donne $\varphi_y$ (à strictement parler, il faudrait
rédiger une preuve que la construction des morphismes vers l'espace projectif
donnée dans
Constructions, section \ref{constructions-section-projective-space},
vaut aussi pour décrire les morphismes d'espaces algébriques vers l'espace
projectif ; nous omettons les détails).

\medskip\noindent
Nous terminerons la démonstration par une astuce ; la démonstration « correcte »
consiste à montrer directement que $\varphi$ est une immersion
fermée après changement de base à un voisinage ouvert de $y$.
En effet, d'après
Cohomologie des espaces, lemme
\ref{spaces-cohomology-lemma-proper-finite-fibre-finite-in-neighbourhood},
nous voyons que $\varphi$ est fini au-dessus d'un voisinage ouvert
de la fibre $\mathbf{P}^r_{\kappa(y)}$ de $\mathbf{P}^r_Y \to Y$
au-dessus de $y$. Comme $\mathbf{P}^r_Y \to Y$ est fermé, quitte à
rétrécir $Y$, nous pouvons supposer que $\varphi$ est fini.
En particulier, $X$ est un schéma.
Alors $\mathcal{L}^{\otimes d} \cong \varphi^*\mathcal{O}_{\mathbf{P}^r_Y}(1)$
est ample d'après le résultat très général donné dans
Morphismes, lemme \ref{morphisms-lemma-pullback-ample-tensor-relatively-ample}.
\end{proof}








\section{Sous-espaces fermés du Proj relatif}
\label{section-closed-in-relative-proj}

\noindent
Quelques lemmes auxiliaires sur les sous-espaces fermés du Proj relatif.
Cette section est l'analogue de
Diviseurs, section \ref{divisors-section-closed-in-relative-proj}.

\begin{lemma}
\label{lemma-closed-subscheme-proj}
Soit $S$ un schéma. Soit $X$ un espace algébrique sur $S$.
Soit $\mathcal{A}$ une $\mathcal{O}_X$-algèbre graduée quasi-cohérente. Soit
$\pi : P = \underline{\text{Proj}}_X(\mathcal{A}) \to X$ le Proj relatif
de $\mathcal{A}$. Soit $i : Z \to P$ un sous-espace fermé. Notons
$\mathcal{I} \subset \mathcal{A}$ le noyau du morphisme canonique
$$
\mathcal{A}
\longrightarrow
\bigoplus\nolimits_{d \geq 0} \pi_*\left((i_*\mathcal{O}_Z)(d)\right)
$$
Si $\pi$ est quasi-compact, il existe un isomorphisme
$Z = \underline{\text{Proj}}_X(\mathcal{A}/\mathcal{I})$.
\end{lemma}

\begin{proof}
Le morphisme $\pi$ est séparé d'après le
lemme \ref{lemma-relative-proj-separated}.
Comme $\pi$ est quasi-compact, $\pi_*$ transforme les modules quasi-cohérents
en modules quasi-cohérents, voir
Morphismes d'espaces, lemme \ref{spaces-morphisms-lemma-pushforward}.
Ainsi, $\mathcal{I}$ est un $\mathcal{O}_X$-module quasi-cohérent.
En particulier, $\mathcal{B} = \mathcal{A}/\mathcal{I}$ est une
$\mathcal{O}_X$-algèbre graduée quasi-cohérente. Le morphisme de fonctorialité
$Z' = \underline{\text{Proj}}_X(\mathcal{B}) \to
\underline{\text{Proj}}_X(\mathcal{A})$ est défini partout et est
une immersion fermée, voir le lemme
\ref{lemma-surjective-graded-rings-map-relative-proj}.
Il suffit donc de montrer que $Z = Z'$ comme sous-espaces fermés de $P$.

\medskip\noindent
Cela étant, la question est locale pour la topologie \'etale sur la base et l'on
se ramène au cas des schémas
(Diviseurs, lemme \ref{divisors-lemma-closed-subscheme-proj})
par localisation \'etale.
\end{proof}

\noindent
Lorsque le sous-espace fermé est localement défini par un nombre fini
d'équations, nous pouvons le définir par un faisceau d'idéaux de type fini de
$\mathcal{A}$.

\begin{lemma}
\label{lemma-closed-subscheme-proj-finite}
Soit $S$ un schéma. Soit $X$ un espace algébrique quasi-compact et quasi-séparé
sur $S$.
Soit $\mathcal{A}$ une $\mathcal{O}_X$-algèbre graduée quasi-cohérente. Soit
$\pi : P = \underline{\text{Proj}}_X(\mathcal{A}) \to X$ le Proj relatif
de $\mathcal{A}$. Soit $i : Z \to P$ un sous-schéma fermé.
Si $\pi$ est quasi-compact et si $i$ est de présentation finie, alors il existe
un $d > 0$ et un sous-$\mathcal{O}_X$-module quasi-cohérent de type fini
$\mathcal{F} \subset \mathcal{A}_d$ tel que
$Z = \underline{\text{Proj}}_X(\mathcal{A}/\mathcal{F}\mathcal{A})$.
\end{lemma}

\begin{proof}
Le lecteur peut reprendre les arguments employés dans le cas des schémas. Nous
allons toutefois montrer, par une astuce, que le lemme résulte de ce cas.
Soit $\mathcal{I} \subset \mathcal{A}$ l'idéal gradué quasi-cohérent
qui définit $Z$ dans le lemme \ref{lemma-closed-subscheme-proj}.
Choisissons un schéma affine $U$ et un morphisme \'etale surjectif
$U \to X$, voir Propriétés des espaces, lemme
\ref{spaces-properties-lemma-quasi-compact-affine-cover}.
D'après le cas des schémas
(Diviseurs, lemme \ref{divisors-lemma-closed-subscheme-proj-finite}),
il existe un $d > 0$ et un
sous-$\mathcal{O}_U$-module quasi-cohérent de type fini
$\mathcal{F}' \subset \mathcal{I}_d|_U \subset \mathcal{A}_d|_U$,
tel que $Z \times_X U$ soit égal à
$\underline{\text{Proj}}_U(\mathcal{A}|_U/\mathcal{F}'\mathcal{A}|_U)$.
D'après Limites des espaces, lemme
\ref{spaces-limits-lemma-directed-colimit-finite-type},
nous pouvons trouver un sous-module quasi-cohérent de type fini
$\mathcal{F} \subset \mathcal{I}_d$ tel que
$\mathcal{F}' \subset \mathcal{F}|_U$. Posons
$Z' = \underline{\text{Proj}}_X(\mathcal{A}/\mathcal{F}\mathcal{A})$.
Alors $Z' \to P$ est une immersion fermée
(lemme \ref{lemma-surjective-generated-degree-1-map-relative-proj})
et $Z \subset Z'$ puisque $\mathcal{F}\mathcal{A} \subset \mathcal{I}$.
D'autre part, $Z' \times_X U \subset Z \times_X U$ d'après notre
choix de $\mathcal{F}$. Ainsi, $Z = Z'$, comme voulu.
\end{proof}

\begin{lemma}
\label{lemma-closed-subscheme-proj-finite-type}
Soit $S$ un schéma. Soit $X$ un espace algébrique quasi-compact et quasi-séparé
sur $S$.
Soit $\mathcal{A}$ une $\mathcal{O}_X$-algèbre graduée quasi-cohérente.
Soit $\pi : P = \underline{\text{Proj}}_X(\mathcal{A}) \to X$ le Proj relatif
de $\mathcal{A}$. Soit $i : Z \to X$ un sous-espace fermé.
Soit $U \subset X$ un ouvert. Supposons que
\begin{enumerate}
\item $\pi$ soit quasi-compact,
\item $i$ soit de présentation finie,
\item $|U| \cap |\pi|(|i|(|Z|)) = \emptyset$,
\item $U$ soit quasi-compact,
\item $\mathcal{A}_n$ soit un $\mathcal{O}_X$-module de type fini pour tout $n$.
\end{enumerate}
Alors il existe un $d > 0$ et un
sous-$\mathcal{O}_X$-module quasi-cohérent de type fini $\mathcal{F} \subset \mathcal{A}_d$ tel que (a)
$Z = \underline{\text{Proj}}_X(\mathcal{A}/\mathcal{F}\mathcal{A})$
et (b) le support de $\mathcal{A}_d/\mathcal{F}$ soit disjoint de $U$.
\end{lemma}

\begin{proof}
Nous employons la même astuce que dans la démonstration du
lemme \ref{lemma-closed-subscheme-proj-finite}
pour nous ramener au cas des schémas.
Soit $\mathcal{I} \subset \mathcal{A}$ l'idéal gradué quasi-cohérent
qui définit $Z$ dans le lemme \ref{lemma-closed-subscheme-proj}.
Choisissons un schéma affine $W$ et un morphisme \'etale surjectif
$W \to X$, voir Propriétés des espaces, lemme
\ref{spaces-properties-lemma-quasi-compact-affine-cover}.
D'après le cas des schémas
(Diviseurs, lemme \ref{divisors-lemma-closed-subscheme-proj-finite-type}),
il existe un $d > 0$ et un
sous-$\mathcal{O}_W$-module quasi-cohérent de type fini
$\mathcal{F}' \subset \mathcal{I}_d|_W \subset \mathcal{A}_d|_W$,
tel que (a) $Z \times_X W$ soit égal à
$\underline{\text{Proj}}_W(\mathcal{A}|_W/\mathcal{F}'\mathcal{A}|_W)$
et que (b) le support de $\mathcal{A}_d|_W/\mathcal{F}'$ soit disjoint de
$U \times_X W$. D'après Limites des espaces, lemme
\ref{spaces-limits-lemma-directed-colimit-finite-type},
nous pouvons trouver un sous-module quasi-cohérent de type fini
$\mathcal{F} \subset \mathcal{I}_d$ tel que
$\mathcal{F}' \subset \mathcal{F}|_W$. Posons
$Z' = \underline{\text{Proj}}_X(\mathcal{A}/\mathcal{F}\mathcal{A})$.
Alors $Z' \to P$ est une immersion fermée
(lemme \ref{lemma-surjective-generated-degree-1-map-relative-proj})
et $Z \subset Z'$ puisque $\mathcal{F}\mathcal{A} \subset \mathcal{I}$.
D'autre part, $Z' \times_X W \subset Z \times_X W$ d'après notre
choix de $\mathcal{F}$. Ainsi, $Z = Z'$.
Enfin, le support de $\mathcal{A}_d/\mathcal{F}$ est contenu dans
$X \setminus U$, car $\mathcal{A}_d|_W/\mathcal{F}|_W$ est un quotient
de $\mathcal{A}_d|_W/\mathcal{F}'$, dont le support est contenu dans
$W \setminus U \times_X W$. Le lemme en résulte.
\end{proof}

\begin{lemma}
\label{lemma-conormal-sheaf-section-projective-bundle}
Soit $S$ un schéma et soit $X$ un espace algébrique sur $S$.
Soit $\mathcal{E}$ un $\mathcal{O}_X$-module quasi-cohérent.
Il existe une bijection
$$
\left\{
\begin{matrix}
\text{sections }\sigma\text{ du } \\
\text{morphisme } \mathbf{P}(\mathcal{E}) \to X
\end{matrix}
\right\}
\leftrightarrow
\left\{
\begin{matrix}
\text{surjections }\mathcal{E} \to \mathcal{L}\text{ où} \\
\mathcal{L}\text{ est un }\mathcal{O}_X\text{-module inversible}
\end{matrix}
\right\}
$$
Dans ce cas, $\sigma$ est une immersion fermée et il existe un
isomorphisme canonique
$$
\Ker(\mathcal{E} \to \mathcal{L})
\otimes_{\mathcal{O}_X} \mathcal{L}^{\otimes -1}
\longrightarrow
\mathcal{C}_{\sigma(X)/\mathbf{P}(\mathcal{E})}
$$
La bijection et l'isomorphisme sont tous deux compatibles avec le changement de base.
\end{lemma}

\begin{proof}
Puisque les constructions sont compatibles avec le changement de base, il suffit de
vérifier l'assertion localement pour la topologie \'etale sur $X$. Nous pouvons donc supposer que $X$ est
un schéma, et le résultat est alors
Diviseurs, lemme \ref{divisors-lemma-conormal-sheaf-section-projective-bundle}.
\end{proof}






\section{Éclatements}
\label{section-blowing-up}

\noindent
L'éclatement est un outil important en géométrie algébrique.

\begin{definition}
\label{definition-blow-up}
Soit $S$ un schéma. Soit $X$ un espace algébrique sur $S$.
Soit $\mathcal{I} \subset \mathcal{O}_X$ un faisceau quasi-cohérent
d'idéaux, et soit $Z \subset X$ le sous-espace fermé correspondant
à $\mathcal{I}$
(Morphismes d'espaces, lemme
\ref{spaces-morphisms-lemma-closed-immersion-ideals}).
L'{\it éclatement de $X$ le long de $Z$}, ou l'
{\it éclatement de $X$ suivant le faisceau d'idéaux $\mathcal{I}$} est
le morphisme
$$
b :
\underline{\text{Proj}}_X
\left(\bigoplus\nolimits_{n \geq 0} \mathcal{I}^n\right)
\longrightarrow
X
$$
Le {\it diviseur exceptionnel} de l'éclatement est l'image inverse
$b^{-1}(Z)$. Le sous-espace $Z$ est parfois appelé le {\it centre} de l'éclatement.
\end{definition}

\noindent
Nous verrons plus loin que le diviseur exceptionnel est un diviseur de Cartier
effectif. De plus, l'éclatement se caractérise comme le « plus petit »
espace algébrique au-dessus de $X$ tel que l'image inverse de $Z$ soit un
diviseur de Cartier effectif.

\medskip\noindent
Si $b : X' \to X$ est l'éclatement de $X$ le long de $Z$, on note souvent
$\mathcal{O}_{X'}(n)$ les faisceaux tordus du faisceau structural. Notons que ceux-ci
sont des $\mathcal{O}_{X'}$-modules inversibles et que
$\mathcal{O}_{X'}(n) = \mathcal{O}_{X'}(1)^{\otimes n}$,
car $X'$ est le Proj relatif d'une algèbre graduée quasi-cohérente
sur $\mathcal{O}_X$ qui est engendrée en degré $1$, voir le
lemme \ref{lemma-relative-proj-generated-in-degree-1}.

\begin{lemma}
\label{lemma-blowing-up-affine}
Soit $S$ un schéma. Soit $X$ un espace algébrique sur $S$.
Soit $\mathcal{I} \subset \mathcal{O}_X$ un
faisceau quasi-cohérent d'idéaux. Soit $U = \Spec(A)$ un schéma affine
\'etale sur $X$, et soit $I \subset A$ l'idéal correspondant à
$\mathcal{I}|_U$. Si $X' \to X$ est l'éclatement de $X$ suivant $\mathcal{I}$,
il existe alors un isomorphisme canonique
$$
U \times_X X' = \text{Proj}(\bigoplus\nolimits_{d \geq 0} I^d)
$$
de schémas au-dessus de $U$, où le membre de droite est
le spectre homogène de l'algèbre de Rees de $I$ dans $A$.
De plus, $U \times_X X'$ possède un recouvrement ouvert affine par les
spectres des algèbres d'éclatement affines $A[\frac{I}{a}]$.
\end{lemma}

\begin{proof}
Notons que la restriction $\mathcal{I}|_U$ est égale à l'image inverse
de $\mathcal{I}$ par le morphisme $U \to X$, voir
Propriétés des espaces, section \ref{spaces-properties-section-modules}.
Le lemme résulte donc de la combinaison du lemme \ref{lemma-relative-proj} et de
Diviseurs, lemme \ref{divisors-lemma-blowing-up-affine}.
\end{proof}

\begin{lemma}
\label{lemma-flat-base-change-blowing-up}
Soit $S$ un schéma.
Soit $X_1 \to X_2$ un morphisme plat d'espaces algébriques sur $S$.
Soit $Z_2 \subset X_2$ un sous-espace fermé.
Soit $Z_1$ l'image inverse de $Z_2$ dans $X_1$.
Soit $X'_i$ l'éclatement de centre $Z_i$ dans $X_i$. Il existe alors un diagramme
cartésien
$$
\xymatrix{
X_1' \ar[r] \ar[d] & X_2' \ar[d] \\
X_1 \ar[r] & X_2
}
$$
d'espaces algébriques sur $S$.
\end{lemma}

\begin{proof}
Soit $\mathcal{I}_2$ le faisceau d'idéaux de $Z_2$ dans $X_2$.
Notons $g : X_1 \to X_2$ le morphisme donné. Le faisceau d'idéaux
$\mathcal{I}_1$ de $Z_1$ est alors l'image de
$g^*\mathcal{I}_2 \to \mathcal{O}_{X_1}$
(voir Morphismes d'espaces, définition
\ref{spaces-morphisms-definition-inverse-image-closed-subspace}
et la discussion qui suit cette définition).
D'après le lemme \ref{lemma-relative-proj-base-change},
$X_1 \times_{X_2} X_2'$ est le Proj relatif de
$\bigoplus_{n \geq 0} g^*\mathcal{I}_2^n$. Comme $g$ est plat, le morphisme
$g^*\mathcal{I}_2^n \to \mathcal{O}_{X_1}$ est injectif et son image est
$\mathcal{I}_1^n$. Ainsi, $X_1 \times_{X_2} X_2' = X_1'$.
\end{proof}

\begin{lemma}
\label{lemma-blowing-up-gives-effective-Cartier-divisor}
Soit $S$ un schéma. Soit $X$ un espace algébrique sur $S$.
Soit $Z \subset X$ un sous-espace fermé.
L'éclatement $b : X' \to X$ de centre $Z$ dans $X$
possède les propriétés suivantes :
\begin{enumerate}
\item $b|_{b^{-1}(X \setminus Z)} : b^{-1}(X \setminus Z) \to X \setminus Z$
est un isomorphisme,
\item le diviseur exceptionnel $E = b^{-1}(Z)$ est un diviseur de Cartier effectif
sur $X'$,
\item il existe un isomorphisme canonique
$\mathcal{O}_{X'}(-1) = \mathcal{O}_{X'}(E)$
\end{enumerate}
\end{lemma}

\begin{proof}
Soit $U$ un schéma et soit $U \to X$ un morphisme \'etale surjectif.
Comme la formation de l'éclatement commute au changement de base plat
(lemme \ref{lemma-flat-base-change-blowing-up}),
nous pouvons démontrer chacune de ces assertions après changement de base à $U$.
Cela nous ramène au cas des schémas.
Dans ce cas, le résultat est
Diviseurs, lemme
\ref{divisors-lemma-blowing-up-gives-effective-Cartier-divisor}.
\end{proof}

\begin{lemma}[Propriété universelle de l'éclatement]
\label{lemma-universal-property-blowing-up}
\begin{slogan}
Éclater un fermé pour en faire un diviseur de Cartier.
\end{slogan}
Soit $S$ un schéma.
Soit $X$ un espace algébrique sur $S$.
Soit $Z \subset X$ un sous-espace fermé.
Soit $\mathcal{C}$ la sous-catégorie pleine de $(\textit{Espaces}/X)$ composée
des $Y \to X$ tels que l'image inverse de $Z$ soit un
diviseur de Cartier effectif sur $Y$. Alors l'éclatement $b : X' \to X$ de centre $Z$ dans $X$
est un objet final de $\mathcal{C}$.
\end{lemma}

\begin{proof}
Le morphisme $b : X' \to X$ est un objet de $\mathcal{C}$ d'après le
lemme \ref{lemma-blowing-up-gives-effective-Cartier-divisor}.
Soit $f : Y \to X$ un objet de $\mathcal{C}$. Nous devons montrer qu'il existe
un unique morphisme $Y \to X'$ au-dessus de $X$. Posons $D = f^{-1}(Z)$.
Soit $\mathcal{I} \subset \mathcal{O}_X$ le faisceau d'idéaux de $Z$
et soit $\mathcal{I}_D$ le faisceau d'idéaux de $D$. Alors
$f^*\mathcal{I} \to \mathcal{I}_D$ est une surjection
vers un $\mathcal{O}_Y$-module inversible. Elle se prolonge en un morphisme
$\psi : \bigoplus f^*\mathcal{I}^d \to \bigoplus \mathcal{I}_D^d$
d'algèbres graduées sur $\mathcal{O}_Y$. (Nous observons que
$\mathcal{I}_D^d = \mathcal{I}_D^{\otimes d}$ puisque $D$ est un
diviseur de Cartier effectif.) D'après le
lemme \ref{lemma-relative-proj-generated-in-degree-1},
le triplet $(f : Y \to X, \mathcal{I}_D, \psi)$ définit un
morphisme $Y \to X'$ au-dessus de $X$. La restriction
$$
Y \setminus D \longrightarrow X' \setminus b^{-1}(Z) = X \setminus Z
$$
est unique. L'ouvert $Y \setminus D$ est schématiquement dense dans $Y$
d'après le lemme \ref{lemma-complement-effective-Cartier-divisor}. 
Ainsi, le morphisme $Y \to X'$ est unique d'après
Morphismes d'espaces, lemme \ref{spaces-morphisms-lemma-equality-of-morphisms}
(de plus, $b$ est séparé d'après le lemme
\ref{lemma-relative-proj-separated}).
\end{proof}

\begin{lemma}
\label{lemma-blow-up-effective-Cartier-divisor}
Soit $S$ un schéma. Soit $X$ un espace algébrique sur $S$.
Soit $Z \subset X$ un diviseur de Cartier effectif.
L'éclatement de $X$ le long de $Z$ est le morphisme identité de $X$.
\end{lemma}

\begin{proof}
Cela résulte immédiatement de la propriété universelle des éclatements
(lemme \ref{lemma-universal-property-blowing-up}).
\end{proof}

\begin{lemma}
\label{lemma-blow-up-reduced-space}
Soit $S$ un schéma. Soit $X$ un espace algébrique sur $S$.
Soit $\mathcal{I} \subset \mathcal{O}_X$ un
faisceau quasi-cohérent d'idéaux. Si $X$ est réduit, alors
l'éclatement $X'$ de $X$ suivant $\mathcal{I}$ est réduit.
\end{lemma}

\begin{proof}
Soit $U$ un schéma et soit $U \to X$ un morphisme \'etale surjectif.
Comme la formation de l'éclatement commute au changement de base plat
(lemme \ref{lemma-flat-base-change-blowing-up}),
nous pouvons démontrer l'assertion après changement de base à $U$.
Cela nous ramène au cas des schémas.
Dans ce cas, le résultat est
Diviseurs, lemme \ref{divisors-lemma-blow-up-reduced-scheme}.
\end{proof}

\begin{lemma}
\label{lemma-blowup-finite-nr-irreducibles}
Soit $S$ un schéma. Soit $X$ un espace algébrique sur $S$. Soit
$b : X' \to X$ l'éclatement de $X$ le long d'un sous-espace fermé. Si
$X$ vérifie les conditions équivalentes de
Morphismes d'espaces, lemme \ref{spaces-morphisms-lemma-prepare-normalization},
alors il en est de même de $X'$.
\end{lemma}

\begin{proof}
Cela résulte immédiatement du lemme cité dans l'énoncé,
de la description locale pour la topologie \'etale des éclatements donnée au
lemme \ref{lemma-blowing-up-affine}, et de
Diviseurs, lemme \ref{divisors-lemma-blow-up-and-irreducible-components}.
\end{proof}

\begin{lemma}
\label{lemma-blow-up-pullback-effective-Cartier}
Soit $S$ un schéma. Soit $X$ un espace algébrique sur $S$.
Soit $b : X' \to X$ un éclatement de $X$ le long d'un sous-espace fermé.
Pour tout diviseur de Cartier effectif $D$ sur $X$, l'image inverse
$b^{-1}D$ est définie (voir la définition
\ref{definition-pullback-effective-Cartier-divisor}).
\end{lemma}

\begin{proof}
D'après les lemmes \ref{lemma-blowing-up-affine} et
\ref{lemma-characterize-effective-Cartier-divisor},
cela se ramène au fait algébrique suivant :
soient $A$ un anneau, $I \subset A$ un idéal, $a \in I$, et $x \in A$
un non-diviseur de zéro. L'image de $x$ dans $A[\frac{I}{a}]$ est alors un
non-diviseur de zéro. En effet, supposons que $x (y/a^n) = 0$ dans $A[\frac{I}{a}]$.
Alors $a^mxy = 0$ dans $A$ pour un certain $m$. Donc $a^my = 0$ puisque $x$ est un
non-diviseur de zéro. Ainsi $y/a^n$ est nul dans $A[\frac{I}{a}]$, comme voulu.
\end{proof}

\begin{lemma}
\label{lemma-blowing-up-two-ideals}
Soit $S$ un schéma. Soit $X$ un espace algébrique sur $S$.
Soient $\mathcal{I} \subset \mathcal{O}_X$ et $\mathcal{J}$ des
faisceaux quasi-cohérents d'idéaux. Soit $b : X' \to X$ l'éclatement
de $X$ suivant $\mathcal{I}$. Soit $b' : X'' \to X'$ l'éclatement de
$X'$ suivant $b^{-1}\mathcal{J} \mathcal{O}_{X'}$. Alors $X'' \to X$
est canoniquement isomorphe à l'éclatement de $X$ suivant $\mathcal{I}\mathcal{J}$.
\end{lemma}

\begin{proof}
Soit $E \subset X'$ le diviseur exceptionnel de $b$, qui est un diviseur de
Cartier effectif d'après le
lemme \ref{lemma-blowing-up-gives-effective-Cartier-divisor}.
Alors $(b')^{-1}E$ est un diviseur de Cartier effectif sur $X''$ d'après le
lemme \ref{lemma-blow-up-pullback-effective-Cartier}.
Soit $E' \subset X''$ le diviseur exceptionnel de $b'$ (lui aussi un diviseur de
Cartier effectif). Considérons le diviseur de Cartier effectif
$E'' = E' + (b')^{-1}E$. Par construction, l'idéal de $E''$ est
$(b \circ b')^{-1}\mathcal{I} (b \circ b')^{-1}\mathcal{J} \mathcal{O}_{X''}$.
D'après le lemme \ref{lemma-universal-property-blowing-up},
il existe donc un morphisme canonique de $X''$ vers l'éclatement $c : Y \to X$
de $X$ le long de $\mathcal{I}\mathcal{J}$. Inversement, puisque $\mathcal{I}\mathcal{J}$
a pour image inverse un idéal inversible, on voit que
$c^{-1}\mathcal{I}\mathcal{O}_Y$ définit
un diviseur de Cartier effectif, voir le
lemme \ref{lemma-sum-closed-subschemes-effective-Cartier}.
Il existe donc un morphisme $c' : Y \to X'$ au-dessus de $X$ par le
lemme \ref{lemma-universal-property-blowing-up}.
Alors $(c')^{-1}b^{-1}\mathcal{J}\mathcal{O}_Y = c^{-1}\mathcal{J}\mathcal{O}_Y$,
qui définit lui aussi un diviseur de Cartier effectif. Il existe donc un morphisme
$c'' : Y \to X''$ au-dessus de $X'$. Nous omettons de vérifier que ce
morphisme est l'inverse du morphisme $X'' \to Y$ construit précédemment.
\end{proof}

\begin{lemma}
\label{lemma-blowing-up-projective}
Soit $S$ un schéma. Soit $X$ un espace algébrique sur $S$.
Soit $\mathcal{I} \subset \mathcal{O}_X$ un faisceau quasi-cohérent
d'idéaux. Soit $b : X' \to X$ l'éclatement de $X$
suivant le faisceau d'idéaux $\mathcal{I}$. Si $\mathcal{I}$ est de type fini, alors
$b : X' \to X$ est un morphisme propre.
\end{lemma}

\begin{proof}
Soit $U$ un schéma et soit $U \to X$ un morphisme \'etale surjectif.
Comme la formation de l'éclatement commute au changement de base plat
(lemme \ref{lemma-flat-base-change-blowing-up}),
nous pouvons démontrer l'assertion après changement de base à $U$
(voir Morphismes d'espaces, lemme
\ref{spaces-morphisms-lemma-proper-local}).
Cela nous ramène au cas des schémas.
Dans ce cas, le morphisme $b$ est projectif d'après
Diviseurs, lemme \ref{divisors-lemma-blowing-up-projective},
donc propre d'après
Morphismes, lemme \ref{morphisms-lemma-locally-projective-proper}.
\end{proof}

\begin{lemma}
\label{lemma-composition-finite-type-blowups}
Soit $S$ un schéma et soit $X$ un espace algébrique sur $S$.
Supposons que $X$ soit quasi-compact et quasi-séparé.
Soit $Z \subset X$ un sous-espace fermé de présentation finie.
Soit $b : X' \to X$ l'éclatement de centre $Z$.
Soit $Z' \subset X'$ un sous-espace fermé de présentation finie.
Soit $X'' \to X'$ l'éclatement de centre $Z'$.
Il existe un sous-espace fermé $Y \subset X$ de présentation finie
tel que
\begin{enumerate}
\item $|Y| = |Z| \cup |b|(|Z'|)$, et
\item la composée $X'' \to X$ soit isomorphe à l'éclatement
de $X$ le long de $Y$.
\end{enumerate}
\end{lemma}

\begin{proof}
La condition que $Z \to X$ soit de présentation finie signifie que
$Z$ est défini par un faisceau d'idéaux quasi-cohérent de type fini
$\mathcal{I} \subset \mathcal{O}_X$, voir
Morphismes d'espaces, lemme
\ref{spaces-morphisms-lemma-closed-immersion-finite-presentation}.
Écrivons $\mathcal{A} = \bigoplus_{n \geq 0} \mathcal{I}^n$, de sorte que
$X' = \underline{\text{Proj}}(\mathcal{A})$.
Notons que $X \setminus Z$ est un sous-espace ouvert quasi-compact de $X$ d'après
Limites des espaces, lemme
\ref{spaces-limits-lemma-quasi-coherent-finite-type-ideals}.
Puisque $b^{-1}(X \setminus Z) \to X \setminus Z$ est un isomorphisme
(lemme \ref{lemma-blowing-up-gives-effective-Cartier-divisor}), le même
résultat montre que
$b^{-1}(X \setminus Z) \setminus Z'$ est un sous-espace ouvert quasi-compact de $X'$.
Ainsi, $U = X \setminus (Z \cup b(Z'))$ est un sous-espace ouvert quasi-compact de $X$.
D'après le lemme \ref{lemma-closed-subscheme-proj-finite-type},
il existe un $d > 0$ et un
sous-$\mathcal{O}_X$-module de type fini $\mathcal{F} \subset \mathcal{I}^d$ tel
que $Z' = \underline{\text{Proj}}(\mathcal{A}/\mathcal{F}\mathcal{A})$
et que le support de $\mathcal{I}^d/\mathcal{F}$ soit contenu
dans $X \setminus U$.

\medskip\noindent
Comme $\mathcal{F} \subset \mathcal{I}^d$ est un sous-$\mathcal{O}_X$-module,
nous pouvons considérer $\mathcal{F} \subset \mathcal{I}^d \subset \mathcal{O}_X$
comme un faisceau d'idéaux quasi-cohérent de type fini sur $X$. Notons-le
$\mathcal{J} \subset \mathcal{O}_X$ pour éviter toute confusion. Puisque
$\mathcal{I}^d / \mathcal{J}$ et $\mathcal{O}/\mathcal{I}^d$ ont
leur support dans $|X| \setminus |U|$, on voit que $|V(\mathcal{J})|$ est contenu
dans $|X| \setminus |U|$. Réciproquement, comme $\mathcal{J} \subset \mathcal{I}^d$,
on voit que $|Z| \subset |V(\mathcal{J})|$. Sur
$X \setminus Z \cong X' \setminus b^{-1}(Z)$, le faisceau d'idéaux
$\mathcal{J}$ définit $Z'$ (voir la formule affichée ci-dessous). Ainsi,
$|V(\mathcal{J})|$ est égal à $|Z| \cup |b|(|Z'|)$. Par suite,
$|V(\mathcal{I}\mathcal{J})| = |Z| \cup |b|(|Z'|)$. De plus,
$\mathcal{I}\mathcal{J}$ est un idéal de type fini, comme produit de deux tels idéaux.
Nous affirmons que $X'' \to X$ est isomorphe à l'éclatement de $X$ suivant
$\mathcal{I}\mathcal{J}$, ce qui achève la démonstration du lemme en posant
$Y = V(\mathcal{I}\mathcal{J})$.

\medskip\noindent
Rappelons d'abord que l'éclatement de $X$ suivant $\mathcal{I}\mathcal{J}$
est le même que l'éclatement de $X'$ suivant $b^{-1}\mathcal{J} \mathcal{O}_{X'}$,
voir le lemme \ref{lemma-blowing-up-two-ideals}.
Il suffit donc de montrer que l'éclatement de $X'$ suivant
$b^{-1}\mathcal{J} \mathcal{O}_{X'}$ coïncide avec l'éclatement de $X'$
le long de $Z'$. Nous allons montrer que
$$
b^{-1}\mathcal{J} \mathcal{O}_{X'} = \mathcal{I}_E^d \mathcal{I}_{Z'}
$$
comme faisceaux d'idéaux sur $X''$. Cela donnera le résultat, puisque
$\mathcal{I}_E^d$ définit le diviseur de Cartier effectif $dE$
et que nous pouvons utiliser les lemmes \ref{lemma-blow-up-effective-Cartier-divisor} et
\ref{lemma-blowing-up-two-ideals}.

\medskip\noindent
Pour établir l'égalité affichée des idéaux, nous pouvons travailler localement.
Avec les notations $A$, $I$, $a \in I$ du lemme \ref{lemma-blowing-up-affine},
on voit que $\mathcal{F}$ correspond à un sous-$R$-module $M \subset I^d$
qui s'envoie isomorphiquement sur un idéal $J \subset R$. La condition
$Z' = \underline{\text{Proj}}(\mathcal{A}/\mathcal{F}\mathcal{A})$
signifie que $Z' \cap \Spec(A[\frac{I}{a}])$ est défini par l'idéal
engendré par les éléments $m/a^d$, où $m \in M$. Supposons que $m \in M$
corresponde à la fonction $f \in J$. Alors, dans l'algèbre d'éclatement affine
$A' = A[\frac{I}{a}]$, on a $f = (a^dm)/a^d = a^d (m/a^d)$.
L'égalité en résulte.
\end{proof}









\section{Transformée stricte}
\label{section-strict-transform}

\noindent
Cette section est l'analogue de
Diviseurs, section \ref{divisors-section-strict-transform}.
Soient $S$ un schéma, $B$ un espace algébrique sur $S$, et
$Z \subset B$ un sous-espace fermé.
Soit $b : B' \to B$ l'éclatement de $B$ le long de $Z$ et désignons par $E \subset B'$
le diviseur exceptionnel $E = b^{-1}Z$. Dans la suite, nous considérerons
souvent un espace algébrique $X$ sur $B$ et formerons le diagramme cartésien
$$
\xymatrix{
\text{pr}_{B'}^{-1}E \ar[r] \ar[d] &
X \times_B B' \ar[r]_-{\text{pr}_X} \ar[d]_{\text{pr}_{B'}} &
X \ar[d]^f \\
E \ar[r] & B' \ar[r] & B
}
$$
Puisque $E$ est un diviseur de Cartier effectif
(lemme \ref{lemma-blowing-up-gives-effective-Cartier-divisor}),
on voit que $\text{pr}_{B'}^{-1}E \subset X \times_B B'$
est localement principal
(lemme \ref{lemma-pullback-locally-principal}).
Ainsi, le morphisme d'inclusion du complémentaire de
$\text{pr}_{B'}^{-1}E$ dans $X \times_B B'$
est affine et, en particulier, quasi-compact
(lemme \ref{lemma-complement-locally-principal-closed-subscheme}).
Par conséquent, pour un $\mathcal{O}_{X \times_B B'}$-module quasi-cohérent
$\mathcal{G}$, le sous-faisceau des sections à support dans $|\text{pr}_{B'}^{-1}E|$
est un sous-module quasi-cohérent ; voir
Limites d'espaces, définition
\ref{spaces-limits-definition-subsheaf-sections-supported-on-closed}.
Si $\mathcal{G}$ est un faisceau quasi-cohérent d'algèbres, par exemple
$\mathcal{G} = \mathcal{O}_{X \times_B B'}$, ce sous-faisceau est alors un idéal
de $\mathcal{G}$.

\begin{definition}
\label{definition-strict-transform}
Reprenons $Z \subset B$ et $f : X \to B$ comme ci-dessus.
\begin{enumerate}
\item Étant donné un $\mathcal{O}_X$-module quasi-cohérent $\mathcal{F}$,
la {\it transformée stricte} de $\mathcal{F}$ relativement à l'éclatement
de $B$ le long de $Z$ est le quotient $\mathcal{F}'$ de $\text{pr}_X^*\mathcal{F}$
par le sous-module des sections à support dans $|\text{pr}_{B'}^{-1}E|$.
\item La {\it transformée stricte} de $X$ est le sous-espace fermé
$X' \subset X \times_B B'$ défini par l'idéal quasi-cohérent des
sections de $\mathcal{O}_{X \times_B B'}$ à support dans
$|\text{pr}_{B'}^{-1}E|$.
\end{enumerate}
\end{definition}

\noindent
Notons que la transformée stricte par un éclatement dépend du
sous-espace fermé qui sert de centre à cet éclatement
(et non du seul morphisme $B' \to B$).

\begin{lemma}[Localisation étale et transformée stricte]
\label{lemma-strict-transform-local}
Dans la situation de la définition \ref{definition-strict-transform}.
Soit
$$
\xymatrix{
U \ar[r] \ar[d] & X \ar[d] \\
V \ar[r] & B
}
$$
un diagramme commutatif de morphismes, où $U$ et $V$ sont des schémas et où
les flèches horizontales sont étales. Soit $V' \to V$ l'éclatement de $V$
le long de $Z \times_B V$. Alors
\begin{enumerate}
\item $V' = V \times_B B'$ et les morphismes
$V' \to B'$ et $U \times_V V' \to X \times_B B'$ sont étales,
\item la transformée stricte $U'$ de $U$ relativement à $V' \to V$
est égale à $X' \times_X U$, où $X'$ est la transformée stricte de $X$
relativement à $B' \to B$, et
\item pour un $\mathcal{O}_X$-module quasi-cohérent $\mathcal{F}$, la
restriction de la transformée stricte $\mathcal{F}'$ à
$U \times_V V'$ est la transformée stricte de $\mathcal{F}|_U$ relativement
à $V' \to V$.
\end{enumerate}
\end{lemma}

\begin{proof}
L'assertion (1) résulte de la commutation de l'éclatement au changement de base
plat (lemme \ref{lemma-flat-base-change-blowing-up}), de la platitude des
morphismes étales et de leur stabilité par changement de base. L'assertion (3)
résulte alors de ce que la formation du faisceau des sections à support dans
un fermé commute à l'image inverse
par les morphismes étales ; voir Limites d'espaces, lemme
\ref{spaces-limits-lemma-sections-supported-on-closed-subset}.
L'assertion (2) résulte de (3), appliquée à $\mathcal{F} = \mathcal{O}_X$.
\end{proof}

\begin{lemma}
\label{lemma-strict-transform}
Dans la situation de la définition \ref{definition-strict-transform}.
\begin{enumerate}
\item La transformée stricte $X'$ de $X$ est l'éclatement de $X$ le long du
sous-espace fermé $f^{-1}Z$ de $X$.
\item Pour un $\mathcal{O}_X$-module quasi-cohérent $\mathcal{F}$, la
transformée stricte $\mathcal{F}'$ est canoniquement isomorphe à
l'image directe par $X' \to X \times_B B'$ de la transformée stricte de
$\mathcal{F}$ relativement à l'éclatement $X' \to X$.
\end{enumerate}
\end{lemma}

\begin{proof}
Soit $X'' \to X$ l'éclatement de $X$ le long de $f^{-1}Z$. Par la propriété
universelle de l'éclatement (lemme \ref{lemma-universal-property-blowing-up}),
il existe un diagramme commutatif
$$
\xymatrix{
X'' \ar[r] \ar[d] & X \ar[d] \\
B' \ar[r] & B
}
$$
d'où un morphisme $i : X'' \to X \times_B B'$. La première assertion
du lemme affirme que $i$ est une immersion fermée d'image $X'$.
La seconde affirme que $\mathcal{F}' = i_*\mathcal{F}''$,
où $\mathcal{F}''$ est la transformée stricte de $\mathcal{F}$
relativement à l'éclatement $X'' \to X$. Ces assertions se vérifient
localement pour la topologie étale sur $X$ ; on se ramène donc au cas des schémas
(Diviseurs, lemme \ref{divisors-lemma-strict-transform}).
Certains détails sont omis.
\end{proof}

\begin{lemma}
\label{lemma-strict-transform-flat}
Dans la situation de la définition \ref{definition-strict-transform}.
\begin{enumerate}
\item Si $X$ est plat sur $B$ en tout point situé au-dessus de $Z$, alors
la transformée stricte de $X$ coïncide avec le changement de base $X \times_B B'$.
\item Soit $\mathcal{F}$ un $\mathcal{O}_X$-module quasi-cohérent.
Si $\mathcal{F}$ est plat sur $B$ en tout point situé au-dessus de $Z$, alors
la transformée stricte $\mathcal{F}'$ de $\mathcal{F}$ est égale à
l'image inverse $\text{pr}_X^*\mathcal{F}$.
\end{enumerate}
\end{lemma}

\begin{proof}
Omis. Indication : cela résulte du cas des schémas
(Diviseurs, lemme \ref{divisors-lemma-strict-transform-flat})
par localisation étale
(lemme \ref{lemma-strict-transform-local}).
\end{proof}

\begin{lemma}
\label{lemma-strict-transform-affine}
Soit $S$ un schéma. Soit $B$ un espace algébrique sur $S$.
Soit $Z \subset B$ un sous-espace fermé.
Soit $b : B' \to B$ l'éclatement de centre $Z$ dans $B$. Soit
$g : X \to Y$ un morphisme affine d'espaces algébriques sur $B$.
Soit $\mathcal{F}$ un faisceau quasi-cohérent sur $X$.
Soit $g' : X \times_B B' \to Y \times_B B'$ le changement de base
de $g$. Soit $\mathcal{F}'$ la transformée stricte de $\mathcal{F}$
relativement à $b$. Alors $g'_*\mathcal{F}'$ est la transformée stricte
de $g_*\mathcal{F}$.
\end{lemma}

\begin{proof}
Omis. Indication : cela résulte du cas des schémas
(Diviseurs, lemme \ref{divisors-lemma-strict-transform-affine})
par localisation étale (lemme \ref{lemma-strict-transform-local}).
\end{proof}

\begin{lemma}
\label{lemma-strict-transform-different-centers}
Soit $S$ un schéma. Soit $B$ un espace algébrique sur $S$.
Soit $Z \subset B$ un sous-espace fermé.
Soit $D \subset B$ un diviseur de Cartier effectif.
Soit $Z' \subset B$ le sous-espace fermé défini par le produit
des faisceaux d'idéaux de $Z$ et de $D$.
Soit $B' \to B$ l'éclatement de $B$ le long de $Z$.
\begin{enumerate}
\item L'éclatement de $B$ le long de $Z'$ est isomorphe au morphisme $B' \to B$.
\item Soient $f : X \to B$ un morphisme d'espaces algébriques et $\mathcal{F}$
un $\mathcal{O}_X$-module quasi-cohérent. Si le sous-faisceau de $\mathcal{F}$ des
sections à support dans $|f^{-1}D|$ est nul, alors la
transformée stricte de $\mathcal{F}$ relativement à l'éclatement
le long de $Z$ coïncide avec la transformée stricte de $\mathcal{F}$ relativement
à l'éclatement de $B$ le long de $Z'$.
\end{enumerate}
\end{lemma}

\begin{proof}
Omis. Indication : cela résulte du cas des schémas
(Diviseurs, lemme \ref{divisors-lemma-strict-transform-different-centers})
par localisation étale (lemme \ref{lemma-strict-transform-local}).
\end{proof}

\begin{lemma}
\label{lemma-strict-transform-composition-blowups}
Soit $S$ un schéma. Soit $B$ un espace algébrique sur $S$.
Soit $Z \subset B$ un sous-espace fermé.
Soit $b : B' \to B$ l'éclatement de centre $Z$.
Soit $Z' \subset B'$ un sous-espace fermé.
Soit $B'' \to B'$ l'éclatement de centre $Z'$.
Soit $Y \subset B$ un sous-schéma fermé tel que
$|Y| = |Z| \cup |b|(|Z'|)$ et que la composée $B'' \to B$
soit isomorphe à l'éclatement de $B$ le long de $Y$.
Dans cette situation, pour tout schéma $X$ sur $B$ et tout
$\mathcal{F} \in \QCoh(\mathcal{O}_X)$, on a
\begin{enumerate}
\item la transformée stricte de $\mathcal{F}$ relativement à l'éclatement
de $B$ le long de $Y$ est égale à la transformée stricte, relativement à
l'éclatement $B'' \to B'$ le long de $Z'$, de la transformée stricte de $\mathcal{F}$
relativement à l'éclatement $B' \to B$ de $B$ le long de $Z$, et
\item la transformée stricte de $X$ relativement à l'éclatement
de $B$ le long de $Y$ est égale à la transformée stricte, relativement à
l'éclatement $B'' \to B'$ le long de $Z'$, de la transformée stricte de $X$
relativement à l'éclatement $B' \to B$ de $B$ le long de $Z$.
\end{enumerate}
\end{lemma}

\begin{proof}
Omis. Indication : cela résulte du cas des schémas
(Diviseurs, lemme \ref{divisors-lemma-strict-transform-composition-blowups})
par localisation étale (lemme \ref{lemma-strict-transform-local}).
\end{proof}

\begin{lemma}
\label{lemma-strict-transform-universally-injective}
Dans la situation de la définition \ref{definition-strict-transform}.
Supposons que
$$
0 \to \mathcal{F}_1 \to \mathcal{F}_2 \to \mathcal{F}_3 \to 0
$$
soit une suite exacte de faisceaux quasi-cohérents sur $X$ qui reste
exacte après tout changement de base $T \to B$. Alors les transformées strictes
$\mathcal{F}_i'$ relativement à tout éclatement $B' \to B$
forment elles aussi la suite exacte
$0 \to \mathcal{F}'_1 \to \mathcal{F}'_2 \to \mathcal{F}'_3 \to 0$.
\end{lemma}

\begin{proof}
Omis. Indication : cela résulte du cas des schémas
(Diviseurs, lemme \ref{divisors-lemma-strict-transform-universally-injective})
par localisation étale (lemme \ref{lemma-strict-transform-local}).
\end{proof}

\begin{lemma}
\label{lemma-strict-transform-blowup-fitting-ideal}
Soit $S$ un schéma. Soit $B$ un espace algébrique sur $S$.
Soit $\mathcal{F}$ un $\mathcal{O}_B$-module quasi-cohérent de type fini.
Soit $Z_k \subset S$ le sous-schéma fermé défini par
$\text{Fit}_k(\mathcal{F})$ ; voir la section \ref{section-fitting-ideals}.
Soit $B' \to B$ l'éclatement de $B$ le long de $Z_k$ et soit
$\mathcal{F}'$ la transformée stricte de $\mathcal{F}$.
Alors $\mathcal{F}'$ peut être engendré localement par $\leq k$
sections.
\end{lemma}

\begin{proof}
Omis. Cela résulte du cas des schémas
(Diviseurs, lemme \ref{divisors-lemma-strict-transform-blowup-fitting-ideal})
par localisation étale (lemme \ref{lemma-strict-transform-local}).
\end{proof}

\begin{lemma}
\label{lemma-strict-transform-blowup-fitting-ideal-locally-free}
Soit $S$ un schéma. Soit $B$ un espace algébrique sur $S$.
Soit $\mathcal{F}$ un $\mathcal{O}_B$-module quasi-cohérent de type fini.
Soit $Z_k \subset S$ le sous-schéma fermé défini par
$\text{Fit}_k(\mathcal{F})$ ; voir la section \ref{section-fitting-ideals}.
Supposons que $\mathcal{F}$ soit localement libre de rang $k$ sur $B \setminus Z_k$.
Soit $B' \to B$ l'éclatement de $B$ le long de $Z_k$ et soit
$\mathcal{F}'$ la transformée stricte de $\mathcal{F}$.
Alors $\mathcal{F}'$ est localement libre de rang $k$.
\end{lemma}

\begin{proof}
Omis. Cela résulte du cas des schémas
(Diviseurs, lemme
\ref{divisors-lemma-strict-transform-blowup-fitting-ideal-locally-free})
par localisation étale (lemme \ref{lemma-strict-transform-local}).
\end{proof}












\section{Éclatements admissibles}
\label{section-admissible-blowups}

\noindent
Afin de mieux maîtriser les éclatements considérés, nous introduisons la
terminologie standard suivante.

\begin{definition}
\label{definition-admissible-blowup}
Soit $S$ un schéma. Soit $X$ un espace algébrique sur $S$.
Soit $U \subset X$ un sous-espace ouvert. Un morphisme
$X' \to X$ est appelé {\it éclatement $U$-admissible} s'il existe une
immersion fermée $Z \to X$ de présentation finie, avec $Z$ disjoint de
$U$, telle que $X'$ soit isomorphe à l'éclatement de $X$ le long de $Z$.
\end{definition}

\noindent
Rappelons que $Z \to X$ est de présentation finie si et seulement si le
faisceau d'idéaux $\mathcal{I}_Z \subset \mathcal{O}_X$ est de type fini ; voir
Morphismes d'espaces, lemme
\ref{spaces-morphisms-lemma-closed-immersion-finite-presentation}.
En particulier, un éclatement $U$-admissible est un morphisme propre ; voir le
lemme \ref{lemma-blowing-up-projective}.
Notons que plusieurs centres peuvent donner lieu au même morphisme.
La condition exige donc seulement l'existence d'un centre disjoint de
$U$ qui produise $X'$.
Enfin, comme le morphisme $b : X' \to X$ est un isomorphisme au-dessus de $U$ (voir le
lemme \ref{lemma-blowing-up-gives-effective-Cartier-divisor}), nous identifierons souvent,
par abus de notation, $U$ à un sous-espace ouvert de $X'$.

\begin{lemma}
\label{lemma-composition-admissible-blowups}
Soit $S$ un schéma.
Soit $X$ un espace algébrique quasi-compact et quasi-séparé sur $S$.
Soit $U \subset X$ un sous-espace ouvert quasi-compact.
Soit $b : X' \to X$ un éclatement $U$-admissible.
Soit $X'' \to X'$ un éclatement $U$-admissible.
Alors la composée $X'' \to X$ est un éclatement $U$-admissible.
\end{lemma}

\begin{proof}
Cela résulte immédiatement du
lemme plus précis \ref{lemma-composition-finite-type-blowups}.
\end{proof}

\begin{lemma}
\label{lemma-extend-admissible-blowups}
Soit $S$ un schéma.
Soit $X$ un espace algébrique quasi-compact et quasi-séparé.
Soient $U, V \subset X$ des sous-espaces ouverts quasi-compacts.
Soit $b : V' \to V$ un éclatement $U \cap V$-admissible.
Alors il existe un éclatement $U$-admissible $X' \to X$
dont la restriction à $V$ est $V'$.
\end{lemma}

\begin{proof}
Soit $\mathcal{I} \subset \mathcal{O}_V$ le faisceau quasi-cohérent d'idéaux
de type fini tel que $V(\mathcal{I})$ soit
disjoint de $U \cap V$ et que $V'$ soit isomorphe à
l'éclatement de $V$ le long de $\mathcal{I}$. Soit
$\mathcal{I}' \subset \mathcal{O}_{U \cup V}$ le faisceau quasi-cohérent
d'idéaux dont la restriction à $U$ est $\mathcal{O}_U$ et
dont la restriction à $V$ est $\mathcal{I}$.
Par Limites d'espaces, lemme \ref{spaces-limits-lemma-extend},
il existe un faisceau quasi-cohérent d'idéaux de type fini
$\mathcal{J} \subset \mathcal{O}_X$ dont la restriction à $U \cup V$ est
$\mathcal{I}'$. Le lemme en résulte.
\end{proof}

\begin{lemma}
\label{lemma-dominate-admissible-blowups}
Soit $S$ un schéma.
Soit $X$ un espace algébrique quasi-compact et quasi-séparé sur $S$.
Soit $U \subset X$ un sous-espace ouvert quasi-compact.
Soient $b_i : X_i \to X$, $i = 1, \ldots, n$, des éclatements $U$-admissibles.
Il existe un éclatement $U$-admissible $b : X' \to X$ tel que
(a) $b$ se factorise sous la forme $X' \to X_i \to X$ pour $i = 1, \ldots, n$, et
(b) chacun des morphismes $X' \to X_i$ soit un éclatement $U$-admissible.
\end{lemma}

\begin{proof}
Soit $\mathcal{I}_i \subset \mathcal{O}_X$ le faisceau quasi-cohérent d'idéaux
de type fini tel que $V(\mathcal{I}_i)$ soit
disjoint de $U$ et que $X_i$ soit isomorphe à
l'éclatement de $X$ le long de $\mathcal{I}_i$. Posons
$\mathcal{I} = \mathcal{I}_1 \cdot \ldots \cdot \mathcal{I}_n$
et soit $X'$ l'éclatement de $X$ le long de $\mathcal{I}$. Alors
$X' \to X$ se factorise par $b_i$ d'après le lemme \ref{lemma-blowing-up-two-ideals}.
\end{proof}

\begin{lemma}
\label{lemma-separate-disjoint-opens-by-blowing-up}
Soit $S$ un schéma.
Soit $X$ un espace algébrique quasi-compact et quasi-séparé sur $S$.
Soient $U, V$ des sous-espaces ouverts quasi-compacts disjoints de $X$.
Il existe alors un éclatement $U \cup V$-admissible $b : X' \to X$
tel que $X'$ soit une réunion disjointe de sous-espaces ouverts
$X' = X'_1 \amalg X'_2$, avec $b^{-1}(U) \subset X'_1$ et
$b^{-1}(V) \subset X'_2$.
\end{lemma}

\begin{proof}
Choisissons des faisceaux quasi-cohérents d'idéaux de type fini $\mathcal{I}$
et, respectivement, $\mathcal{J}$ tels que $X \setminus U = V(\mathcal{I})$
et, respectivement, $X \setminus V = V(\mathcal{J})$ ; voir
Limites d'espaces, lemme
\ref{spaces-limits-lemma-quasi-coherent-finite-type-ideals}.
Alors $|V(\mathcal{I}\mathcal{J})| = |X|$. Par suite,
$\mathcal{I}\mathcal{J}$ est un faisceau d'idéaux localement nilpotent.
Comme $\mathcal{I}$ et $\mathcal{J}$ sont de type fini et $X$
est quasi-compact, il existe un $n > 0$ tel que
$\mathcal{I}^n \mathcal{J}^n = 0$. On peut donc remplacer, et l'on remplace, $\mathcal{I}$
par $\mathcal{I}^n$ et $\mathcal{J}$ par $\mathcal{J}^n$. On a alors
$\mathcal{I} \mathcal{J} = 0$. Soit $b : X' \to X$ l'éclatement
le long de $\mathcal{I} + \mathcal{J}$. Cet éclatement est $U \cup V$-admissible
car $|V(\mathcal{I} + \mathcal{J})| = |X| \setminus |U| \cup |V|$.
Nous allons montrer que $X'$ est une réunion disjointe de sous-espaces ouverts
$X' = X'_1 \amalg X'_2$, comme dans l'énoncé du lemme.

\medskip\noindent
Comme $|V(\mathcal{I} + \mathcal{J})|$ est le complémentaire de
$|U \cup V|$, on en déduit que $V \cup U$ est schématiquement
dense dans $X'$ ; voir les
lemmes \ref{lemma-blowing-up-gives-effective-Cartier-divisor} et
\ref{lemma-complement-effective-Cartier-divisor}.
Ainsi, si une telle décomposition $X' = X'_1 \amalg X'_2$
en sous-espaces à la fois ouverts et fermés existe, $X'_1$ est l'adhérence
schématique de $U$ dans $X'$, et, de même, $X'_2$ est
l'adhérence schématique de $V$ dans $X'$. Comme $U \to X'$
et $V \to X'$ sont quasi-compacts, la formation des adhérences
schématiques commute à la localisation étale (Morphismes d'espaces,
lemme \ref{spaces-morphisms-lemma-quasi-compact-scheme-theoretic-image}).
Pour vérifier l'existence de $X'_1$ et $X'_2$, on peut donc travailler localement
pour la topologie étale sur $X$. On se ramène ainsi au cas des schémas, traité
dans la démonstration de Diviseurs, lemme
\ref{divisors-lemma-separate-disjoint-opens-by-blowing-up}.
\end{proof}




















\begin{multicols}{2}[\section{Autres chapitres}]
\noindent
Préliminaires
\begin{enumerate}
\item \hyperref[introduction-section-phantom]{Introduction}
\item \hyperref[conventions-section-phantom]{Conventions}
\item \hyperref[sets-section-phantom]{Théorie des ensembles}
\item \hyperref[categories-section-phantom]{Catégories}
\item \hyperref[topology-section-phantom]{Topologie}
\item \hyperref[sheaves-section-phantom]{Faisceaux sur les espaces}
\item \hyperref[sites-section-phantom]{Sites et faisceaux}
\item \hyperref[stacks-section-phantom]{Champs}
\item \hyperref[fields-section-phantom]{Corps}
\item \hyperref[algebra-section-phantom]{Algèbre commutative}
\item \hyperref[brauer-section-phantom]{Groupes de Brauer}
\item \hyperref[homology-section-phantom]{Algèbre homologique}
\item \hyperref[derived-section-phantom]{Catégories dérivées}
\item \hyperref[simplicial-section-phantom]{Méthodes simpliciales}
\item \hyperref[more-algebra-section-phantom]{Compléments d'algèbre}
\item \hyperref[smoothing-section-phantom]{Lissage des morphismes d'anneaux}
\item \hyperref[modules-section-phantom]{Faisceaux de modules}
\item \hyperref[sites-modules-section-phantom]{Modules sur les sites}
\item \hyperref[injectives-section-phantom]{Objets injectifs}
\item \hyperref[cohomology-section-phantom]{Cohomologie des faisceaux}
\item \hyperref[sites-cohomology-section-phantom]{Cohomologie sur les sites}
\item \hyperref[dga-section-phantom]{Algèbre différentielle graduée}
\item \hyperref[dpa-section-phantom]{Algèbre à puissances divisées}
\item \hyperref[sdga-section-phantom]{Faisceaux différentiels gradués}
\item \hyperref[hypercovering-section-phantom]{Hyperrecouvrements}
\end{enumerate}
Schémas
\begin{enumerate}
\setcounter{enumi}{25}
\item \hyperref[schemes-section-phantom]{Schémas}
\item \hyperref[constructions-section-phantom]{Constructions de schémas}
\item \hyperref[properties-section-phantom]{Propriétés des schémas}
\item \hyperref[morphisms-section-phantom]{Morphismes de schémas}
\item \hyperref[coherent-section-phantom]{Cohomologie des schémas}
\item \hyperref[divisors-section-phantom]{Diviseurs}
\item \hyperref[limits-section-phantom]{Limites de schémas}
\item \hyperref[varieties-section-phantom]{Variétés}
\item \hyperref[topologies-section-phantom]{Topologies sur les schémas}
\item \hyperref[descent-section-phantom]{Descente}
\item \hyperref[perfect-section-phantom]{Catégories dérivées de schémas}
\item \hyperref[more-morphisms-section-phantom]{Compléments sur les morphismes}
\item \hyperref[flat-section-phantom]{Compléments sur la platitude}
\item \hyperref[groupoids-section-phantom]{Schémas en groupoïdes}
\item \hyperref[more-groupoids-section-phantom]{Compléments sur les schémas en groupoïdes}
\item \hyperref[etale-section-phantom]{Morphismes étales de schémas}
\end{enumerate}
Thèmes de théorie des schémas
\begin{enumerate}
\setcounter{enumi}{41}
\item \hyperref[chow-section-phantom]{Homologie de Chow}
\item \hyperref[intersection-section-phantom]{Théorie de l'intersection}
\item \hyperref[pic-section-phantom]{Schémas de Picard des courbes}
\item \hyperref[weil-section-phantom]{Théories cohomologiques de Weil}
\item \hyperref[adequate-section-phantom]{Modules adéquats}
\item \hyperref[dualizing-section-phantom]{Complexes dualisants}
\item \hyperref[duality-section-phantom]{Dualité pour les schémas}
\item \hyperref[discriminant-section-phantom]{Discriminants et différentes}
\item \hyperref[derham-section-phantom]{Cohomologie de de Rham}
\item \hyperref[local-cohomology-section-phantom]{Cohomologie locale}
\item \hyperref[algebraization-section-phantom]{Géométrie algébrique et formelle}
\item \hyperref[curves-section-phantom]{Courbes algébriques}
\item \hyperref[resolve-section-phantom]{Résolution des surfaces}
\item \hyperref[models-section-phantom]{Réduction semi-stable}
\item \hyperref[functors-section-phantom]{Foncteurs et morphismes}
\item \hyperref[equiv-section-phantom]{Catégories dérivées de variétés}
\item \hyperref[pione-section-phantom]{Groupes fondamentaux de schémas}
\item \hyperref[etale-cohomology-section-phantom]{Cohomologie étale}
\item \hyperref[crystalline-section-phantom]{Cohomologie cristalline}
\item \hyperref[proetale-section-phantom]{Cohomologie pro-étale}
\item \hyperref[relative-cycles-section-phantom]{Cycles relatifs}
\item \hyperref[more-etale-section-phantom]{Compléments de cohomologie étale}
\item \hyperref[trace-section-phantom]{La formule des traces}
\end{enumerate}
Espaces algébriques
\begin{enumerate}
\setcounter{enumi}{64}
\item \hyperref[spaces-section-phantom]{Espaces algébriques}
\item \hyperref[spaces-properties-section-phantom]{Propriétés des espaces algébriques}
\item \hyperref[spaces-morphisms-section-phantom]{Morphismes d'espaces algébriques}
\item \hyperref[decent-spaces-section-phantom]{Espaces algébriques décents}
\item \hyperref[spaces-cohomology-section-phantom]{Cohomologie des espaces algébriques}
\item \hyperref[spaces-limits-section-phantom]{Limites d'espaces algébriques}
\item \hyperref[spaces-divisors-section-phantom]{Diviseurs sur les espaces algébriques}
\item \hyperref[spaces-over-fields-section-phantom]{Espaces algébriques sur des corps}
\item \hyperref[spaces-topologies-section-phantom]{Topologies sur les espaces algébriques}
\item \hyperref[spaces-descent-section-phantom]{Descente et espaces algébriques}
\item \hyperref[spaces-perfect-section-phantom]{Catégories dérivées d'espaces}
\item \hyperref[spaces-more-morphisms-section-phantom]{Compléments sur les morphismes d'espaces}
\item \hyperref[spaces-flat-section-phantom]{Platitude sur les espaces algébriques}
\item \hyperref[spaces-groupoids-section-phantom]{Groupoïdes dans les espaces algébriques}
\item \hyperref[spaces-more-groupoids-section-phantom]{Compléments sur les groupoïdes dans les espaces}
\item \hyperref[bootstrap-section-phantom]{Amorçage}
\item \hyperref[spaces-pushouts-section-phantom]{Sommes amalgamées d'espaces algébriques}
\end{enumerate}
Thèmes de géométrie
\begin{enumerate}
\setcounter{enumi}{81}
\item \hyperref[spaces-chow-section-phantom]{Groupes de Chow des espaces}
\item \hyperref[groupoids-quotients-section-phantom]{Quotients de groupoïdes}
\item \hyperref[spaces-more-cohomology-section-phantom]{Compléments sur la cohomologie des espaces}
\item \hyperref[spaces-simplicial-section-phantom]{Espaces simpliciaux}
\item \hyperref[spaces-duality-section-phantom]{Dualité pour les espaces}
\item \hyperref[formal-spaces-section-phantom]{Espaces algébriques formels}
\item \hyperref[restricted-section-phantom]{Algébrisation des espaces formels}
\item \hyperref[spaces-resolve-section-phantom]{Résolution des surfaces revisitée}
\end{enumerate}
Théorie des déformations
\begin{enumerate}
\setcounter{enumi}{89}
\item \hyperref[formal-defos-section-phantom]{Théorie formelle des déformations}
\item \hyperref[defos-section-phantom]{Théorie des déformations}
\item \hyperref[cotangent-section-phantom]{Le complexe cotangent}
\item \hyperref[examples-defos-section-phantom]{Problèmes de déformation}
\end{enumerate}
Champs algébriques
\begin{enumerate}
\setcounter{enumi}{93}
\item \hyperref[algebraic-section-phantom]{Champs algébriques}
\item \hyperref[examples-stacks-section-phantom]{Exemples de champs}
\item \hyperref[stacks-sheaves-section-phantom]{Faisceaux sur les champs algébriques}
\item \hyperref[criteria-section-phantom]{Critères de représentabilité}
\item \hyperref[artin-section-phantom]{Axiomes d'Artin}
\item \hyperref[quot-section-phantom]{Espaces de Quot et de Hilbert}
\item \hyperref[stacks-properties-section-phantom]{Propriétés des champs algébriques}
\item \hyperref[stacks-morphisms-section-phantom]{Morphismes de champs algébriques}
\item \hyperref[stacks-limits-section-phantom]{Limites de champs algébriques}
\item \hyperref[stacks-cohomology-section-phantom]{Cohomologie des champs algébriques}
\item \hyperref[stacks-perfect-section-phantom]{Catégories dérivées de champs}
\item \hyperref[stacks-introduction-section-phantom]{Introduction aux champs algébriques}
\item \hyperref[stacks-more-morphisms-section-phantom]{Compléments sur les morphismes de champs}
\item \hyperref[stacks-geometry-section-phantom]{La géométrie des champs}
\end{enumerate}
Thèmes de théorie des espaces de modules
\begin{enumerate}
\setcounter{enumi}{107}
\item \hyperref[moduli-section-phantom]{Champs de modules}
\item \hyperref[moduli-curves-section-phantom]{Espaces de modules de courbes}
\end{enumerate}
Divers
\begin{enumerate}
\setcounter{enumi}{109}
\item \hyperref[examples-section-phantom]{Exemples}
\item \hyperref[exercises-section-phantom]{Exercices}
\item \hyperref[guide-section-phantom]{Guide bibliographique}
\item \hyperref[desirables-section-phantom]{Desiderata}
\item \hyperref[coding-section-phantom]{Style de programmation}
\item \hyperref[obsolete-section-phantom]{Obsolète}
\item \hyperref[fdl-section-phantom]{Licence de documentation libre GNU}
\item \hyperref[index-section-phantom]{Index généré automatiquement}
\end{enumerate}
\end{multicols}


\bibliography{my}
\bibliographystyle{amsalpha}

\end{document}
